\documentclass[12pt,a4paper,oneside]{report}

% --- Pakete ---
\usepackage[utf8]{inputenc}
\usepackage[T1]{fontenc}
\usepackage[ngerman]{babel}
\usepackage[a4paper,margin=2.5cm]{geometry}
\usepackage{amsmath}
\usepackage{amssymb}    % \mathbb in Kapitel 2 (LoRA-Notation)
\usepackage{setspace}
\usepackage{csquotes}

% Bessere Umbrüche / Vermeidung von Overfull hbox
\usepackage{microtype}      % feine Laufweite, bessere Silbentrennung
\emergencystretch=2em       % LaTeX darf Zeilen bei Bedarf "ziehen"

% =========================
% BibLaTeX / APA-Zitierstil
% =========================
\usepackage[
  backend=biber,
  style=apa,
  natbib=true,
  doi=true,
  url=true,
  eprint=true
]{biblatex}

\addbibresource{references.bib}

% BibLaTeX: ganzer Zitiertext als Link, nicht nur das Jahr
\DeclareFieldFormat{citehyperref}{%
  \DeclareFieldAlias{bibhyperref}{noformat}%
  \bibhyperref{#1}}

\DeclareFieldFormat{textcitehyperref}{%
  \DeclareFieldAlias{bibhyperref}{noformat}%
  \bibhyperref{%
    #1%
    \ifbool{cbx:parens}
      {\bibcloseparen\global\boolfalse{cbx:parens}}
      {}}}

\savebibmacro{cite}
\savebibmacro{textcite}

\renewbibmacro*{cite}{%
  \printtext[citehyperref]{%
    \restorebibmacro{cite}%
    \usebibmacro{cite}}}

\renewbibmacro*{textcite}{%
  \ifboolexpr{
    ( not test {\iffieldundef{prenote}} and
      test {\ifnumequal{\value{citecount}}{1}} )
    or
    ( not test {\iffieldundef{postnote}} and
      test {\ifnumequal{\value{citecount}}{\value{citetotal}}} )
  }
    {\DeclareFieldAlias{textcitehyperref}{noformat}}
    {}%
  \printtext[textcitehyperref]{%
    \restorebibmacro{textcite}%
    \usebibmacro{textcite}}}

\DeclareLanguageMapping{ngerman}{ngerman-apa}

\DefineBibliographyStrings{ngerman}{
  retrieved = {Zuletzt geprüft am}
}

% Damit \cite{...} wie ein APA-Klammerzitat erscheint: (Autor, Jahr)
\let\cite\parencite

\usepackage{xcolor}
\usepackage{xurl} % bessere Zeilenumbrüche in URLs

\usepackage[
  colorlinks=true,
  linkcolor=black,
  citecolor=blue!60!black,
  urlcolor=blue!60!black,
  breaklinks=true
]{hyperref}

\usepackage[noforwardlinks]{acronym}
\usepackage{titlesec}
\usepackage{graphicx}
\usepackage{tikz}
\usetikzlibrary{positioning,arrows.meta,calc,fit,backgrounds}
\usepackage{pgfplots}
\pgfplotsset{compat=1.18}
\usepackage{lmodern}

\usepackage{tabularx}
\usepackage{booktabs}   % Tabellenlinien (alle Tabellen)
\usepackage{listings}
\onehalfspacing
\setcounter{secnumdepth}{3}
\setcounter{tocdepth}{3}
\setlength{\parindent}{0pt}
\usepackage{parskip}    % Absatzabstand statt Einzug; eine Zeile, bei Bedarf entfernen

% Kapitelüberschriften wie "1 Einleitung" in einer Zeile
\titleformat{\chapter}[hang]
  {\normalfont\Huge\bfseries}
  {\thechapter}
  {1em}
  {}

\titlespacing*{\chapter}{0pt}{-10pt}{25pt}

% Sections wie im Beispiel
\titleformat{\section}[hang]
  {\normalfont\LARGE\bfseries}
  {\thesection}
  {1em}
  {}

\titlespacing*{\section}{0pt}{18pt}{12pt}

% Subsections etwas kompakter
\titleformat{\subsection}[hang]
  {\normalfont\large\bfseries}
  {\thesubsection}
  {1em}
  {}

\titlespacing*{\subsection}{0pt}{12pt}{8pt}

% Gleitobjekte: weniger fast leere Seiten
\renewcommand{\topfraction}{0.85}
\renewcommand{\bottomfraction}{0.5}
\renewcommand{\textfraction}{0.1}
\renewcommand{\floatpagefraction}{0.75}

\begin{document}

% =========================
% Vorderteil
% =========================
\pagenumbering{roman}
\pagestyle{empty}

% Deckblatt
\hypersetup{pageanchor=false}
\begin{titlepage}
\thispagestyle{empty}
\thispdfpagelabel{Deckblatt}% PDF-Seitenbezeichnung statt eines zweiten "i" (titlepage setzt die Zählung zurück)

\begin{tikzpicture}[remember picture, overlay]

    % --- Logos oben ---
    \node[anchor=north west] at ([xshift=1.9cm,yshift=-0.6cm]current page.north west)
        {\includegraphics[width=5.1cm]{frontmatter/logos/tu_berlin.png}};

    \node[anchor=north east] at ([xshift=-1.9cm,yshift=-1.25cm]current page.north east)
        {\includegraphics[width=7.4cm]{frontmatter/logos/dfki.jpg}};


    % --- 1. Block ---
    \node[
        anchor=center,
        align=center,
        text width=14cm
    ] (block1) at ([yshift=5.6cm]current page.center) {%
        \fontsize{23}{27}\selectfont
        \textbf{Bachelorarbeit}\\[0.5cm]
        {\fontsize{19}{23}\selectfont an der Technischen Universität Berlin}
    };

    % --- 2. Block ---
    \node[
        anchor=north,
        inner sep=0pt
    ] (block2) at ([yshift=-2.2cm]block1.south) {%
        \parbox{19cm}{\centering\bfseries\fontsize{18}{27}\selectfont
        Ein empirischer Vergleich zwischen\\
        Encoder-basierten Modellen und LoRA-finetuned LLMs\\
        für Named Entity Recognition\par}%
    };

    % --- 3. Block ---
    \node[
        anchor=north,
        align=center,
        text width=13cm
    ] (block3) at ([yshift=-2.1cm]block2.south) {%
        \fontsize{19}{23}\selectfont
        Fakultät IV\\[0.35cm]
        {\fontsize{17}{20}\selectfont Fachgebiet Quality \& Usability}
    };

    % --- Abgabedatum ---
    \node[anchor=south, align=center] (datum) at ([yshift=2.0cm]current page.south) {%
        \fontsize{14}{18}\selectfont
        Berlin, 26.~November 2026
    };

    % Vertikale Mitte zwischen Fachgebiet und Datum
    \coordinate (mitte) at ($(block3.south)!0.5!(datum.north)$);

    % --- Linker Block unten (vertikal mittig zwischen Block 3 und Datum) ---
    \node[
        anchor=west,
        align=left,
        text width=6.2cm,
        xshift=3.3cm
    ] (verfasser) at (current page.west |- mitte) {%
        \fontsize{14}{18}\selectfont
        Verfasser: \\[0.10cm]
        Luaj Osman \\[0.10cm]
        B.Sc. Wirtschaftsinformatik \\[0.10cm]
        Matrikelnummer: 456094
    };

    % --- Rechter Block unten (Oberkante bündig mit dem linken Block) ---
    \node[
        anchor=north east,
        align=left,
        text width=7.1cm,
        xshift=-2.7cm
    ] at (current page.east |- verfasser.north) {%
        \fontsize{14}{18}\selectfont
        Prüfer: \\[0.10cm]
        Prof.\ Dr.\ Sebastian Möller \\[0.10cm]
        Prof.\ Dr.\ Axel Küpper \\[0.10cm]
        Betreuer: Dr.\ Philippe Thomas
    };

\end{tikzpicture}

\end{titlepage}
\clearpage
\hypersetup{pageanchor=true}

% Eigenständigkeitserklärung
\chapter*{Eigenständigkeitserklärung}
\thispagestyle{empty}

Hiermit versichere ich, dass ich die vorliegende Arbeit ohne Hilfe Dritter und ausschließlich unter Verwendung der aufgeführten Quellen und Hilfsmittel angefertigt habe. Alle Stellen, die den benutzten Quellen und Hilfsmitteln wörtlich oder inhaltlich entnommen sind, habe ich als solche kenntlich gemacht. Sofern KI-Tools verwendet wurden, habe ich Produktnamen, Hersteller, die jeweils verwendete Softwareversion und die Art der Nutzung (z.\,B. sprachliche Überprüfung und Verbesserung der Texte, systematische Recherche) benannt. Ich verantworte die Auswahl, die Übernahme und sämtliche Ergebnisse des von mir verwendeten KI-generierten Outputs vollumfänglich selbst. IB / Stand: April 2025 Die Satzung zur Sicherung guter wissenschaftlicher Praxis an der TU Berlin vom 15. Februar 2023.\\
\url{https://www.static.tu.berlin/fileadmin/www/10002457/K3-AMBl/Amtsblatt_2023/Amtliches_Mitteilungsblatt_Nr._16_vom_30.05.2023.pdf}\\
habe ich zur Kenntnis genommen.

\vspace{1cm}

Ich erkläre weiterhin, dass ich die Arbeit in gleicher oder ähnlicher Form noch keiner anderen Prüfungsbehörde vorgelegt habe.

\vspace{1.5cm}

Berlin, den XX.XX.XXXX\\
Luaj Osman
\clearpage

% Abstract
\chapter*{Abstract}
\thispagestyle{plain}

Named Entity Recognition (NER) findet benannte Entitäten wie Personen,
Orte oder Organisationen in Texten.
Encoder ordnen dazu jedem Token ein Label zu, Large Language Models
(LLMs) erzeugen die Entitäten als strukturierten Text.
Diese Arbeit vergleicht beide Wege: DeBERTa-v3-base und -large mit
vollständigem Fine-Tuning sowie Qwen3.5-0.8B, -4B und -27B im Zero-Shot-
und QLoRA-Regime auf dem englischen Teil von MultiNERD mit
15 Entitätstypen und automatisch erzeugter Referenzannotation.
Die fünf trainierten Konfigurationen werden mit je drei Seeds
ausgewertet.
Erfasst werden strikte Metriken auf Entitätsebene, die Formtreue der
JSON-Ausgaben, Trainingsdauer, Inferenzlatenz und GPU-Spitzenspeicher
der Inferenz; hinzu kommt ein Transfertest auf CrossNER.

DeBERTa-v3-large erreicht einen Mikro-F$_1$ von 95,8 Punkten,
Qwen3.5-27B im QLoRA-Regime zunächst 94,5.
Doch das Limit von 256 Ausgabetokens schneidet lange Antworten ab;
abgeschnittene zählen als leere Vorhersagen.
Die selektive Nachinferenz mit 1\,024 Tokens hebt das LLM auf 95,6 und
verringert den Abstand von 1,34 auf 0,18 Punkte (gepaartes
95-\%-Bootstrap-Intervall: 0,02 bis 0,33).
Im Makro-F$_1$ wird gegenüber dem großen Encoder kein eindeutiger
Unterschied nachgewiesen.
Auf CrossNER mit den gemeinsamen Typen PER, ORG und LOC kehrt sich die
Rangfolge bei höherem Limit um: Das LLM erreicht 87,1 Punkte, der Encoder
85,0 (Intervall der gepoolten Differenz: 1,47 bis 2,80 Punkte).

Die Encoder benötigen auf MultiNERD 9 bis 17~ms je Satz und 0,7 bis
1,7~GiB Inferenzspeicher, die QLoRA-Läufe mit 256 Ausgabetokens 1,9 bis
5,5~s und bis zu 18,7~GiB.
Auf gleicher Hardware ergibt sich für Qwen3.5-0.8B gegenüber
DeBERTa-v3-base ein Latenzfaktor von etwa 210.
Die Zero-Shot-Konfigurationen erreichen 6,2 bis 69,4 F$_1$-Punkte.
Prompt und Decoding dieser Läufe wurden nach Beobachtungen auf dem
Testsplit gewählt.
Da sich die Regime zudem über das Training hinaus in Prompt, Decoding und
numerischer Präzision unterscheiden, ist kein isolierter
Fine-Tuning-Effekt bestimmbar.

Auf MultiNERD erzielt bei jedem Latenz- und Speicherbudget, das eine
getestete Konfiguration zulässt, ein Encoder den höchsten Mikro-F$_1$.
Auf CrossNER benötigt das LLM rund 13~s je Satz und 18,6~GiB
Inferenzspeicher; ob der Qualitätsgewinn das rechtfertigt, hängt von der
Anwendung ab.
Der Qualitätsvergleich hängt jedoch wesentlich von Ausgabelimit,
Datensatz und der Typabbildung auf PER, ORG und LOC ab.
Die Arbeit zeigt damit zweierlei: den Ressourcenunterschied der konkreten
Pipelines und die Bedeutung ausreichender Ausgabelimits für die Bewertung
generativer NER.
Eine allgemeine Überlegenheit einer Modellklasse wird aus den
untersuchten Konfigurationen nicht abgeleitet.

\addcontentsline{toc}{chapter}{Abstract}
\clearpage

% Inhaltsverzeichnis
\begingroup
\makeatletter
\let\ps@plain\ps@empty
\makeatother
\begin{singlespace}
\tableofcontents
\clearpage
\listoffigures
\addcontentsline{toc}{chapter}{Abbildungsverzeichnis}
\clearpage
\listoftables
\addcontentsline{toc}{chapter}{Tabellenverzeichnis}
\clearpage
\end{singlespace}
\endgroup

% Abkürzungsverzeichnis
\chapter*{Abkürzungsverzeichnis}
\thispagestyle{plain}

% Alphabetisch sortiert (Groß-/Kleinschreibung ignoriert).
% Aufgenommen sind die fachspezifischen Abkürzungen, die in der Arbeit
% mehrfach verwendet werden. Bewusst nicht aufgenommen:
% - Eigennamen von Datensätzen, Konferenzen, Software und Institutionen
%   (MultiNERD, CrossNER, OntoNotes, SemEval, CUDA, SLURM, DFKI, ...);
%   Modellnamen nur, wenn sie selbst Akronyme sind (BERT, GPT, ...),
% - allgemein bekannte IT-Abkürzungen (API, JSON, XML, YAML, CPU, ...),
% - Abkürzungen, die nur an einer Stelle vorkommen und dort im Text
%   erklärt werden (NSP, DoRA, rsLoRA, KV, ...),
% - seltene Tagging-Schema-Varianten (IOBES, BILOU, IOE), die in
%   Kapitel 2 erläutert werden; das Tagging-Schema IOB2 (auch BIO) steht
%   unter IOB, da BIO in MultiNERD zugleich ein Entitätstyp ist,
% - Maßeinheiten (GB, MiB, ms) und Entitätstypen-Labels (PER, LOC, ...).

\begin{singlespace}
\begin{acronym}[RoBERTa]
    \setlength{\parskip}{3pt}  % Zeilenabstand 1,5 plus etwas Luft zwischen den Einträgen
    \acro{BERT}{Bidirectional Encoder Representations from Transformers}
    \acro{BF16}{BFloat16 (Brain Floating Point, 16 Bit)}
    \acro{BiLSTM}{Bidirectional Long Short-Term Memory}
    \acro{CoNLL}{Conference on Computational Natural Language Learning}
    \acro{CRF}{Conditional Random Field}
    \acro{DeBERTa}{Decoding-enhanced BERT with Disentangled Attention}
    \acro{ELECTRA}{Efficiently Learning an Encoder that Classifies Token Replacements Accurately}
    \acro{F1}[F$_1$]{F$_1$-Score, harmonisches Mittel aus Precision und Recall}
    \acro{FF}{Forschungsfrage}
    \acro{FN}{False Negatives}
    \acro{FP}{False Positives}
    \acro{FP32}{32-Bit-Gleitkommaformat (Single Precision)}
    \acro{GPT}{Generative Pre-trained Transformer}
    \acro{GPU}{Graphics Processing Unit}
    \acro{HBM3}{High Bandwidth Memory 3 (GPU-Speichertechnologie)}
    \acro{IE}{Information Extraction}
    \acro{IOB}{Inside, Outside, Beginning (Tagging-Schema; Varianten IOB1 und IOB2, auch BIO genannt)}
    \acro{LLaMA}{Large Language Model Meta AI}
    \acro{LLM}{Large Language Model}
    \acro{LoRA}{Low-Rank Adaptation}
    \acro{MLM}{Masked Language Modeling}
    \acro{MRC}{Machine Reading Comprehension}
    \acro{MTP}{Multi-Token-Prediction}
    \acro{NER}{Named Entity Recognition}
    \acro{NF4}{4-Bit-NormalFloat}
    \acro{NLP}{Natural Language Processing}
    \acro{PEFT}{Parameter-Efficient Fine-Tuning}
    \acro{QLoRA}{Quantized Low-Rank Adaptation}
    \acro{RE}{Relation Extraction}
    \acro{RoBERTa}{Robustly Optimized BERT Pre-training Approach}
    \acro{SFT}{Supervised Fine-Tuning}
    \acro{TP}{True Positives}
    \acro{UIE}{Universal Information Extraction}
\end{acronym}
\end{singlespace}

\addcontentsline{toc}{chapter}{Abkürzungsverzeichnis}
\clearpage

% =========================
% Hauptteil
% =========================
\pagenumbering{arabic}
\pagestyle{plain}

% chapters/01_einleitung.tex
%
% Einleitung mit Forschungsfragen FF1 bis FF3 (Abschnitt 1.2).

\chapter{Einleitung}
\label{chap:einleitung}

Dieses Kapitel erläutert Motivation, Problemstellung und Zielsetzung der
vorliegenden Arbeit.
Abschnitt~\ref{sec:motivation_kontext} beschreibt den Anwendungskontext
und die beiden untersuchten Modellklassen.
Abschnitt~\ref{sec:problem_ziel} ordnet den Forschungsstand ein und leitet
daraus die Forschungsfragen und den Beitrag der Arbeit ab.
Abschnitt~\ref{sec:abgrenzung} legt den Untersuchungsumfang fest,
Abschnitt~\ref{sec:aufbau_arbeit} gibt einen Überblick über die weiteren
Kapitel.

\section{Motivation und Kontext}
\label{sec:motivation_kontext}

In Krankenhäusern, Nachrichtenredaktionen und Behörden entstehen täglich
zahlreiche Textdokumente, darunter Arztbriefe, Nachrichtenartikel und
Gerichtsentscheidungen.
Die darin enthaltenen Informationen liegen überwiegend unstrukturiert vor,
und ihre manuelle Erschließung ist aufwendig.
Natural Language Processing (NLP) umfasst Verfahren, mit denen solche
Texte automatisch verarbeitet und analysiert
werden~\cite{jurafsky2024speech,manning1999foundations}.
Eine grundlegende Aufgabe ist die Named Entity Recognition (NER):
Erwähnungen benannter Entitäten werden im Text lokalisiert und
vorgegebenen Kategorien zugeordnet~\cite{li2018survey}.
Die Ergebnisse dienen unter anderem als Grundlage für
Relationsextraktion, den Aufbau von Wissensgraphen und
Frage-Antwort-Systeme~\cite{jurafsky2024speech,keraghel2024recent}.

Welche Kategorien eine Anwendung benötigt, hängt vom Einsatzzweck ab.
Klassische Benchmarks wie CoNLL-2003 unterscheiden nur Personen,
Organisationen, Orte und eine Restkategorie~\cite{sang2003conll}.
In vielen Texten sind jedoch weitere Typen relevant, in Arztbriefen etwa
Krankheiten, in Nachrichtenartikeln Ereignisse oder Werke.
Feingranulare Datensätze wie MultiNERD bilden solche Typen ab und
unterscheiden 15 Kategorien, darunter Krankheiten, Ereignisse und
Medien~\cite{tedeschi2022multinerd}.

Die vorliegende Arbeit entstand am Deutschen Forschungszentrum für
Künstliche Intelligenz (DFKI) in der Forschungsgruppe Speech and Language
Technology (SLT).
Dort bildet die strukturierte Extraktion von Informationen aus
unstrukturierten Textbeständen ein zentrales Anwendungsfeld.

Für NER stehen heute zwei Klassen vortrainierter Sprachmodelle zur
Verfügung.
Encoder-Modelle wie BERT oder DeBERTa berechnen für jedes Token eine
kontextabhängige Repräsentation; ein darauf trainierter Klassifikator
weist jedem Token ein Label zu, aus dem sich die Entitäten ergeben
(Sequenzlabeling)~\cite{devlin2019bert,he2021debertav3}.
Large Language Models (LLMs) sind deutlich größere, generative
Sprachmodelle, die ein breites Spektrum von Aufgaben mit demselben Modell
bearbeiten~\cite{brown2020language,zhao2023survey}.
Sie lösen NER, indem sie die erkannten Entitäten als strukturierten Text
ausgeben~\cite{zhan2026generativeNER,semin2026spanLabeling}.
Im einfachsten Fall erhält das Modell dazu lediglich eine
Aufgabenbeschreibung im Prompt, ohne auf annotierten Daten trainiert zu
werden (\emph{Zero-Shot})~\cite{brown2020language}; alternativ wird es mit
annotierten Beispielen an die Aufgabe angepasst~\cite{zhan2026generativeNER}.

Die beiden Modellklassen unterscheiden sich nicht nur in ihrer
Arbeitsweise, sondern auch im Ressourcenbedarf.
Encoder wie DeBERTa-v3 umfassen einige hundert Millionen
Parameter~\cite{he2021debertav3}, aktuelle LLMs dagegen Milliarden bis
Hunderte Milliarden~\cite{qwen2026qwen35blog}.
Zudem erzeugen LLMs ihre Ausgabe Token für Token, sodass die Inferenzzeit
mit der Länge der Ausgabe wächst~\cite{lu2024padellm}, während ein Encoder
die Labels eines Satzes in einem einzigen Vorwärtsdurchlauf
bestimmt~\cite{devlin2019bert}.
Auch die Anpassung ist bei LLMs aufwendiger: Ein vollständiges
Fine-Tuning aller Gewichte erfordert zusätzlich Speicher für Gradienten
und Optimiererzustände~\cite{hu2021lora}.
Parameter-Efficient Fine-Tuning (PEFT) umfasst Verfahren, die stattdessen
nur eine kleine Zahl zusätzlicher Parameter trainieren~\cite{han2024parameter};
mit QLoRA lassen sich so auch Modelle mit Dutzenden Milliarden Parametern
auf einer einzelnen GPU anpassen~\cite{dettmers2023qlora}.
Für die praktische Modellwahl stellt sich damit die Frage, wann der
zusätzliche Aufwand eines LLM durch eine bessere Erkennungsleistung
gerechtfertigt ist und wann ein kompakter Encoder die effizientere Wahl
darstellt.

\section{Problemstellung, Forschungsfragen und Beitrag}
\label{sec:problem_ziel}

Mehrere Arbeiten vergleichen generative LLMs bereits mit Encoder-basierten
NER-Verfahren.
Ohne aufgabenspezifische Anpassung liegen LLMs dabei deutlich zurück:
In der Vergleichsstudie von \textcite{keraghel2024recent} erreicht GPT-4o
mit Few-Shot-Prompting auf CoNLL-2003 einen Makro-F$_1$ von 65{,}17, ein
vollständig angepasstes DeBERTa-v3-large dagegen 90{,}98.
Mit LoRA angepasste LLMs schließen diesen Abstand.
\textcite{zhan2026generativeNER} berichten für LLaMA3.1-8B im besten
Ausgabeformat einen F$_1$ von 93{,}85 auf CoNLL-2003 gegenüber 93{,}04 für
das Encoder-Verfahren BERT-MRC und zeigen zugleich, dass das Ausgabeformat
die Erkennungsleistung stark beeinflusst.
\textcite{obeidat2025llm} vergleichen auf fünf biomedizinischen
Datensätzen Encoder mit vollständigem Fine-Tuning und mittels QLoRA
angepasste LLMs mit 7 bis 8~Mrd. Parametern.
Die LLMs liegen auf vier Datensätzen um 2 bis 8 F$_1$-Punkte vorn,
benötigen je Testbeispiel aber das 40- bis 220-Fache der Inferenzzeit des
besten Encoders.

In der hier gesichteten Literatur lassen diese Vergleiche vier Fragen
offen, die Abschnitt~\ref{sec:research_gap} ausführlich herleitet.
Erstens verwenden sie entweder allgemeine Datensätze mit wenigen groben
Typen oder domänenspezifische Korpora; ein allgemeiner, feingranularer
Datensatz wurde für den direkten Vergleich nicht herangezogen.
Zweitens stammen die Encoder-Referenzwerte teilweise aus anderen
Publikationen statt aus derselben Pipeline.
Drittens werden Erkennungsleistung und Ressourcenbedarf nur teilweise
gemeinsam erfasst: \textcite{obeidat2025llm} messen Inferenzzeiten, aber
keinen Speicherbedarf, \textcite{zhan2026generativeNER} berichten nur
Trainingskosten.
Viertens stammen die untersuchten LLMs aus einer einzigen Größenklasse
zwischen 3 und 9~Mrd. Parametern, sodass offen bleibt, wie deutlich
kleinere oder größere Modelle im Vergleich zu Encodern abschneiden.

Die vorliegende Arbeit setzt an diesen Punkten an.
Sie vergleicht auf dem englischen Teil von MultiNERD mit seinen
15 Entitätstypen~\cite{tedeschi2022multinerd} drei Regime:
DeBERTa-v3-base und DeBERTa-v3-large mit vollständigem
Fine-Tuning~\cite{he2021debertav3}, Qwen3.5-0.8B, Qwen3.5-4B und
Qwen3.5-27B~\cite{qwen2026qwen35blog} im Zero-Shot-Einsatz sowie dieselben
Qwen3.5-Modelle nach einer Anpassung mit QLoRA.
Das Zero-Shot-Regime zeigt, was die veröffentlichten Modelle ohne
Trainingsaufwand leisten, und dient als Bezugspunkt für den Nutzen der
Anpassung.
QLoRA trainiert Low-Rank-Adapter auf einem in 4~Bit quantisierten,
eingefrorenen Basismodell~\cite{hu2021lora,dettmers2023qlora} und macht
so auch die Anpassung des 27B-Modells auf einer einzelnen GPU möglich.
Alle Konfigurationen verwenden dieselben offiziellen Datensplits und werden
mit derselben strikten Metrik auf Entitätsebene bewertet; die generierten
Ausgaben der LLMs werden dazu in dieselbe Label-Repräsentation überführt
wie die Vorhersagen der Encoder.
Neben der Erkennungsleistung werden die Trainingsdauer, die
Inferenzlatenz und der GPU-Spitzenspeicher während der Inferenz erfasst.

Daraus ergeben sich drei Forschungsfragen (FF1--FF3):
\begin{description}
  \item[FF1 (Erkennungsleistung):] Wie unterscheiden sich
  DeBERTa-v3-Encoder mit vollständigem Fine-Tuning und Qwen3.5-LLMs im
  Zero-Shot- und im QLoRA-Regime in ihrer Erkennungsleistung auf dem
  englischen MultiNERD-Testsplit, und wie zuverlässig halten die LLMs das
  vorgegebene Ausgabeformat ein (Formtreue)?
  \item[FF2 (Modellgröße und Ressourcen):] Wie verändern sich
  Erkennungsleistung, Trainingsdauer, Inferenzlatenz und
  GPU-Spitzenspeicher der Inferenz mit der Modellgröße innerhalb beider
  Modellfamilien?
  \item[FF3 (Empfehlungen):] Welche Konfigurationen bieten unter gegebenen
  Speicher- und Latenzbeschränkungen das günstigste Verhältnis von
  Erkennungsleistung und Ressourcenbedarf, und unter welchen Bedingungen
  rechtfertigt ein LLM den Mehraufwand gegenüber einem Encoder?
\end{description}

Die Arbeit leistet damit folgende Beiträge:
\begin{itemize}
  \item einen Vergleich von acht Konfigurationen beider Modellklassen in
  einer gemeinsamen Pipeline mit eigener Encoder-Referenz statt
  übernommener Werte; die fünf trainierten Konfigurationen werden mit je
  drei Seeds ausgewertet, sodass sich insgesamt 18 Testevaluationen
  ergeben;
  \item eine Größenstaffelung über zwei Modellfamilien von rund 184~Mio.
  bis 27~Mrd. Parametern mit gemeinsamer Erfassung von
  Erkennungsleistung und Ressourcenbedarf;
  \item ein JSON-Ausgabeformat, das Entitäten über Wortindizes adressiert
  und ohne Constrained Decoding erzeugt wird, sodass die Formtreue der
  LLM-Ausgaben eigens gemessen werden kann~\cite{semin2026spanLabeling};
  \item einen ergänzenden Transfertest auf CrossNER, der prüft, wie sich
  die jeweils größten auf MultiNERD angepassten Modelle beider Familien
  auf Wikipedia-Texte aus fünf anderen Domänen übertragen
  lassen~\cite{liu2021crossner};
  \item eine Absicherung der Befunde durch gepaarte
  Bootstrap-Konfidenzintervalle über die Testsätze und durch zwei
  Sensitivitätsanalysen zum Ausgabelimit der LLM-Inferenz, die den
  wichtigsten Störfaktor des Vergleichs messen statt schätzen.
\end{itemize}
Der Transfertest verwendet ausschließlich die Testsplits von CrossNER,
ohne weiteres Training, und wertet nur die gemeinsamen Grobtypen Person,
Organisation und Ort aus.
Er dient der Einordnung der Hauptergebnisse, nicht dem Vergleich mit
Bestwerten der Literatur.

\section{Abgrenzung der Arbeit}
\label{sec:abgrenzung}

Der angemeldete Titel der Arbeit nennt BERT-basierte Modelle und LoRA.
Auf der Encoder-Seite werden mit DeBERTa-v3-base und DeBERTa-v3-large
zwei Weiterentwicklungen von BERT untersucht, die sich in vergleichenden
NER-Studien als besonders leistungsfähig erwiesen haben
(Abschnitt~\ref{sec:encoder_ner}); auf der LLM-Seite wird LoRA in seiner
speicheroptimierten Form QLoRA eingesetzt, die das Training des
27B-Modells auf einer einzelnen GPU erst möglich macht
(Abschnitt~\ref{sec:lora}), und das Zero-Shot-Regime kommt als
Bezugspunkt ohne Training hinzu.
Im Folgenden werden die Verfahren daher mit diesen präziseren
Bezeichnungen benannt.

Die Untersuchung beschränkt sich auf NER in englischsprachigen Texten.
Hauptdatenbasis sind die offiziellen Splits des englischen Teils von
MultiNERD mit Texten aus Wikipedia und Wikinews~\cite{tedeschi2022multinerd}.
Da diese Splits automatisch annotiert wurden, messen die berichteten Werte
die Übereinstimmung mit einer automatisch erzeugten Referenz
(Abschnitte~\ref{sec:datasets_corpora} und~\ref{sec:datenauswahl}).
CrossNER wird ausschließlich als Transfertest auf den Testsplits und für
die Typen Person, Organisation und Ort ausgewertet~\cite{liu2021crossner}.
CoNLL-2003 und OntoNotes dienen nur der Einordnung verwandter Arbeiten und
werden nicht experimentell untersucht.
Eine Übertragung der Schlussfolgerungen auf andere Sprachen oder
Anwendungsdomänen wird nicht abgeleitet.

Auf der Encoder-Seite werden DeBERTa-v3-base und DeBERTa-v3-large
eingesetzt, die sich in vergleichenden NER-Untersuchungen als besonders
leistungsfähig erwiesen haben~\cite{keraghel2024recent}.
Sie werden um einen linearen Klassifikationskopf ergänzt und
ausschließlich mit vollständigem Fine-Tuning angepasst, bei dem
Encoder-Gewichte und Klassifikationskopf gemeinsam trainiert
werden~\cite{devlin2019bert}; ein Zero-Shot-Einsatz ist ohne trainierten
Klassifikationskopf nicht möglich.
Auf der LLM-Seite werden drei Größen der Qwen3.5-Familie untersucht, die
einen Vergleich innerhalb einer Modellfamilie erlauben~\cite{qwen2026qwen35blog}.
Sie werden ohne Training im Zero-Shot-Regime sowie nach einer Anpassung
mit QLoRA eingesetzt, bei der nur die Low-Rank-Adapter trainiert
werden~\cite{hu2021lora,dettmers2023qlora}.
Nicht untersucht werden ein vollständiges Fine-Tuning der LLMs,
alternative PEFT-Verfahren, weitere Modellfamilien, Few-Shot-Prompting
mit Demonstrationsbeispielen, Constrained Decoding und der Thinking-Modus
der Qwen3.5-Modelle; andere Modellfamilien und Verfahren werden nur im
Literaturüberblick berücksichtigt.

Die Regime unterscheiden sich in mehreren Faktoren zugleich, darunter
Vortraining, Architektur, Modellgröße, Quantisierung und
Anpassungsverfahren; zudem liefen die Konfigurationen auf
unterschiedlichen GPU-Typen (Abschnitt~\ref{sec:validitaet}).
Die Ergebnisse sind daher als Vergleich konkreter Modell- und
Anpassungskonfigurationen zu lesen; isolierte kausale Effekte einzelner
Faktoren werden mit diesem Versuchsdesign nicht bestimmt.

Die Erkennungsleistung wird mit Precision, Recall und F$_1$-Score auf
Entitätsebene gemessen; eine Entität gilt nur dann als korrekt erkannt,
wenn ihre Grenzen und ihr Typ mit der Referenz
übereinstimmen~\cite{manning2008introduction,sang2003conll}.
Zentrale Kennzahl ist der Mikro-F$_1$-Score.
Für die LLMs wird zusätzlich die Formtreue ihrer Ausgaben erfasst
(Abschnitt~\ref{sec:robustheit-ausgaben}).
Der Ressourcenbedarf umfasst die Trainingsdauer, die Inferenzlatenz bei
einzeln verarbeiteten Sätzen und den GPU-Spitzenspeicher während der
Inferenz (Abschnitt~\ref{sec:effizienzmetriken}).
Nicht erfasst werden der Speicherbedarf während des Trainings, ein durch
Batching optimierter Durchsatz sowie Energieverbrauch und Kosten.
Eine mögliche Kontamination der Vortrainingsdaten der LLMs mit
MultiNERD-Texten wird diskutiert, aber nicht empirisch geprüft
(Abschnitt~\ref{sec:research_gap}).
Interpretierbarkeit, Fairness und Robustheit gegenüber adversarialen
Eingaben sind ebenso wenig Gegenstand der Arbeit wie die Erhebung und
Annotation eines neuen Datensatzes.

\section{Aufbau der Arbeit}
\label{sec:aufbau_arbeit}

Die weiteren Kapitel führen von den Grundlagen und dem Forschungsstand
über die Methodik zu den Ergebnissen und Schlussfolgerungen.

Kapitel~\ref{chap:grundlagen} führt in NLP und NER ein und beschreibt die
Transformer-Architektur, Encoder-Modelle für NER, generative LLMs sowie
parameter-effizientes Fine-Tuning mit LoRA und QLoRA.
Nach einer Einordnung von Skalierungsgesetzen stellt es die verwendeten
Datensätze, darunter MultiNERD und CrossNER, sowie die
Evaluationsmetriken vor.

Kapitel~\ref{chap:verwandte_arbeiten} stellt Encoder-basierte und
LLM-basierte NER-Ansätze einschließlich verschiedener Ausgabeformate vor.
Es behandelt Arbeiten, die mit PEFT angepasste LLMs direkt mit Encodern
vergleichen, fasst Befunde zum Einfluss der Modellgröße zusammen und
leitet daraus die Forschungslücke ab.

Kapitel~\ref{chap:methodik} beschreibt die drei Vergleichsregime, die
Auswahl von Modellen und Daten sowie die Encoder- und die LLM-Pipeline
einschließlich Prompt-Design, QLoRA-Konfiguration und Nachverarbeitung der
generierten Ausgaben.
Darauf folgen das Evaluationskonzept für Erkennungsleistung,
Ressourcenbedarf und Formtreue, das Versuchsdesign des CrossNER-Transfertests
und eine Zusammenfassung der bekannten Abweichungen vom idealisierten
Versuchsdesign.

Kapitel~\ref{chap:ergebnisse} präsentiert die Erkennungsleistung aller
Konfigurationen insgesamt und je Entitätstyp sowie die Formtreue der
LLM-Ausgaben, untersucht den Einfluss der Modellgröße und den
Ressourcenbedarf, analysiert die Fehlerprofile der Modellklassen,
berichtet die Ergebnisse des CrossNER-Transfertests und diskutiert die
Befunde.

Kapitel~\ref{chap:fazit} beantwortet die Forschungsfragen, erläutert die
Grenzen der Untersuchung und skizziert weiterführende Arbeiten.
Der Anhang dokumentiert die Reproduzierbarkeit der Läufe mit
Konfigurationen, Codeständen, Modellrevisionen und nachrangigen
Einstellungen sowie den vollständigen Zero-Shot-Prompt.

% chapters/02_grundlagen.tex
%
% Grundlagenkapitel.
% Benötigt in der Präambel: amsmath, graphicx, booktabs.

\chapter{Grundlagen}
\label{chap:grundlagen}

Dieses Kapitel führt die Konzepte ein, die das Versuchsdesign in
Kapitel~\ref{chap:methodik} voraussetzt.
Es erklärt Begriffe und Mechanismen; die Datensätze sowie die
Einordnung einzelner Forschungsarbeiten und ihrer Ergebnisse folgen in
Kapitel~\ref{chap:verwandte_arbeiten}.

\section{Natural Language Processing (NLP)}
\label{sec:nlp}

Natural Language Processing (NLP) umfasst Verfahren, mit denen Computer
natürliche Sprache in geschriebener oder gesprochener Form verarbeiten,
analysieren und erzeugen~\cite{jurafsky2024speech,manning1999foundations}.
Typische Aufgaben sind die Klassifikation ganzer Texte, die Zuweisung von
Labels an einzelne Wörter einer Sequenz, die Extraktion strukturierter
Informationen und die Erzeugung von Text.
Named Entity Recognition gehört zur zweiten Gruppe und wird in
Abschnitt~\ref{sec:ner} eingeführt.

Heutige Systeme beruhen auf kontextualisierten Repräsentationen: Ein
neuronales Netz liest die gesamte Eingabesequenz und berechnet für jedes
Token einen Vektor, der von den umgebenden Tokens
abhängt~\cite{devlin2019bert}.
Daraus ist das heute dominierende Vorgehen in zwei Schritten
entstanden: Ein Modell wird zunächst auf großen, nicht annotierten
Textkorpora vortrainiert und anschließend an eine konkrete Aufgabe
angepasst, sei es durch Fine-Tuning seiner Gewichte oder durch
Anweisungen im Eingabetext~\cite{devlin2019bert,brown2020language}.
Vortrainierte Modelle lassen sich dabei auf zwei Arten nutzen.
Encoder liefern Repräsentationen, auf denen ein Klassifikator Labels
vorhersagt (Abschnitt~\ref{sec:encoder_ner}); generative Modelle erzeugen
die Lösung einer Aufgabe als Text (Abschnitt~\ref{sec:llms}).

\section{Named Entity Recognition (NER)}
\label{sec:ner}

Named Entity Recognition (NER) ist die Aufgabe, Erwähnungen benannter
Entitäten in einem Text zu lokalisieren und einer vordefinierten
Kategorie zuzuordnen~\cite{li2018survey}.
Formal erhält ein NER-System eine Tokenfolge $t_1, \dots, t_n$ und gibt
eine Menge von Tripeln $(I_s, I_e, \tau)$ zurück, wobei $I_s$ und $I_e$
den Start- und Endindex einer Erwähnung und $\tau$ ihren Typ
bezeichnen; die Formalisierung folgt \textcite{li2018survey}.
Eine Erwähnung ist damit ein zusammenhängender Abschnitt der Tokenfolge,
ein \emph{Span}; korrekt ist eine Vorhersage nur, wenn Grenzen und Typ
stimmen.
Das Prinzip zeigt Abbildung~\ref{fig:ner-grundprinzip} an einem
Beispiel.
Ob der Endindex das letzte Token einschließt oder auf das Token danach
zeigt, ist Konvention; die Abbildung verwendet einschließende Indizes,
die LLM-Pipeline in Abschnitt~\ref{sec:prompt-design} einen
ausschließenden Endindex.

\begin{figure}[htbp]
  \centering
  % figures/ner_grundprinzip.tex -- Grundprinzip der NER (eigene Darstellung,
% TikZ). Beispiel wie in Tabelle 2.1; Endindizes einschliessend.
\begin{tikzpicture}[
  font=\sffamily\small,
  tok/.style={draw=gray!60, fill=gray!8, rounded corners=2pt,
    minimum width=1.75cm, minimum height=0.8cm, inner sep=3pt},
  per/.style={tok, fill=orange!22}, loc/.style={tok, fill=blue!14},
  box/.style={draw=gray!60, fill=green!12, rounded corners=3pt,
    minimum width=5.2cm, minimum height=1.25cm, align=center,
    font=\sffamily\normalsize},
  outbox/.style={draw=gray!60, fill=white, rounded corners=3pt, inner sep=7pt},
  arr/.style={-{Latex[length=2.5mm]}, thick, gray!60!black},
  lbl/.style={font=\sffamily\footnotesize, text=gray!70!black},
]
  \node[lbl, anchor=west] at (-4.95, 1.25) {Eingabe: Tokenfolge $t_1, \dots, t_5$};
  \foreach \i/\w/\s in {1/Barack/per, 2/Obama/per, 3/visited/tok, 4/Berlin/loc, 5/./tok} {
    \node[\s] (t\i) at ({(\i-3)*2.05}, 0) {\w};
    \node[lbl] at ({(\i-3)*2.05}, 0.68) {$t_\i$};
  }
  \node[box] (ner) at (0, -2.05) {Named Entity Recognition\\(NER)};
  \draw[arr] (0, -0.42) -- (ner.north);
  \node[outbox, anchor=north] (res) at (0, -3.35) {%
    \begin{tabular}{@{}l@{\hspace{1.4em}}l@{}}
      $(t_1,\; t_2,\; \mathrm{PER})$ & \colorbox{orange!22}{\strut Barack Obama}\\[2pt]
      $(t_4,\; t_4,\; \mathrm{LOC})$ & \colorbox{blue!14}{\strut Berlin}
    \end{tabular}};
  \draw[arr] (ner.south) -- (res.north);
  \node[lbl, anchor=north] at ($(res.south)+(0,-0.12)$)
    {Ausgabe: Tripel (Startindex, Endindex, Typ), Endindex einschließend};
\end{tikzpicture}

  \caption[Grundprinzip der Named Entity Recognition]{Grundprinzip der Named Entity Recognition: Aus einer
  Tokenfolge werden Erwähnungen als Tripel aus Startindex, Endindex und
  Typ extrahiert.
  Eigene Darstellung in Anlehnung an
  \textcite[Abb.~1]{li2018survey}; der Endindex gehört zur Erwähnung.}
  \label{fig:ner-grundprinzip}
\end{figure}

Welche Kategorien unterschieden werden, legt der jeweilige Datensatz
fest.
Klassische Benchmarks verwenden wenige grobe Typen wie Person,
Organisation und Ort, feingranulare Datensätze wie MultiNERD
unterscheiden mehr als ein Dutzend Typen
(Abschnitt~\ref{sec:datasets_corpora}).

In der Praxis wird NER häufig als Sequenzlabeling formuliert: Jedes Token
erhält genau ein Label, das zugleich die Segmentierung in Spans und deren
Typ kodiert~\cite{li2018survey}.
Das Tagset $\{\mathrm{I}, \mathrm{O}, \mathrm{B}\}$ geht auf
\textcite{ramshaw1995text} zurück.
In der verbreiteten Variante IOB2, häufig kurz BIO-Schema genannt,
beginnt jede Entität mit einem B-Tag, I markiert ihre weiteren Tokens
und O alle Tokens außerhalb~\cite{tjong1999representing}.
Für NER wird der Typ an das Tag angehängt, etwa \texttt{B-PER} und
\texttt{I-PER}.
Die Kodierung setzt voraus, dass sich Entitäten weder überlappen noch
ineinander verschachtelt sind; CoNLL-2003 trifft genau diese
Annahme~\cite{sang2003conll}, und auch die in dieser Arbeit verwendeten
Datensätze enthalten nur flache Annotationen.
Ein Beispiel zeigt Tabelle~\ref{tab:bio-beispiel}.
Diese Arbeit verwendet IOB2 in der Encoder-Pipeline
(Abschnitt~\ref{sec:encoder-pipeline}) ebenso wie in der Auswertung
(Abschnitt~\ref{sec:evaluation}).

\begin{table}[htbp]
  \centering
  \small
  \renewcommand{\arraystretch}{1.3}
  \caption[IOB2-Kodierung eines Satzes]{IOB2-Kodierung eines Satzes.}
  \label{tab:bio-beispiel}
  \begin{tabular}{lccccc}
    \toprule
    \textbf{Token} & Barack & Obama & visited & Berlin & . \\
    \midrule
    \textbf{IOB2} & \texttt{B-PER} & \texttt{I-PER}
                  & \texttt{O} & \texttt{B-LOC} & \texttt{O} \\
    
    \bottomrule
  \end{tabular}
\end{table}

Lange Zeit waren Conditional Random Fields (CRF) für diese
Formulierung verbreitet; sie modellieren die bedingte
Wahrscheinlichkeit der gesamten Labelfolge bei gegebener Eingabe, sodass
Übergänge zwischen benachbarten Labels mitgelernt werden und unzulässige
Folgen unwahrscheinlich werden~\cite{lafferty2001crf}.
Neuronale Sequenzmodelle wie das BiLSTM-CRF ersetzten von Hand
entworfene Merkmale durch gelernte
Repräsentationen~\cite{lample2016neural,ma2016end}.

\section{Transformer-Architektur}
\label{sec:transformer}

Der Transformer von \textcite{vaswani2017attention} beruht
vollständig auf Attention: Jedes Token kann in einem Schritt auf jede
andere Position der Sequenz zugreifen.

Grundbaustein ist die \emph{Self-Attention}.
Aus der Repräsentation jedes Tokens werden durch gelernte lineare
Abbildungen drei Vektoren erzeugt: eine Query, ein Key und ein Value.
Das Skalarprodukt zwischen der Query eines Tokens und den Keys aller
Tokens misst deren Relevanz; die Softmax-Funktion normiert die
Relevanzwerte zu Gewichten, und die Ausgabe für jedes Token ist die
gewichtete Summe der
Value-Vektoren~\cite{vaswani2017attention}.
Transformer verwenden mehrere parallele \emph{Köpfe} mit jeweils
eigenen Projektionen.
Da Attention die Reihenfolge der Tokens nicht berücksichtigt, wird die
Position jedes Tokens eigens kodiert
(Abschnitt~\ref{sec:encoder_ner}).

Das Originalmodell besteht aus einem Encoder und einem Decoder.
Eine Transformer-Schicht des Encoders besteht aus zwei Teilschichten: der
Multi-Head-Self-Attention und einem positionsweisen Feed-Forward-Netz,
das jeden Tokenvektor unabhängig transformiert.
Der Decoder ergänzt jede Schicht um eine dritte Teilschicht, die auf die
Encoder-Ausgabe zugreift, und maskiert seine Self-Attention so, dass
Position $i$ nur auf Positionen bis $i$ zugreifen kann.
Diese \emph{kausale Maske} stellt sicher, dass die Vorhersage eines
Tokens nur von den bereits bekannten Tokens abhängt.
Den Aufbau zeigt Abbildung~\ref{fig:transformer-architektur}.

\begin{figure}[htbp]
  \centering
  % figures/vaswani_transformer.tex -- Transformer-Architektur (eigene
% Darstellung in Anlehnung an Vaswani et al. 2017, TikZ).
\begin{tikzpicture}[
  font=\sffamily\footnotesize,
  blk/.style={draw=gray!60, rounded corners=2pt, minimum width=3.5cm,
    minimum height=0.64cm, align=center, inner sep=2pt},
  attn/.style={blk, fill=orange!22}, ffn/.style={blk, fill=blue!12},
  norm/.style={blk, fill=green!16, minimum height=0.5cm},
  emb/.style={blk, fill=red!10}, lin/.style={blk, fill=violet!14},
  smx/.style={blk, fill=yellow!22, minimum height=0.5cm},
  stack/.style={draw=gray!50, rounded corners=4pt, fill=gray!6, inner sep=5pt},
  arr/.style={-{Latex[length=2mm]}, thick, black!75},
  res/.style={-{Latex[length=2mm]}, thick, black!75, rounded corners=2pt},
  pe/.style={circle, draw=gray!60, inner sep=1pt, minimum size=0.42cm, fill=white},
]
  % ---------------- Encoder (x = 0)
  \node[emb]  (ein) at (0, -0.15) {Eingabe-Embedding};
  \node[attn] (esa) at (0, 1.9)  {Multi-Head-\\Self-Attention};
  \node[norm] (en1) at (0, 2.75) {Add \& Norm};
  \node[ffn]  (eff) at (0, 3.75) {Feed-Forward-\\Netz};
  \node[norm] (en2) at (0, 4.6)  {Add \& Norm};
  \begin{scope}[on background layer]
    \node[stack, fit={(esa)(en2)(-2.05, 1.3)}] (enc) {};
  \end{scope}
  \node[font=\sffamily\small, anchor=east] at ($(enc.west)+(-0.25,0)$) {$N\times$};
  \node (einp) at (0, -1.3) {Eingaben};
  \node[pe] (eplus) at (0, 0.72) {$+$};
  \node[pe] (epos) at (-2.0, 0.72) {$\sim$};
  \node[align=center, font=\sffamily\scriptsize, anchor=east] at (-2.3, 0.72) {Positions-\\kodierung};
  \draw[arr] (einp) -- (ein);
  \draw[arr] (ein) -- (eplus); \draw[arr] (epos) -- (eplus); \draw[arr] (eplus) -- (esa);
  \draw[arr] (esa) -- (en1); \draw[arr] (en1) -- (eff); \draw[arr] (eff) -- (en2);
  \draw[res] (0, 1.3)  -- ++(-2.05, 0) |- (en1.west);
  \draw[res] (0, 3.2)  -- ++(-2.05, 0) |- (en2.west);
  % ---------------- Decoder (x = 6.6)
  \node[emb]  (aus) at (6.6, -0.15) {Ausgabe-Embedding};
  \node[attn] (dma) at (6.6, 1.9)  {Maskierte Multi-Head-\\Self-Attention};
  \node[norm] (dn1) at (6.6, 2.75) {Add \& Norm};
  \node[attn] (dca) at (6.6, 3.75) {Encoder-Decoder-\\Attention};
  \node[norm] (dn2) at (6.6, 4.6)  {Add \& Norm};
  \node[ffn]  (dff) at (6.6, 5.6)  {Feed-Forward-\\Netz};
  \node[norm] (dn3) at (6.6, 6.45) {Add \& Norm};
  \begin{scope}[on background layer]
    \node[stack, fit={(dma)(dn3)(8.65, 1.3)}] (dec) {};
  \end{scope}
  \node[font=\sffamily\small, anchor=west] at ($(dec.east)+(0.25,0)$) {$N\times$};
  \node[lin]  (lin)  at (6.6, 7.55) {Lineare Schicht};
  \node[smx]  (soft) at (6.6, 8.35) {Softmax};
  \node[align=center] (outp) at (6.6, 9.3) {Ausgabe-\\wahrscheinlichkeiten};
  \node[align=center] (ausp) at (6.6, -1.3) {Ausgaben\\(rechtsverschoben)};
  \node[pe] (dplus) at (6.6, 0.72) {$+$};
  \node[pe] (dpos) at (8.6, 0.72) {$\sim$};
  \node[align=center, font=\sffamily\scriptsize, anchor=west] at (8.9, 0.72) {Positions-\\kodierung};
  \draw[arr] (ausp) -- (aus);
  \draw[arr] (aus) -- (dplus); \draw[arr] (dpos) -- (dplus); \draw[arr] (dplus) -- (dma);
  \draw[arr] (dma) -- (dn1); \draw[arr] (dn1) -- (dca); \draw[arr] (dca) -- (dn2);
  \draw[arr] (dn2) -- (dff); \draw[arr] (dff) -- (dn3); \draw[arr] (dn3) -- (lin);
  \draw[arr] (lin) -- (soft); \draw[arr] (soft) -- (outp);
  \draw[res] (6.6, 1.3)  -- ++(2.05, 0) |- (dn1.east);
  \draw[res] (6.6, 3.2)  -- ++(2.05, 0) |- (dn2.east);
  \draw[res] (6.6, 5.05) -- ++(2.05, 0) |- (dn3.east);
  % ---------------- Encoder-Ausgabe -> Encoder-Decoder-Attention (Keys, Values)
  \draw[arr] (en2.north) -- ++(0, 0.45) -| (3.3, 3.75) -- (dca.west);
  \node[font=\sffamily\scriptsize, anchor=south, inner sep=2pt] at (1.65, 5.3) {Encoder-Ausgabe (Keys, Values)};
\end{tikzpicture}

  \caption[Transformer-Architektur mit Encoder und Decoder]{Transformer-Architektur mit Encoder (links) und Decoder
  (rechts).
  Jede Encoder-Schicht enthält Multi-Head-Self-Attention und ein
  Feed-Forward-Netz, jeweils mit Residualverbindung und
  Layer-Normalisierung; der Decoder ersetzt die Self-Attention durch eine
  maskierte und ergänzt eine Attention über die Encoder-Ausgabe.
  Eigene Darstellung in Anlehnung an \textcite[Abb.~1]{vaswani2017attention}.}
  \label{fig:transformer-architektur}
\end{figure}

Aus diesem Bauplan sind drei Modellfamilien hervorgegangen.
\emph{Encoder-only}-Modelle wie BERT verwenden nur den Encoder; jedes
Token sieht den Kontext in beide Richtungen~\cite{devlin2019bert}.
\emph{Decoder-only}-Modelle wie GPT-3 verwenden nur den Decoder ohne
Encoder-Attention; durch die kausale Maske eignen sie sich zur
Texterzeugung Token für Token~\cite{brown2020language}.
Encoder-Decoder-Modelle behalten den Aufbau des Originals bei.
Für NER folgt daraus der zentrale Unterschied der beiden verglichenen
Modellklassen: Ein Encoder liefert für jedes Token eine
beidseitig kontextualisierte Repräsentation, aus der ein Label
klassifiziert werden kann.
Ein Decoder verarbeitet die Sequenz kausal und gibt in der hier
verwendeten Pipeline die Entitäten als Text aus. Da der ganze
Eingabesatz im Prompt liegt, kann er dabei dennoch auch den Kontext
rechts von einer Entität nutzen.

Die Kosten der Self-Attention wachsen quadratisch mit der
Sequenzlänge~\cite{vaswani2017attention}.
Bei der Texterzeugung werden zudem die Keys und Values aller bisherigen
Tokens zwischengespeichert; Grouped-Query-Attention lässt mehrere
Query-Köpfe denselben Key- und Value-Kopf nutzen und verkleinert so
diesen Zwischenspeicher~\cite{ainslie2023gqa}.
Lineare Attention-Verfahren wie Gated DeltaNet ersetzen die
Relevanzmatrix durch einen Zustand fester Größe, sodass der Aufwand nur
linear mit der Sequenzlänge wächst~\cite{yang2025gated}.
Die in dieser Arbeit verwendeten LLMs wechseln lineare Schichten mit
einzelnen Schichten voller Attention ab~\cite{qwen2026card08b,qwen2026card4b,qwen2026card27b}; diese
hybride Attention ist für die Platzierung der Adapter relevant
(Abschnitt~\ref{sec:lora-konfiguration}).

\section{Encoder-basierte Modelle für NER}
\label{sec:encoder_ner}

Vortrainierte Transformer-Encoder sind seit BERT~\cite{devlin2019bert}
die übliche Architektur für überwachte NER.
Bevor ein Encoder einen Text verarbeiten kann, muss dieser in Einheiten
aus einem festen Vokabular zerlegt werden.
Da ein Vokabular aus ganzen Wörtern seltene Wörter nicht abdecken kann,
arbeiten alle hier behandelten Modelle mit \emph{Subwort-Einheiten}:
Häufige Wörter bleiben ganz, seltene werden in mehrere Teilstücke
zerlegt.
DeBERTa-v3 verwendet ein
SentencePiece-Vokabular~\cite{kudo2018sentencepiece} mit
128\,000 Einheiten~\cite{he2021debertav3}.
Für NER hat die Zerlegung eine praktische Folge: Die Labels liegen je
Wort vor, die Modelleingabe besteht aber aus Subwörtern.
\textcite{devlin2019bert} verwenden deshalb die Repräsentation des ersten
Subworts eines Wortes als Eingabe des Token-Klassifikators; in der
üblichen Umsetzung erhalten die übrigen Subwörter kein Label und gehen
nicht in den Verlust ein.
Die Encoder-Pipeline dieser Arbeit folgt diesem Vorgehen
(Abschnitt~\ref{sec:datenvorverarbeitung}).

\textcite{devlin2019bert} trainieren BERT auf nicht annotiertem Text vor allem mit
\emph{Masked Language Modeling} (MLM): Ein Teil der Tokens wird maskiert
und aus dem beidseitigen Kontext rekonstruiert.
Für eine Zielaufgabe wird das vortrainierte Modell um eine
aufgabenspezifische Ausgabeschicht ergänzt und vollständig feinjustiert,
das heißt, alle Encoder-Gewichte und die neue Schicht werden gemeinsam
mit dem Verlust der Zielaufgabe optimiert.
Für NER ist diese Ausgabeschicht eine lineare Klassifikationsschicht, die
aus der Repräsentation $h_i$ jedes Tokens Werte $W h_i + b$ für alle
Labels des Tagsets berechnet; eine Softmax-Funktion macht daraus eine
Verteilung, und bei der Vorhersage wird das Label mit dem höchsten Wert
gewählt.
Trainiert wird mit der Kreuzentropie zwischen dieser Verteilung und dem
Referenzlabel, gemittelt über alle Tokens mit Label.
Da die Positionskodierung von BERT auf 512 Positionen festgelegt
ist~\cite[Anh.~A.2]{devlin2019bert},
müssen längere Eingaben gekürzt oder geteilt werden; auch für
DeBERTa-v3 ist die maximale Eingabelänge auf 512 Positionen konfiguriert
(Parameter \texttt{max\_position\_embeddings} in den
Konfigurationsdateien der Revisionen in
Tabelle~\ref{tab:app-revisionen}),
sodass eine Längenbegrenzung ebenso für die Encoder-Pipeline dieser Arbeit
gilt (Abschnitt~\ref{sec:datenvorverarbeitung}).

DeBERTa verändert zwei Komponenten des
Encoders~\cite{he2020deberta}.
Erstens trennt die \emph{Disentangled Attention} Inhalt und Position:
Jedes Token wird durch einen Inhaltsvektor und einen Vektor für seine
relative Position zu den anderen Tokens repräsentiert, sodass die
Attention getrennt lernen kann, welche Inhalte und welche Abstände
zusammengehören.
Zweitens bezieht der \emph{Enhanced Mask Decoder} die absoluten
Positionen erst unmittelbar vor der Vorhersage der maskierten Tokens
ein.
DeBERTa-v3 behält diese Architektur bei und ändert das Vortraining in
zwei Punkten~\cite{he2021debertav3}.
Statt MLM verwendet es \emph{Replaced Token Detection} nach ELECTRA: Ein
kleines Generatormodell ersetzt einige Tokens durch plausible
Alternativen, und das eigentliche Modell lernt als Diskriminator, für
jedes Token zu entscheiden, ob es ersetzt wurde.
Das Trainingssignal entsteht so an allen Positionen statt nur an den
maskierten, was das Vortraining dateneffizienter
macht~\cite{clark2020electra}.
Zudem teilen Generator und Diskriminator beim \emph{Gradient-Disentangled
Embedding Sharing} ihre Token-Embeddings, ohne dass die Gradienten des
Diskriminators in die gemeinsamen Embeddings
zurückfließen~\cite{he2021debertav3}.

Die Base-Variante hat zwölf Schichten mit 86~Mio.
Parametern im Backbone, die Large-Variante 24 Schichten mit 304~Mio.,
jeweils zuzüglich der Embedding-Parameter des großen
Vokabulars~\cite{he2021debertav3card,he2021debertav3cardlarge}.

\section{Large Language Models (LLMs)}
\label{sec:llms}

Large Language Models (LLMs) sind Sprachmodelle mit typischerweise
Milliarden von Parametern, die auf sehr großen Textkorpora vortrainiert werden und als einzelnes
Modell ein breites Spektrum von Aufgaben bearbeiten~\cite{brown2020language,zhao2023survey}.
Die heute verbreiteten LLMs sind Decoder-only-Transformer
(Abschnitt~\ref{sec:transformer}).
Sie modellieren die Wahrscheinlichkeit einer Tokenfolge autoregressiv
als Produkt bedingter Wahrscheinlichkeiten,
\begin{equation}
  p(x_1, \dots, x_T) = \prod_{t=1}^{T} p(x_t \mid x_1, \dots, x_{t-1}),
  \label{eq:lm}
\end{equation}
und werden darauf trainiert, an jeder Position das nächste Token
vorherzusagen~\cite{brown2020language}.
Text wird erzeugt, indem das Modell wiederholt eine Verteilung über das
nächste Token berechnet, ein Token auswählt und es an die Eingabe
anhängt.
Jedes erzeugte Token erfordert einen eigenen Vorwärtsdurchlauf, der mit
wachsendem Kontext teurer wird (Abschnitt~\ref{sec:transformer}).
Die Ausgabelänge ist neben der Eingabelänge ein wesentlicher Faktor der
Inferenzzeit, was bei generativer NER die Wahl des Ausgabeformats zu
einer Effizienzfrage macht (Abschnitt~\ref{sec:llm_ner_related}).
Wie das Token ausgewählt wird, legt das \emph{Decoding} fest.
Greedy-Decoding wählt stets das wahrscheinlichste Token und ist als
Auswahlregel deterministisch; Gleitkommaarithmetik und Batching können
hardwareseitig Abweichungen verursachen.
Beim Sampling wird ein Token zufällig aus der Verteilung gezogen; eine
Temperatur skaliert die Verteilung, Top-$k$ beschränkt die Auswahl auf
die $k$ wahrscheinlichsten Tokens, und Top-$p$ (\emph{Nucleus Sampling})
auf die kleinste Menge von Tokens, deren Wahrscheinlichkeiten sich zu
mindestens $p$ summieren~\cite{holtzman2020curious}.
Gesampelte Ausgaben variieren zwischen Durchläufen, was bei der
Interpretation von Ergebnissen zu berücksichtigen ist
(Abschnitt~\ref{sec:llm-pipeline}).

Das Vortraining liefert ein Basismodell, das Text fortsetzt, aber
Anweisungen nicht zuverlässig befolgt.
\textcite{ouyang2022instructgpt} zeigen, wie ein \emph{Nachtraining}
(\emph{Post-Training}) dies ändert, unter anderem durch überwachtes
Fine-Tuning auf von Menschen geschriebenen Antworten
(\emph{Supervised Fine-Tuning}, SFT).
Nachtrainierte Modelle verarbeiten Eingaben in einem Dialogformat aus
Nachrichten mit Rollen wie System, Nutzer und Assistent, das ein
\emph{Chat-Template} mit speziellen Steuertokens in eine Tokenfolge
überführt~\cite{qwen2025qwen3,qwen2026card27b}.
Qwen-Modelle ab Qwen3 können vor der eigentlichen Antwort einen
abgegrenzten Block mit Zwischenüberlegungen erzeugen; dieser
\emph{Thinking-Modus} lässt sich ein- und
ausschalten~\cite{qwen2025qwen3}.

Nachtrainierte LLMs lösen Aufgaben, die im Eingabetext beschrieben
werden: Beim \emph{In-Context Learning}~\cite{brown2020language} erhält
das Modell eine Anweisung und optional einige
Demonstrationsbeispiele im Prompt und löst die Aufgabe ohne Änderung
seiner Parameter.
Enthält der Prompt keine Beispiele, spricht man von Zero-Shot, sonst von
Few-Shot.
Einige Fähigkeiten treten dabei erst ab einer gewissen Modellgröße auf;
\textcite{wei2022emergent} bezeichnen sie als emergent.

Viele LLMs werden nur mit frei herunterladbaren Gewichten und einer
Lizenz veröffentlicht (\emph{offene Gewichte}, \emph{open weights}).
Zu dieser Gruppe gehört die Qwen-Familie von Alibaba mit den Generationen
Qwen2.5, Qwen3 und Qwen3.5~\cite{qwen2024qwen25,qwen2025qwen3}; zur im
Februar 2026 vorgestellten Generation
Qwen3.5~\cite{qwen2026qwen35blog} gehören die hier verwendeten Modelle
Qwen3.5-27B, Qwen3.5-4B und Qwen3.5-0.8B, deren Modellkarten im
Schichtaufbau Feed-Forward-Netze und keine
Mixture-of-Experts-Schichten ausweisen~\cite{qwen2026card27b,qwen2026card4b,qwen2026card08b}
(Abschnitt~\ref{sec:llm-auswahl}).

Für NER ergibt sich aus der generativen Arbeitsweise eine andere
Formulierung als beim Sequenzlabeling.
Das Modell erhält die Aufgabenbeschreibung und den Text als Prompt und
erzeugt die Entitäten als Text, etwa als strukturierte Liste im
JSON-Format~\cite{dagdelen2024structuredIE}, je nach Verfahren mit
Angaben zur Position~\cite[Abschn.~3.2]{semin2026spanLabeling}.
Diese Ausgabe muss anschließend auf die Tokenfolge abgebildet werden,
damit sie mit den Referenzlabels verglichen werden kann.
Ausgabeformat und Nachverarbeitung werden damit Teil des Systems, und
die Ausgabe kann formal fehlerhaft sein, etwa durch ungültige Indizes
oder Typen außerhalb des Inventars~\cite{semin2026spanLabeling}.
Verfahren des \emph{Constrained Decoding} binden die Erzeugung an eine
formale Grammatik~\cite{geng2023grammar}; ob die Entitäten inhaltlich
stimmen und im Eingabetext vorkommen, bleibt dabei ungesichert.
Alternativ wird frei generiert und die Ausgabe nachträglich geprüft
(Abschnitte~\ref{sec:llm_ner_related} und~\ref{sec:prompt-design}).

\section{Parameter-Efficient Fine-Tuning (LoRA und QLoRA)}
\label{sec:lora}

Vollständiges Fine-Tuning passt alle Gewichte eines vortrainierten
Modells an.
Bei LLMs stößt das an Speichergrenzen, weil während des Trainings zu
jedem Parameter auch sein Gradient und die Zustände des Optimierers
gehalten werden müssen; Adam speichert je Parameter zwei zusätzliche Größen.
\emph{Parameter-Efficient Fine-Tuning} (PEFT) hält die vortrainierten
Gewichte eingefroren und trainiert nur eine kleine Zahl zusätzlicher
oder ausgewählter Parameter~\cite{han2024parameter}.

LoRA (\emph{Low-Rank Adaptation}) trainiert dazu eine
niedrigdimensionale Darstellung der
Gewichtsänderung~\cite{hu2021lora}.
Für eine vortrainierte Gewichtsmatrix $W_0 \in \mathbb{R}^{d \times k}$
wird die Änderung $\Delta W$ durch das Produkt zweier kleiner Matrizen
dargestellt, sodass sich für eine Eingabe $x$ die Ausgabe
\begin{equation}
  h = W_0 x + \Delta W x = W_0 x + \frac{\alpha}{r}\, B A x,
  \qquad B \in \mathbb{R}^{d \times r},\;
  A \in \mathbb{R}^{r \times k},\;
  r \ll \min(d, k)
  \label{eq:lora}
\end{equation}
ergibt.
$W_0$ bleibt eingefroren; trainiert werden nur $A$ und $B$.
Der Rang $r$ bestimmt die Kapazität des Adapters und die Zahl der
trainierbaren Parameter $r\,(d + k)$ anstelle von $d \cdot k$; die
Konstante $\alpha$ skaliert den Beitrag des
Adapters~\cite[Abschn.~4.1]{hu2021lora}.
In gängigen Implementierungen dient ein Dropout auf dem Adapterpfad als
Regularisierung.
$A$ wird zufällig initialisiert und $B$ mit null, sodass $\Delta W$ zu
Beginn des Trainings null ist und das Modell unverändert
startet~\cite{hu2021lora}.
Den Aufbau zeigt Abbildung~\ref{fig:lora-adapter}.
Nach dem Training kann $\Delta W$ zu $W_0$ addiert werden, sodass
gegenüber dem Ausgangsmodell keine zusätzliche Inferenzlatenz entsteht;
alternativ bleibt der Adapter getrennt und wird bei der Inferenz
mitberechnet.
\textcite{hu2021lora} adaptieren in ihren Experimenten die Query- und
Value-Projektionen der Attention; \textcite{dettmers2023qlora} setzen
Adapter auf alle linearen Schichten
(Abschnitt~\ref{sec:lora-konfiguration}).
Auf RoBERTa, DeBERTa, GPT-2 und GPT-3 erreicht LoRA bei einem Bruchteil
der trainierbaren Parameter eine Qualität, die der des vollständigen
Fine-Tunings entspricht oder sie übertrifft~\cite{hu2021lora}.

\begin{figure}[htbp]
  \centering
  % Abbildung 2.3: Aufbau eines LoRA-Adapters (TikZ)
\begin{tikzpicture}[
  >=Latex, font=\small,
  frozen/.style={draw, rounded corners=2pt, fill=gray!15, minimum width=2.6cm, minimum height=1.1cm, align=center},
  train/.style={draw, rounded corners=2pt, fill=blue!10, minimum width=1.6cm, minimum height=0.9cm, align=center},
  io/.style={draw, circle, minimum size=0.75cm, inner sep=0pt},
  sum/.style={draw, circle, minimum size=0.5cm, inner sep=0pt},
  arr/.style={->, thick}]
  \node[io] (x) at (0,0) {$x$};
  \node[frozen] (w) at (3.4,1.1) {$W_0 \in \mathbb{R}^{d \times k}$\\[1pt]{\footnotesize eingefroren}};
  \node[train] (a) at (2.5,-1.2) {$A \in \mathbb{R}^{r \times k}$};
  \node[train] (b) at (4.7,-1.2) {$B \in \mathbb{R}^{d \times r}$};
  \node[sum] (plus) at (6.6,0) {$+$};
  \node[io] (h) at (8.3,0) {$h$};
  \draw[arr] (x) -- (1.1,0) |- (w);
  \draw[arr] (x) -- (1.1,0) |- (a);
  \draw[arr] (a) -- (b);
  \draw[arr] (w) -| (plus);
  \draw[arr] (b) -| (plus);
  \draw[arr] (plus) -- (h);
  \node[font=\footnotesize, anchor=south] at (3.6,-0.55) {Rang $r \ll \min(d,k)$};
  \node[font=\footnotesize, anchor=north] at (3.6,-1.72) {trainierbarer Low-Rank-Pfad $\Delta W = \tfrac{\alpha}{r} BA$};
  \node[font=\footnotesize, anchor=south] at (8.3,0.65) {$h = W_0 x + \tfrac{\alpha}{r} B A x$};
\end{tikzpicture}

  \caption[Aufbau eines LoRA-Adapters]{Aufbau eines LoRA-Adapters.
  Die eingefrorene Gewichtsmatrix $W_0$ wird um den trainierbaren
  Low-Rank-Pfad $\Delta W = \tfrac{\alpha}{r} BA$ mit Rang $r$ ergänzt; beide Pfade
  verarbeiten dieselbe Eingabe $x$, und ihre Ausgaben werden addiert
  (Gleichung~\ref{eq:lora}).
  Eigene Darstellung in Anlehnung an \textcite[Abb.~1]{hu2021lora}.}
  \label{fig:lora-adapter}
\end{figure}

LoRA senkt den Speicherbedarf für Gradienten und Optimiererzustände,
nicht aber den für die eingefrorenen Gewichte selbst.
Hier setzt die \emph{Quantisierung} an.
Sie bildet Gewichte, die üblicherweise mit 16 oder 32~Bit gespeichert
werden, auf eine kleine Zahl von Stufen ab, etwa 256 Stufen bei 8~Bit.
Dazu wird eine Gruppe von Gewichten durch eine Skalierungskonstante auf
den Wertebereich des Zieltyps normiert; zur Berechnung wird sie
mit dieser Konstante wieder in höhere Präzision
zurückgerechnet~\cite{dettmers2022llmint8}.
Der Speicherbedarf sinkt entsprechend der Bitbreite.

QLoRA verbindet beides: Das Basismodell wird eingefroren und mit 4~Bit
gespeichert, gerechnet wird in 16~Bit, und die LoRA-Adapter werden in
16-Bit-Präzision trainiert~\cite{dettmers2023qlora}.
\textcite{dettmers2023qlora} führen dazu drei Komponenten ein.
Der Datentyp \emph{4-Bit-NormalFloat} (NF4) legt seine 16 Stufen so
fest, dass bei normalverteilten Gewichten in jede Stufe gleich viele
Werte fallen; er nutzt die vier Bit damit besser aus als
gleichabständige Stufen.
\emph{Double Quantization} quantisiert zusätzlich die
Skalierungskonstanten der Blöcke, die bei kleinen Blöcken selbst
spürbar Speicher belegen.
\emph{Paged Optimizers} lagern die Optimiererzustände bei
Speicherspitzen automatisch in den Hauptspeicher aus und verringern so
das Risiko von Abbrüchen wegen Speichermangels.
Im Vorwärts- und Rückwärtsdurchlauf werden die 4-Bit-Gewichte blockweise
in den Rechendatentyp BFloat16 zurückgerechnet; die Gradienten fließen
durch das eingefrorene quantisierte Modell, aktualisiert werden nur die
Adapter.
Mit diesem Verfahren lässt sich ein Modell mit 65~Mrd. Parametern auf
einer einzelnen GPU mit 48~GB feinjustieren; auf den von
\textcite{dettmers2023qlora} geprüften Benchmarks erreicht QLoRA die
Aufgabenleistung des 16-Bit-Fine-Tunings (volles Fine-Tuning bei GLUE
und Super-NaturalInstructions, LoRA bei MMLU).
In dieser Arbeit werden die Encoder vollständig feinjustiert und die
LLMs mit QLoRA angepasst; welche der genannten Komponenten dabei aktiv
sind, dokumentiert Abschnitt~\ref{sec:lora-konfiguration}.

\clearpage  % 2.7 beginnt auf neuer Seite
\section{Modellgröße und Parameteranzahl}
\label{sec:model_size}

Sprachmodelle unterscheiden sich um Größenordnungen in ihrer
Parameterzahl.
Wie sich die Größe auf die Qualität auswirkt, beschreiben
\emph{Skalierungsgesetze} in Form von Potenzgesetzen: Der Verlust fällt
als Potenzfunktion der Parameterzahl, der Datenmenge und des
Rechenaufwands~\cite{kaplan2020scaling}.
Die Gesetze beschreiben allerdings den Verlust des Vortrainings von
Modellen, die eigens für die Untersuchung trainiert wurden; über
Kennzahlen einer nachgelagerten Aufgabe wie den F$_1$-Score bei NER
sagen sie nichts unmittelbar aus.
Empirische Befunde zum Größeneinfluss bei NER fasst
Abschnitt~\ref{sec:model_size_related} zusammen; die Folgen für den
Größenvergleich dieser Arbeit nennt
Abschnitt~\ref{sec:einfluss-parameteranzahl}.

\section{Evaluationsmetriken}
\label{sec:evaluation}

NER-Systeme werden mit Precision, Recall und dem F-Maß bewertet, die
aus dem Information Retrieval
stammen~\cite{vanrijsbergen1979information,manning2008introduction}.
Auf Entitätsebene bezeichnet $TP$ die Zahl korrekt erkannter Entitäten,
$FP$ die Zahl vorhergesagter Entitäten ohne Entsprechung in der Referenz
und $FN$ die Zahl der Referenzentitäten, die das System nicht gefunden
hat.
Precision ist der Anteil korrekter Vorhersagen an allen Vorhersagen,
Recall der Anteil gefundener Entitäten an allen Referenzentitäten:
\begin{equation}
  \text{Precision} = \frac{TP}{TP + FP}, \qquad
  \text{Recall}    = \frac{TP}{TP + FN}, \qquad
  F_1 = 2 \cdot \frac{\text{Precision} \cdot \text{Recall}}%
                      {\text{Precision} + \text{Recall}} .
  \label{eq:prf1}
\end{equation}
Das F-Maß, das auf das Effektivitätsmaß $E$ bei van Rijsbergen
zurückgeht, fasst beide in einer Kennzahl
zusammen~\cite{vanrijsbergen1979information,manning2008introduction}; der F$_1$-Score ist das
harmonische Mittel beider Größen und die in der NER-Forschung
üblicherweise berichtete Kennzahl~\cite{sang2003conll,keraghel2024recent}.

Entscheidend ist, was als Instanz gezählt wird.
Bei einer Evaluation auf Tokenebene wird jedes Tokenlabel einzeln
bewertet; bei einer Evaluation auf Entitätsebene ist die Instanz die Entität, und eine
Vorhersage gilt nur dann als korrekt, wenn Spangrenzen und Typ exakt
mit der Referenz übereinstimmen~\cite{sang2003conll}.
Diese strenge Zählweise ist seit CoNLL-2003 üblich.
\textcite{segura2013semeval} definieren für die SemEval-2013-Aufgabe zur
Erkennung von Arzneimittelinteraktionen daneben weniger strenge Modi, die
nur die exakten Grenzen, eine teilweise Überlappung der Grenzen oder den
Typ bei überlappenden Grenzen als Treffer werten.
Solche Modi machen partielle Treffer sichtbar, erschweren aber den
Vergleich zwischen Arbeiten; diese Arbeit verwendet ausschließlich den
strikten Modus (Abschnitt~\ref{sec:performanzmetriken}).

In der Praxis werden Metriken auf Entitätsebene aus Tagfolgen berechnet;
wie dabei Spans aus den Tags gewonnen werden, hängt von der verwendeten
Bibliothek ab.
Die in dieser Arbeit verwendete Bibliothek \texttt{seqeval} bietet zwei
Modi~\cite{nakayama2018seqeval}.
Der Standardmodus ahmt das CoNLL-Auswertungsskript nach und toleriert formal unzulässige Tagfolgen; ein
\texttt{I-}Tag ohne vorangehenden Entitätsbeginn eröffnet dort dennoch
eine Entität.
Der strikte Modus wertet die Tagfolgen nach einem explizit angegebenen
Schema wie IOB2 aus, sodass unzulässige Folgen keine Entität bilden.

Bei mehreren Entitätstypen werden die Kennzahlen aggregiert.
Der Mikro-F$_1$-Score summiert $TP$, $FP$ und $FN$ über alle Typen und
berechnet daraus eine Kennzahl, sodass häufige Typen das Ergebnis
dominieren; der Makro-F$_1$-Score mittelt die je Typ berechneten
F$_1$-Scores und gewichtet alle Typen gleich~\cite{manning2008introduction}.
Bei stark unausgewogenen Typverteilungen wie in MultiNERD, wo drei Typen
den größten Teil der Entitäten stellen, spiegelt der Mikro-F$_1$-Score
vor allem die Leistung auf diesen Typen wider; typbezogene Werte
ergänzen ihn (Abschnitt~\ref{sec:performanzmetriken}).
Für seltene Typen mit geringem Support, also wenigen
Referenzentitäten, sind typbezogene Werte allerdings mit großer Unsicherheit
behaftet, da wenige Fehler den Wert stark verändern.

\textcite{hripcsak2005agreement} zeigen zudem, dass das F-Maß beim
Vergleich zweier Annotationen der positiven spezifischen Übereinstimmung
entspricht.
Ein F$_1$-Score misst damit die Übereinstimmung mit der Referenz und
nur mittelbar die Richtigkeit; bei einer automatisch erzeugten Referenz
wie der von MultiNERD (Abschnitt~\ref{sec:datasets_corpora}) ist diese
Unterscheidung wesentlich.

% chapters/03_verwandte_arbeiten.tex
%
% Überarbeitete Fassung vom 16.10.2026.
% Alle Angaben zu zitierten Arbeiten wurden gegen die Volltexte geprüft
% (arXiv-HTML/PDF, ACL Anthology, ICLR-Proceedings, PMC). Angaben zur
% eigenen Arbeit folgen Kapitel 1 und 4.
% Benötigt in der Präambel: booktabs, array (über tabularx geladen).

\chapter{Verwandte Arbeiten}
\label{chap:verwandte_arbeiten}

Dieses Kapitel ordnet die vorliegende Arbeit in den Forschungsstand zur
Named Entity Recognition (NER) mit Encoder-Modellen und generativen Large
Language Models (LLMs) ein.
Berücksichtigt werden Arbeiten, die mindestens eines von vier Kriterien
erfüllen: Sie vergleichen Encoder und generative LLMs auf derselben
NER-Aufgabe, sie untersuchen Ausgabeformate oder Anpassungsverfahren für
LLM-basierte NER, sie verwenden MultiNERD oder CrossNER, oder sie
variieren die Modellgröße.
Begutachtete Fassungen werden bevorzugt; Preprints werden einbezogen,
wenn keine begutachtete Fassung vorliegt und die Arbeit unmittelbar
einschlägig ist, darunter zwei Arbeiten aus dem Jahr
2026~\cite{zhan2026generativeNER,semin2026spanLabeling}.
Die Auswahl beruht auf einer narrativen Recherche in Google Scholar,
arXiv und der ACL Anthology (2019 bis 2025) mit Stichwörtern wie
\enquote{LLM named entity recognition}, \enquote{generative NER},
\enquote{instruction tuning NER}, \enquote{LoRA/QLoRA NER},
\enquote{DeBERTa NER}, \enquote{MultiNERD} und \enquote{CrossNER}.
Ergänzend wurde entlang der Zitationen der Schlüsselarbeiten per
Schneeballsuche gesucht.
Trefferzahlen wurden nicht protokolliert; die Auswahl beansprucht keine
systematische Vollständigkeit.
Grundlagen zu Architekturen, LoRA und QLoRA, Skalierungsgesetzen,
Datensätzen und Metriken wurden in Kapitel~\ref{chap:grundlagen}
eingeführt und werden hier nicht wiederholt.

Abschnitt~\ref{sec:encoder_ner_related} behandelt Encoder-basierte
Ansätze, einschließlich der MultiNERD-Originalarbeit und solcher Encoder,
die mit LLM-erzeugten Trainingsdaten als Generalisten eingesetzt werden.
Abschnitt~\ref{sec:llm_ner_related} gliedert LLM-basierte Ansätze nach
Ausgabestrategie und Anpassungsart und begründet daraus die in dieser
Arbeit gewählte Ausgabestrategie.
Abschnitt~\ref{sec:peft_ner_related} stellt Arbeiten vor, die mit
Parameter-Efficient Fine-Tuning (PEFT) angepasste LLMs direkt mit
Encodern vergleichen.
Abschnitt~\ref{sec:model_size_related} fasst Befunde zum Einfluss der
Modellgröße zusammen.
Abschnitt~\ref{sec:research_gap} leitet daraus die Forschungslücke ab und
ordnet sie den Forschungsfragen aus Abschnitt~\ref{sec:problem_ziel} zu.
Eigene Ergebnisse werden nicht vorweggenommen.

\section{Encoder-basierte Ansätze für NER}
\label{sec:encoder_ner_related}

Encoder mit Klassifikationskopf bilden seit BERT den Standard für
überwachte NER (Abschnitt~\ref{sec:encoder_ner}).
Für diese Arbeit sind drei Fragen leitend: Wie ist die Datenbasis
MultiNERD entstanden und welche Referenzwerte liegen für sie vor?
Wie leistungsfähig sind feinjustierte Encoder bei feingranularer NER im
Vergleich zu per Prompting genutzten LLMs?
Und wie weit tragen Encoder, wenn sie ohne aufgabenspezifische
Annotationen eingesetzt werden?

\textcite{tedeschi2022multinerd} stellen mit MultiNERD einen Datensatz
vor, der zehn Sprachen, die Genres Wikipedia und Wikinews sowie
15 Entitätstypen abdeckt; der englische Teil umfasst rund 164\,000 Sätze.
Die Annotation erfolgt automatisch: Verlinkte Textstellen in Wikipedia-
und Wikinews-Artikeln werden über BabelNet-Synsets einem Entitätstyp
zugeordnet, auf nicht verlinkte Vorkommen übertragen und anschließend
durch ein iteratives Verfahren bereinigt, das Sätze verwirft, bei denen
ein auf den Silberdaten trainiertes Modell (mBERT mit BiLSTM und CRF) der
wissensbasierten Annotation widerspricht.
Dieses mBERT-Modell ist zugleich die einzige in der Arbeit trainierte
Architektur.
Evaluiert wird nicht auf den eigenen Splits, sondern einerseits nach
Abbildung auf die vier groben CoNLL-Typen auf Gold-Benchmarks wie
CoNLL-2003, andererseits auf 1\,000 englischen Sätzen, deren automatische
Annotation manuell geprüft wurde.
Auf dieser Stichprobe stimmt die Silberreferenz auf Tokenebene mit
einem F$_1$ von 0{,}85 mit der manuellen Korrektur überein; PER, ORG,
LOC, CEL, DIS, EVE und MEDIA zählen zu den zuverlässig annotierten
Klassen, während ANIM, BIO, FOOD und PLANT häufig verwechselt oder nicht
verlinkt werden.
Für die offiziellen Splits, die diese Arbeit verwendet
(Abschnitt~\ref{sec:datenauswahl}), berichtet die Originalarbeit somit
keine Referenzwerte.
Veröffentlichte Werte auf MultiNERD finden sich in der hier gesichteten
Literatur nur in Benchmarks instruktionsgetunter LLMs
(Abschnitt~\ref{sec:llm_ner_related}) und in der unten behandelten
GLiNER-Arbeit; sie verwenden eigene Datenaufbereitungen, oder ihr
Testumfang lässt sich nicht auf die offiziellen Splits abbilden.

Wie leistungsfähig Encoder bei feingranularer NER sind, zeigt der
MultiCoNER-II-Wettbewerb mit 33 feingranularen Typen in zwölf Sprachen.
\textcite{tan2023damo} gewinnen dort mit U-RaNER neun der 13 Tracks.
Das System kombiniert XLM-RoBERTa-large mit einer CRF-Schicht und einem
Retrieval-Modul, das zu jedem Eingabesatz Textpassagen aus Wikipedia
sowie Entitäten und Typinformationen aus Wikidata abruft und dem Encoder
als zusätzlichen Kontext mitgibt; auf dem englischen Track erreicht es
einen Makro-F$_1$ von 85{,}53.
Ergänzend prompten die Autoren ChatGPT (\texttt{gpt-3.5-turbo}) in drei
Varianten auf dem Validierungssplit.
Selbst die beste Variante mit Demonstrationsbeispielen erreicht im Mittel
über die Sprachen nur einen Mikro-F$_1$ von 14{,}78, gegenüber 92{,}77
für das eigene System auf demselben Split~\cite[Tab.~4]{tan2023damo}.
Der Befund bezieht sich auf ein Modell ohne jede aufgabenspezifische
Anpassung und auf ein sehr feines Typinventar; er zeigt aber, wie groß
der Abstand zwischen reinem Prompting und überwachten Encodern bei
feingranularer NER ausfallen kann.

Eine breitere Einordnung liefert die Vergleichsstudie von
\textcite{keraghel2024recent}, die als Teil eines Surveys zehn
NER-Datensätze aus allgemeinen und spezialisierten Domänen mit mehreren
Frameworks und Modellen auswertet, darunter BERT, XLM-RoBERTa-large,
DeBERTa-v3-base und DeBERTa-v3-large, ein feinjustiertes GLiNER-Modell
sowie GPT-4o mit Beispielentitäten im Prompt.
MultiNERD ist nicht enthalten.
Unter den Transformer-Modellen erweisen sich die DeBERTa-Varianten als am
effektivsten; nach einem Nemenyi-Test teilen sich DeBERTa-v3-large und
das feinjustierte GLiNER-Modell die besten Ränge, ohne sich signifikant
zu unterscheiden.
GPT-4o bleibt auf neun der zehn Datensätze deutlich hinter den
feinjustierten Encodern zurück, auf CoNLL-2003 etwa mit einem
Makro-F$_1$ von 65{,}17 gegenüber 90{,}98 für
DeBERTa-v3-large~\cite[Tab.~5]{keraghel2024recent}.
Auf WNUT2017 liegt GPT-4o mit 44{,}73 nahezu gleichauf mit
DeBERTa-v3-large (44{,}81), bleibt aber auch dort hinter dem
feinjustierten GLiNER-Modell zurück.
Die Autoren folgern, dass LLMs allein für NER nur eingeschränkt geeignet
sind und eine Kombination mit feinjustierten Modellen vielversprechender
erscheint.
Die Studie stützt die Wahl von DeBERTa-v3 als Encoder-Referenz in dieser
Arbeit (Abschnitt~\ref{sec:encoder-auswahl}); sie liegt bislang nur als
Preprint vor.

Eine zweite Entwicklungslinie nutzt Encoder als Generalisten, die
Entitätstypen zur Laufzeit aus natürlichsprachlichen Bezeichnungen
ableiten.
\textcite{zaratiana2024gliner} stellen mit GLiNER ein Modell auf Basis
von DeBERTa-v3 vor, das Typbezeichnungen und Eingabetext in einer
gemeinsamen Sequenz verarbeitet, Kandidatenspans bis zu einer
Maximallänge repräsentiert und jeden Span über ein Skalarprodukt mit den
Typrepräsentationen bewertet.
Trainiert wird auf Pile-NER, den von ChatGPT erzeugten Annotationen, die
auch UniversalNER verwendet (Abschnitt~\ref{sec:llm_ner_related}).
Im Zero-Shot-Vergleich auf sieben Datensätzen außerhalb der
Trainingsverteilung, darunter die fünf CrossNER-Domänen, erreicht die
mittlere Variante mit rund 90~Mio. Parametern im Mittel einen F$_1$ von
55{,}4 und damit das Niveau des etwa 140-mal größeren UniNER-13B; die
größte Variante liegt im Mittel rund fünf Punkte darüber, und selbst die
kleinste Variante mit 50~Mio. Parametern übertrifft ChatGPT.
Das ebenfalls verglichene GoLLIE-7B ordnet sich zwischen der mittleren
und der größten Variante ein.
Inferenzzeiten messen die Autoren nicht; sie verweisen auf die parallele
Span-Klassifikation gegenüber der sequenziellen Generierung von LLMs und
auf die Lauffähigkeit auf CPUs.

\textcite{bogdanov2024nuner} verfolgen mit NuNER einen verwandten Weg,
setzen aber auf ein Pre-Training für die anschließende Feinjustierung
statt auf das Zero-Shot-Regime.
Ein RoBERTa-base-Modell wird kontrastiv auf 4{,}38~Mio.
Entitätsannotationen vortrainiert, die GPT-3.5 für Sätze aus dem
C4-Korpus über rund 200\,000 verschiedene Konzepte erzeugt hat.
Zwei Vergleiche mit LLMs sind dabei zu unterscheiden.
Gegenüber UniNER-7B mit 56-mal mehr Parametern zeigt NuNER im
Few-Shot-Fine-Tuning auf den vollständigen Testsets von vier Datensätzen
(Mikro-F$_1$ auf Entitätsebene) ähnliche
Transferfähigkeiten~\cite[Tab.~2]{bogdanov2024nuner}.
GPT-3.5 und GPT-4 per In-Context-Learning werden nur auf zwei
Datensätzen mit auf 1\,000 Beispiele verkleinerten Testsets und
Makro-F$_1$ auf Tokenebene verglichen; NuNER übertrifft dort GPT-3.5 bei
größeren Trainingsmengen, während die Autoren das Verhältnis zu GPT-4
selbst als unklar bezeichnen~\cite[Abschn.~6]{bogdanov2024nuner}.

Zusammengenommen zeigen diese Arbeiten zweierlei.
Erstens bleiben vollständig feinjustierte Encoder im überwachten Setting
die Referenz, auch bei feingranularen Typinventaren, und per Prompting
genutzte LLMs ohne Anpassung liegen in den hier berichteten Vergleichen deutlich
zurück~\cite{tan2023damo,keraghel2024recent}.
Zweitens lässt sich Wissen aus LLMs über erzeugte Trainingsdaten in
Encoder übertragen, sodass Modelle mit einem Bruchteil der Parameter mit
LLMs im Bereich von 7 bis 13~Mrd. Parametern
gleichziehen~\cite{zaratiana2024gliner,bogdanov2024nuner}.
Diese Gleichstände beruhen jedoch auf Zero- und Few-Shot-Benchmarks, bei
GLiNER unter anderem auf den domänenspezifischen Typinventaren von
CrossNER; keine der Arbeiten stellt vollständig feinjustierte Encoder
und feinjustierte generative LLMs in derselben Pipeline auf den
offiziellen MultiNERD-Splits mit gemeinsamer Ressourcenmessung
gegenüber.
Für MultiNERD kommt hinzu, dass die Originalarbeit keine Referenzwerte
für die offiziellen Splits liefert und die Testreferenz selbst ein Silver
Standard ist.
Vergleiche auf diesem Datensatz müssen daher innerhalb einer gemeinsamen
Pipeline erfolgen und absolute Werte als Übereinstimmung mit einer
automatisch erzeugten Referenz lesen (Abschnitt~\ref{sec:datenauswahl}).

\section{LLM-basierte Ansätze für NER}
\label{sec:llm_ner_related}

Wird NER als Textgenerierung formuliert, sind zwei Entwurfsentscheidungen
zu treffen, die in der Literatur getrennt untersucht werden: das
Ausgabeformat, in dem das Modell Entitäten benennt, und die Art der
Anpassung, von reinem Prompting über Destillation bis zum Fine-Tuning.
Dieser Abschnitt folgt dieser Zweiteilung.
Zunächst werden Ausgabestrategien und ihre Wirkung auf Erkennungsleistung
und Ausgabelänge dargestellt, woraus die Begründung der in dieser Arbeit
gewählten Strategie folgt.
Anschließend werden Ansätze ohne aufgabenspezifisches Training und
Ansätze mit Destillation oder Fine-Tuning vorgestellt.

\textcite{semin2026spanLabeling} systematisieren die Ausgabestrategien
für Span-Labeling mit LLMs in drei Familien.
Beim \emph{Tagging} wird der Eingabetext mit Markierungen reproduziert,
sodass die Ausgabelänge linear mit der Eingabelänge wächst.
Beim \emph{Indexing} gibt das Modell für jeden Span numerische Start- und
Endindizes aus, sodass die Ausgabelänge nur von der Zahl der Spans
abhängt.
Beim \emph{Matching} erzeugt das Modell den Text der Spans, der
anschließend per Zeichenkettenabgleich in der Eingabe lokalisiert wird.
Die Autoren vergleichen die Strategien mit vier per Prompting genutzten
offenen Modellen zwischen 8 und 70~Mrd. Parametern unter anderem auf dem
mehrsprachigen UNER-Benchmark mit den Typen Person, Organisation und
Ort.
Dieser Benchmark ist der Datensatz Universal NER von
\textcite{mayhew2024uner}, nicht das UniversalNER-Modell von
\textcite{zhou2024universalner}.
Tagging erweist sich als robusteste, aber tokenintensivste Strategie;
Matching ist konkurrenzfähig und tokeneffizient; Indexing mit
Zeichenindizes ist nur dann konkurrenzfähig, wenn die Eingabe selbst mit
Zeichenindizes angereichert wird, weil die Modelle Zeichenpositionen
sonst nur indirekt schätzen können.
Mit LogitMatch schlagen die Autoren zudem ein Constrained-Decoding-Verfahren
vor, das Matching-Ausgaben auf gültige Eingabespans beschränkt und so
Fehler beim Zeichenkettenabgleich beseitigt.
Eine auf \textcite[Tab.~9]{ramponi2025fallacy} zurückgehende
Indexing-Variante mit nummerierten Wörtern statt Zeichenpositionen wird
beschrieben, aber nicht evaluiert.
Die Arbeit liegt bislang nur als Preprint vor.

Dass das Format auch bei feinjustierten Modellen ein Faktor erster
Ordnung ist, zeigen \textcite{zhan2026generativeNER}.
Sie trainieren acht offene LLMs mit LoRA auf fünf Ausgabeformaten: zwei
Inline-Formate, die Entitäten mit Klammern oder XML-Tags im
reproduzierten Satz markieren, und drei JSON-Formate, die Entitäten nach
Typ gruppieren, mit einem Vorkommensindex versehen oder mit Zeichenoffsets
angeben.
Die Inline-Formate erreichen über alle Modelle die besten Mittelwerte,
auf den flachen Datensätzen rund 92 F$_1$, während JSON mit Zeichenoffsets
dort im Mittel auf rund 54 fällt; die Autoren führen dies darauf zurück,
dass das Berechnen von Zeichenpositionen der generativen Arbeitsweise der
Modelle schlecht entspricht.
\textcite{dagdelen2024structuredIE} berichten für die Extraktion aus
materialwissenschaftlichen Texten mit feinjustiertem GPT-3 und einem
LoRA-angepassten Llama-2-Modell, dass JSON-Ausgaben in rund 99\,\% der
Fälle parsebar sind, die Formatierung aber nicht erzwungen wird und JSON
mehr Tokens benötigt als satzförmige Schemata.
Auf die Kosten der Ausgabelänge weisen \textcite{lu2024padellm} hin: Die
Latenz generativer NER entsteht vor allem durch die sequenzielle
Dekodierung langer Ausgaben, und Inline-Formate verlängern die Ausgabe,
weil sie den gesamten Eingabetext wiederholen.
Ihr paralleles Decoding, das Erwähnungen je Typ gleichzeitig statt
nacheinander erzeugt, verkürzt die Inferenz um das 1{,}76- bis
10{,}22-Fache bei unveränderter Erkennungsleistung.

Vor diesem Hintergrund lässt sich die in Abschnitt~\ref{sec:prompt-design}
beschriebene Ausgabestrategie dieser Arbeit einordnen.
Das Modell erhält eine nummerierte Tokenliste und gibt ein JSON-Array mit
Start- und Endindex, Oberflächentext und Typ je Entität aus.
Das Format gehört zur Indexing-Familie, allerdings in der angereicherten
Variante auf Wortebene: Die nummerierten Eingabetokens und die
Wortindizes in der Ausgabe folgen \textcite{ramponi2025fallacy}, die
Einordnung in die drei Familien \textcite{semin2026spanLabeling};
JSON-Schema und Nachverarbeitung sind eigene Umsetzung dieser Arbeit.
Daraus ergeben sich drei erwartete Vorteile.
Erstens stehen die Indizes explizit in der Eingabe, sodass das Modell
keine Zeichen zählen muss, was den von \textcite{zhan2026generativeNER}
beobachteten Einbruch bei Zeichenoffsets vermeiden soll.
Zweitens hängt die Ausgabelänge von der Zahl der Entitäten, der Länge
ihrer Spantexte und der Serialisierung ab; anders als bei Tagging- und
Inline-Formaten entfällt lediglich die Reproduktion der
Nicht-Entitätstokens.
Drittens erlauben die Indizes eine eindeutige Rückabbildung auch bei
mehrfach vorkommenden Oberflächenformen, die beim Matching eine
Fehlerquelle sind.
Das redundante Textfeld ermöglicht die Konsistenzprüfung in der
Nachverarbeitung (Abschnitt~\ref{sec:postprocessing}).
Einschränkend ist festzuhalten, dass \textcite{semin2026spanLabeling} die
Wortindex-Variante nicht evaluiert haben und ihre Befunde auf
per Prompting genutzten, nicht feinjustierten Modellen beruhen; ob die Vorteile der
angereicherten Indizierung auch im QLoRA-Regime und bei 15 Typen tragen,
ist eine empirische Frage dieser Arbeit.
Auf Constrained Decoding wird verzichtet, damit die Formtreue der
Ausgaben selbst messbar bleibt (Abschnitt~\ref{sec:robustheit-ausgaben}).

Die frühen LLM-Ansätze für NER lassen die Gewichte des Modells
unverändert und passen es nicht an die Aufgabe an.
\textcite{wang2025gptner} formulieren mit GPT-NER die Aufgabe als
Generierung des Eingabesatzes, in dem Entitäten eines Typs durch die
Marker \texttt{@@} und \texttt{\#\#} eingefasst sind, und wählen
Demonstrationsbeispiele über eine Nachbarschaftssuche auf
Entitätsrepräsentationen aus.
Diese Repräsentationen liefert in der stärksten Variante ein
feinjustierter NER-Tagger~\cite[Abschn.~3.2.2]{wang2025gptner}.
Das verwendete GPT-3-Modell neigt dazu, auch solche Eingaben
überkonfident mit Entitäten zu versehen, die gar keine Entität des
abgefragten Typs enthalten; die Autoren bezeichnen diese Fälle als
\emph{NULL inputs} und ergänzen eine Selbstverifikation, in der das
Modell jede extrahierte Entität mit Ja oder Nein bestätigt.
Auf CoNLL-2003 erreicht GPT-NER so einen F$_1$ von 90{,}91 gegenüber
92{,}8 für einen feinjustierten BERT-Tagger~\cite[Tab.~1]{wang2025gptner};
bei verschachtelten Datensätzen ist der Abstand größer.
\textcite{xie2023zeroshot} untersuchen Zero-Shot-NER mit ChatGPT und
zerlegen die Aufgabe in eine Abfrage je Entitätstyp, ergänzt um
syntaktische Hinweise im Prompt und ein zweistufiges Mehrheitsvotum über
mehrere Antworten.
Die Zerlegung verbessert die Ergebnisse auf allen sieben untersuchten
Benchmarks, vervielfacht aber die Zahl der Modellaufrufe je Satz; die in
den Tabellen der Arbeit mitgeführten überwachten Bestwerte liegen
weiterhin weit über den Zero-Shot-Ergebnissen.

\textcite{ashok2023promptner} erweitern das Prompting mit PromptNER um
textuelle Typdefinitionen und verlangen vom Modell eine Kandidatenliste
mit Begründungen, warum ein Kandidat zur Definition passt oder nicht.
Mit GPT-4 und wenigen Demonstrationsbeispielen verbessert PromptNER in
der zweiten Fassung der Arbeit die besten bisherigen
Few-Shot-Ergebnisse auf CoNLL-2003 um rund vier, auf GENIA um rund neun
und auf FewNERD um rund vier F$_1$-Punkte.
Im Cross-Domain-Setting mit nur zwei Beispielen aus der Zieldomäne
erreicht es auf drei der fünf CrossNER-Domänen neue Bestwerte, etwa
84{,}26 auf \emph{Music}, obwohl weniger als 2\,\% der verfügbaren
Trainingsdaten verwendet werden; in den Domänen \emph{AI} und
\emph{Science} bleiben die bisherigen Verfahren vorn.

Einen anderen Weg geht CodeIE von \textcite{li2023codeie}: NER und
Relationsextraktion werden als Vervollständigung von Python-Code
formuliert und mit dem Code-Modell Codex gelöst.
Im Few-Shot-Setting übertrifft CodeIE sowohl natürlichsprachliche Prompts
als auch feinjustierte UIE-Modelle, die nur dieselben wenigen Beispiele
sehen; auf CoNLL-2003 steht einem Wert von 82{,}32 mit fünf Beispielen je
Typ jedoch ein vollständig überwachter UIE-Wert von 92{,}99 gegenüber.

Eine zweite Gruppe passt offene LLMs gezielt an NER an.
\textcite{zhou2024universalner} lassen ChatGPT Passagen aus dem
Pile-Korpus mit Entitäten frei gewählter Typen annotieren und trainieren
damit LLaMA-Modelle mit 7 und 13~Mrd. Parametern als UniversalNER.
Auf einem Benchmark aus 43 Datensätzen in neun Domänen übertreffen die
Studentenmodelle im Zero-Shot-Regime ihren Lehrer deutlich, mit
mittleren F$_1$-Werten von 41{,}7 bzw. 43{,}4 gegenüber 34{,}9 für
ChatGPT.
Das Training folgt dem Rezept von Vicuna; Adapterverfahren werden nicht
erwähnt, Ressourcen nicht gemessen.

In einem zusätzlichen überwachten Setting wird UniNER-7B auf den
Trainingsdaten von 20 Datensätzen weiter feinjustiert und erreicht im
Mittel 84{,}78 gegenüber 80{,}09 für BERT-base; die BERT-Werte sind aus
der InstructUIE-Arbeit~\cite{wang2023instructuie} übernommen und nicht neu
gemessen.
Zu diesen 20 Datensätzen gehört auch eine Variante von MultiNERD, auf
der UniNER-7B 93{,}73 und BERT-base 91{,}25
erreichen~\cite[Tab.~2]{zhou2024universalner}.
Da InstructUIE eigene Datenaufbereitungen verwendet, lassen sich diese
Werte nicht auf die offiziellen MultiNERD-Splits übertragen; sie zeigen
aber, dass auf diesem Datensatz sowohl Encoder als auch feinjustierte
LLMs hohe Werte erreichen und der Abstand zwischen beiden klein ist.

\textcite{sainz2024gollie} verfolgen mit GoLLIE das Ziel,
Annotationsrichtlinien explizit nutzbar zu machen.
Code-LLaMA-Modelle mit 7, 13 und 34~Mrd. Parametern werden mit QLoRA auf
Informationsextraktionsaufgaben feinjustiert, deren Schemata samt
Richtlinien als Python-Klassen im Prompt stehen.
Im Zero-Shot-Regime auf nicht im Training enthaltenen Datensätzen hebt
das Richtlinientraining den mittleren F$_1$ gegenüber einer Baseline mit
demselben Backbone ohne Richtlinien (und ohne die zugehörigen
Kandidatenlisten und Regularisierung) um rund 13 Punkte auf 55{,}3 für
das 7B-Modell; die größeren Modelle legen nur wenig zu.
Im überwachten Setting dagegen unterscheiden sich GoLLIE und Baseline
kaum.
Auf den fünf CrossNER-Domänen erreicht GoLLIE-7B Werte zwischen 55{,}5
(\emph{Science}) und 67{,}8 (\emph{Music}) und übertrifft damit das bis
dahin beste Zero-Shot-System InstructUIE in jeder
Domäne~\cite[Tab.~3]{sainz2024gollie}.

Auch MultiNERD ist Teil der Zero-Shot-Evaluation von GoLLIE.
Das 7B-Modell erreicht dort einen F$_1$ von 76{,}0 und liegt damit
deutlich über der Baseline ohne Richtlinien~\cite[Tab.~3]{sainz2024gollie},
wobei acht der 15 Typen in anderer Benennung bereits im Training
vorkamen.
Die Arbeit gibt an, für alle Datensätze die Standardsplits zu verwenden,
und nennt für MultiNERD 32\,908 Testbeispiele, genau das Doppelte der
16\,454 Sätze des englischen Testsplits dieser Arbeit
(Abschnitt~\ref{sec:datenaufteilung}); wie beide Angaben zusammenhängen,
bleibt ungeklärt, sodass der Wert nur eingeschränkt vergleichbar
ist~\cite[Anh.~D.2, Tab.~9]{sainz2024gollie}.

Einen weiteren Weg zur Zero-Shot-Generalisierung geht GNER von
\textcite{ding2024gner}.
Flan-T5-Modelle bis 11~Mrd. Parameter und LLaMA-7B lernen auf Pile-NER,
jedes Eingabewort mit seinem Label zu wiederholen, also auch
Nicht-Entitäten als negative Instanzen mitzuführen, was nach den Autoren
Kontext liefert und Labelgrenzen schärft.
Ein Abgleich über die längste gemeinsame Teilfolge ordnet die Ausgabe
auch bei ausgelassenen oder hinzugefügten Wörtern der Eingabe zu.
Im Zero-Shot-Regime auf den CrossNER- und MIT-Datensätzen berichten die
Autoren einen Gewinn von 9 F$_1$-Punkten gegenüber dem bisherigen Stand
der Technik.
Verglichen wird mit Werten anderer Systeme, darunter GLiNER, die teils
übernommen und teils mit veröffentlichten Checkpoints neu evaluiert
wurden; gemessen werden Inferenzdurchsatz und gesamte GPU-Stunden des
Trainings, aber kein Encoder in derselben Pipeline.
Das Label-je-Wort-Format reproduziert den gesamten Eingabetext und
gehört damit zur Tagging-Familie, anders als das hier verwendete
Indexing-Format.

Über die Einzelarbeiten hinaus ergeben sich drei Beobachtungen.
Erstens ist das Ausgabeformat kein Implementierungsdetail, sondern
bestimmt Erkennungsleistung und Ausgabelänge zugleich; Formate, die den
Eingabetext reproduzieren, sind robust, aber teuer, Formate mit
Zeichenpositionen
fehleranfällig~\cite{semin2026spanLabeling,zhan2026generativeNER,lu2024padellm}.
Zweitens bleiben per Prompting genutzte LLMs ohne Anpassung der Gewichte
auf etablierten Benchmarks hinter überwachten Encodern zurück, selbst
wenn wie bei GPT-NER ein feinjustierter Tagger die Demonstrationen
auswählt; ihre Stärke zeigen sie dort, wo wenige Beispiele vorliegen
oder Domänen
wechseln~\cite{wang2025gptner,xie2023zeroshot,ashok2023promptner,li2023codeie}.
Drittens verringern Destillation und Fine-Tuning diesen Abstand, doch die
Vergleiche mit Encodern stützen sich auf übernommene Werte oder auf
Zero-Shot-Benchmarks~\cite{zhou2024universalner,sainz2024gollie}.
UniversalNER misst keinen Trainings- oder Inferenzaufwand; GoLLIE
dokumentiert Trainingszeit und CO$_2$ je Größe (7B: eine A100,
44{,}5~h; 13B: 79{,}5~h; 34B: zwei A100, 94{,}6~h), aber keine
Inferenzlatenz und keinen
Encodervergleich~\cite[Anh.~D.3, Tab.~10]{sainz2024gollie}.
Für den CrossNER-Ergänzungsversuch dieser Arbeit
(Abschnitt~\ref{sec:crossner}) liefern PromptNER und GoLLIE Orientierung,
aber keine Vergleichswerte: PromptNER sieht zwei Beispiele aus der
Zieldomäne, GoLLIE trainiert auf den Richtlinien vieler Datensätze, und
beide werten die domänenspezifischen Typinventare von CrossNER aus.
Der Versuch dieser Arbeit ist dagegen ein reiner Transfertest ohne
Zielbeispiele, beschränkt auf die Grobtypen PER, ORG und LOC.
Die berichteten Werte zeigen deshalb lediglich, dass der Domänenwechsel
auch für generative Modelle eine erhebliche Hürde bleibt, nicht aber,
welches Niveau auf dem reduzierten Typinventar zu erwarten ist.

\section[PEFT-angepasste LLMs im direkten Vergleich mit Encodern]{PEFT-angepasste LLMs im direkten \mbox{Vergleich} mit Encodern}
\label{sec:peft_ner_related}

LoRA und QLoRA wurden in Abschnitt~\ref{sec:lora} als Verfahren
eingeführt.
Dieser Abschnitt behandelt Arbeiten, die generative LLMs mit solchen
Adaptern an NER anpassen und in derselben Studie mit Encodern
vergleichen.
Soweit die hier gesichtete Literatur reicht, ist diese Kombination noch
selten untersucht; die vier im Folgenden dargestellten Arbeiten decken
allgemeine und domänenspezifische Datensätze ab und unterscheiden sich
darin, welche Ressourcen sie erfassen.

Den breitesten Vergleich auf allgemeinen Benchmarks legen
\textcite{zhan2026generativeNER} vor.
Acht offene LLMs der Familien LLaMA, Qwen, Granite und Gemma mit 3 bis
9~Mrd. Parametern werden mit LoRA auf CoNLL-2003 und OntoNotes 5.0 sowie
auf den verschachtelten Datensätzen ACE2005 und GENIA feinjustiert und
spezialisierten Modellen wie BERT-MRC gegenübergestellt.
Im besten Format erreicht LLaMA3.1-8B auf CoNLL-2003 einen F$_1$ von
93{,}85 gegenüber 93{,}04 für BERT-MRC~\cite[Tab.~3]{zhan2026generativeNER};
auf GENIA liegt das beste LLM 4{,}02 Punkte hinter dem besten
spezialisierten Modell.
Für den Kostenvergleich trainieren die Autoren LLaMA3.1-8B und BERT-MRC
selbst auf einer einzelnen L40-GPU und berichten dort dieselben
F$_1$-Werte wie in der Haupttabelle~\cite[Anh.~G, Tab.~13]{zhan2026generativeNER};
für die übrigen Baselines nennt die Arbeit keine Herkunft der Werte.
Eine Fehleranalyse für LLaMA3.1-8B und BERT-MRC zeigt unterschiedliche
Profile: Der Encoder übersieht vor allem Entitäten, das LLM weist ihnen
vor allem falsche Typen zu.
Ein Kontrollexperiment, in dem die Typbezeichnungen auf CoNLL-2003 und
CleanCoNLL durch bedeutungslose Symbole ersetzt werden, verändert die
Ergebnisse kaum, was die Autoren gegen eine bloße Memorisierung von
Entität-Typ-Paaren anführen.
Zum Ressourcenaufwand enthält die aktuelle Fassung nur einen Vergleich
der Trainingskosten von LLaMA3.1-8B und BERT-MRC auf einer einzelnen GPU,
bei dem Trainingszeit und Spitzenspeicher des LLM moderat höher
ausfallen; Inferenzlatenz und Decoding-Einstellungen werden nicht
berichtet.
Die Arbeit ist zur Veröffentlichung in den Findings der EMNLP 2026
angenommen.

Den Ressourcenaspekt stellen \textcite{obeidat2025llm} in den
Mittelpunkt.
Für fünf biomedizinische Datensätze, darunter ein Social-Media-Korpus zu
Substanzkonsum, vergleichen sie vollständig feinjustierte Encoder mit
rund 340 bis 435~Mio. Parametern (BERT-large, BiomedBERT-large und
DeBERTa-v3-large) mit Mistral-7B-Instruct und Llama-3.1-8B-Instruct, die
mit QLoRA und 4-Bit-Quantisierung angepasst werden; der LoRA-Rang ist so
gewählt, dass die trainierbaren Parameter der LLMs etwa der Größe der
Encoder entsprechen.
Beide Modellklassen geben IOB2-Tags aus und werden strikt auf
Entitätsebene bewertet.
Die LLMs liegen auf vier Datensätzen um 2 bis 8 F$_1$-Punkte vorn, auf
dem fünften knapp (0{,}5 Punkte) hinter dem besten
Encoder~\cite[Tab.~III]{obeidat2025llm}; der Vorsprung
ist bei Entitäten ab drei Tokens am größten und geht vor allem auf eine
höhere Precision zurück.
Dem stehen Inferenzzeiten gegenüber, die je Testbeispiel um das 40- bis
220-Fache über denen des besten Encoders
liegen~\cite[Tab.~IV]{obeidat2025llm}; die LLMs benötigten zudem zwei
statt einer H100-GPU.
Die Autoren führen die Latenz unter anderem darauf zurück, dass das
IOB2-Format für jedes Token ein Tag generieren lässt und lange
Typbezeichnungen in mehrere Tokens zerfallen.

Dass LLMs für tokenbasierte NER ohne Anpassung kaum geeignet sind, zeigen
\textcite{lu2024llmstruggle} am klinischen Korpus RareDis-v1 mit fünf
Typen zu seltenen Krankheiten.
Lokale LLMs mit 7~Mrd. Parametern sowie GPT-3.5 und GPT-4 werden
aufgefordert, IOB2-Tags zu erzeugen.
Im Zero-Shot-Regime erreicht LLaMA-2-7B für den Typ \emph{Rare Disease}
einen F$_1$ von 0{,}149 und GPT-4 einen von 0{,}576.
Nach einem LoRA-Fine-Tuning nähert sich das medizinisch vortrainierte
Llama2-MedTuned-7B nach Darstellung der Autoren dem feinjustierten
BioClinicalBERT an, bleibt aber auf allen Typen darunter; die Autoren
vermuten, dass vollständiges Fine-Tuning den Abstand schließen könnte.
Diese Vermutung relativieren \textcite[Tab.~11]{zhan2026generativeNER}:
In ihren Prüffällen auf CoNLL-2003 und GENIA liegt LoRA um 0{,}14 bis
0{,}48 F$_1$-Punkte über vollständigem Fine-Tuning, was für Qwen3.5 und
MultiNERD jedoch keine Gleichwertigkeit belegt.
Als Gründe für die Schwierigkeiten nennen sie, dass tokenbasierte NER auf
die Repräsentation einzelner Tokens angewiesen ist, dass Decoder im
Pre-Training kein auf Repräsentationen ausgerichtetes Ziel verfolgen und
dass IOB2-Tags für die Modelle ungesehene Symbole sind.
Die Modellgröße variiert die Arbeit nicht.

Einen konstruktiven Umgang mit genau diesem Unterschied zeigen
\textcite{dukic2024lookingright}.
Sie setzen auf Llama-2-7B und Mistral-7B einen klassischen
Klassifikationskopf für Sequence Labeling und entfernen während eines
QLoRA-Fine-Tunings die kausale Aufmerksamkeitsmaske schichtweise, sodass
die Decoder bidirektional auf den Kontext zugreifen können.
Auf CoNLL-2003 steigt der F$_1$ von Llama-2-7B durch die Entfernung der
Maske in allen Schichten von 0{,}744 auf 0{,}908, und beide Decoder
übertreffen RoBERTa-large mit 0{,}900; bei einem Modell mit nur 68~Mio.
Parametern bringt die Maskenentfernung dagegen keinen Gewinn.
Die Decoder werden hier also nicht generativ, sondern wie Encoder
eingesetzt, was die Formatfrage umgeht, aber eine Suche über
Maskierungskonfigurationen erfordert.

Drei der vier Arbeiten konvergieren in einem Punkt: Mit Adaptern
feinjustierte LLMs erreichen auf etablierten Benchmarks das Niveau
feinjustierter Encoder oder übertreffen es
leicht~\cite{zhan2026generativeNER,obeidat2025llm,dukic2024lookingright};
\textcite{lu2024llmstruggle} finden im klinischen Setting mit LoRA
dagegen einen bleibenden Rückstand.
Generativ arbeiten dabei nur \textcite{zhan2026generativeNER} und
\textcite{obeidat2025llm}, während \textcite{dukic2024lookingright} die
Decoder mit Klassifikationskopf einsetzen; den Mehraufwand bei der
Inferenz beziffern allein \textcite{obeidat2025llm}, und dort fällt der
Leistungsgewinn im Verhältnis dazu klein aus.
Zugleich markieren sie die Grenzen des bisherigen Wissens.
Die Datensätze sind allgemeine Benchmarks wie CoNLL-2003 mit vier und
OntoNotes mit 18 Typen oder Fachkorpora; ein genreübergreifender
Datensatz mit 15 Typen wie MultiNERD ist nicht darunter.
Die generativen Ausgabeformate wachsen mit der Eingabelänge: Die
IOB2-Ausgaben bei \textcite{obeidat2025llm} führen jedes Eingabetoken
mit seinem Tag auf, die Inline-Markierungen bei
\textcite{zhan2026generativeNER} wiederholen den Satz; beides erhöht
nach den Befunden aus Abschnitt~\ref{sec:llm_ner_related} die Latenz.
Die Ressourcenerfassung ist unvollständig: \textcite{obeidat2025llm}
messen Inferenzzeiten, aber keinen Speicher, \textcite{zhan2026generativeNER}
berichten nur Trainingskosten, und \textcite{lu2024llmstruggle} sowie
\textcite{dukic2024lookingright} messen keine Ressourcen.
Die untersuchten LLMs liegen durchweg zwischen 3 und 9~Mrd. Parametern,
sodass weder deutlich kleinere noch deutlich größere Decoder im Vergleich
stehen.
Und die Encoder-Referenzwerte stammen teilweise aus anderen Publikationen
statt aus derselben Pipeline.
Diese Punkte bilden den Ausgangspunkt der Forschungslücke in
Abschnitt~\ref{sec:research_gap}.

\section{Einfluss der Modellgröße in verwandten Arbeiten}
\label{sec:model_size_related}

Die in Abschnitt~\ref{sec:model_size} dargestellten Skalierungsgesetze
beschreiben den Pre-Training-Verlust in Abhängigkeit von Parameterzahl,
Datenmenge und Rechenaufwand.
Welche Rolle die Modellgröße für eine nachgelagerte Aufgabe wie NER
spielt, lässt sich daraus nicht ableiten.
Die zuvor vorgestellten Arbeiten erlauben hierzu nur indirekte Aussagen,
weil die Modellgröße fast immer gemeinsam mit anderen Faktoren variiert.

Innerhalb einer LLM-Familie sind Größeneffekte sichtbar, aber klein.
UniversalNER gewinnt beim Übergang vom 7B- zum 13B-Studenten im
Zero-Shot-Regime leicht hinzu, was die Autoren als bessere
Generalisierung größerer Modelle deuten~\cite{zhou2024universalner}.
GoLLIE zeigt von 7 über 13 bis 34~Mrd. Parametern im Mittel steigende
Zero-Shot-Werte, wobei einzelne Datensätze stark, andere gar nicht
profitieren; der Vergleich der berichteten Mittelwerte zeigt zudem, dass
der Gewinn durch die Vergrößerung des Backbones deutlich hinter dem
Gewinn durch das Richtlinientraining zurückbleibt, und auf MultiNERD
liegt das 13B-Modell nicht hinter dem 34B-Modell~\cite{sainz2024gollie}.
\textcite{zhan2026generativeNER} unterscheiden eine leichte und eine
mittlere Größenklasse, ziehen aber keinen Skalierungsschluss.
Zwischen den Modellklassen fallen die Vergleiche noch uneindeutiger aus:
GLiNER und NuNER erreichen mit Encodern weit unterhalb einer Milliarde
Parameter das Niveau von LLMs mit 7 bis 13~Mrd.
Parametern~\cite{zaratiana2024gliner,bogdanov2024nuner}.
Dabei teilen nur GLiNER und UniNER die Pile-NER-Trainingsdaten, NuNER
nutzt GPT-3.5-Annotationen des C4-Korpus, sodass der Größenunterschied
von Architektur, Trainingsziel und Daten überlagert wird.
Bei \textcite{obeidat2025llm} ist die Zahl trainierbarer Parameter über
den LoRA-Rang auf eine ähnliche Größenordnung wie die Encoder abgestimmt
(rund 341 bzw. 353~Mio. Adapterparameter gegenüber 336 bis 435~Mio.),
und dennoch liegen die LLMs vorn.
Dieser Vorsprung dürfte dann eher auf die Kapazität des eingefrorenen
Basismodells als auf die Zahl angepasster Parameter zurückgehen.
Bei Encodern selbst ist die größere Variante nicht durchgehend besser:
In der Vergleichsstudie von \textcite{keraghel2024recent} liegt
DeBERTa-v3-base auf CoNLL-2003 und NCBI Disease knapp vor
DeBERTa-v3-large, auf den meisten anderen Datensätzen ist es umgekehrt.
Und \textcite{dukic2024lookingright} beobachten, dass ihre
Maskenentfernung erst ab einer gewissen Modellgröße wirkt.

Für die vorliegende Arbeit folgt daraus zweierlei.
Erstens wird die Modellgröße innerhalb jeder Familie getrennt betrachtet:
auf der Encoder-Seite mit DeBERTa-v3-base (rund 184~Mio. Parameter) und
DeBERTa-v3-large (rund 434~Mio.), auf der LLM-Seite mit Qwen3.5-0.8B,
-4B und -27B, deren Textmodelle rund 0{,}75, 4{,}21 und 26{,}90~Mrd.
Parameter umfassen (Abschnitt~\ref{sec:groessenklassen}).
Die Spanne reicht damit vom kleinsten Encoder bis zum größten LLM über
mehr als zwei Größenordnungen und deckt auf der LLM-Seite sowohl ein
Modell unterhalb als auch ein Modell oberhalb des in den direkten
Vergleichen der Literatur abgedeckten Größenbereichs ab.
Zweitens bleibt die Auswertung deskriptiv: Mit zwei Encoder- und drei
LLM-Größen, größenabhängig gewählten Hyperparametern und
unterschiedlichen GPU-Typen (Abschnitt~\ref{sec:methodischer-ansatz})
lassen sich Größeneffekte gegen die Streuung über Seeds abgrenzen, aber
keine Skalierungsgesetze schätzen.
Die Frage, ob der Mehraufwand größerer LLMs in einem angemessenen
Verhältnis zum Leistungsgewinn steht, wird deshalb für beide LLM-Regime
je Größenklasse beantwortet.

\section[Forschungslücke und Einordnung der eigenen Arbeit]{Forschungslücke und Einordnung der \mbox{eigenen} Arbeit}
\label{sec:research_gap}

Die vorangegangenen Abschnitte zeigen, dass direkte Vergleiche zwischen
Encodern und generativen LLMs für NER existieren und in ihren
Ergebnissen weitgehend übereinstimmen.
Die Forschungslücke liegt daher nicht im Fehlen solcher Vergleiche,
sondern in ihrer Reichweite.
Tabelle~\ref{tab:related-overview} stellt die behandelten Arbeiten nach
Modellen, Daten, Anpassungsverfahren, Art des Encoder-LLM-Vergleichs und
erfassten Ressourcen zusammen und ordnet die vorliegende Arbeit ein.

\begin{table}[p]
  \centering
  \scriptsize
  \setlength{\tabcolsep}{3pt}
  \renewcommand{\arraystretch}{1.15}
  \caption[Übersicht der behandelten Arbeiten]{Übersicht der behandelten Arbeiten.
  Vergleich Encoder--LLM: „ja“ = beide Modellklassen in derselben Studie
  trainiert bzw. ausgewertet; „berichtete Werte“ = Vergleich mit aus
  anderen Publikationen übernommenen Werten; „--“ = kein Vergleich.
  Ressourcen gemessen: in der Studie erhobene Trainings- oder
  Inferenzgrößen; „--“ = keine Messung. Full-FT: vollständiges
  Fine-Tuning aller Gewichte.}
  \label{tab:related-overview}
  \begin{tabular}{@{}%
      >{\raggedright\arraybackslash}p{2.4cm}%
      >{\raggedright\arraybackslash}p{3.0cm}%
      >{\raggedright\arraybackslash}p{2.9cm}%
      >{\raggedright\arraybackslash}p{2.0cm}%
      >{\raggedright\arraybackslash}p{2.0cm}%
      >{\raggedright\arraybackslash}p{2.2cm}@{}}
    \toprule
    \textbf{Arbeit} & \textbf{Modelle} & \textbf{Daten/Domäne} &
    \textbf{Anpassung} & \textbf{Vergleich Encoder--LLM} &
    \textbf{Ressourcen gemessen} \\
    \midrule
    \textcite{tedeschi2022multinerd} & mBERT + BiLSTM-CRF &
      MultiNERD (10 Spr., 15 Typen); Gold-Benchmarks & Fine-Tuning &
      -- & -- \\
    \textcite{tan2023damo} & XLM-R-large + Retrieval; ChatGPT &
      MultiCoNER II (12 Spr., 33 Typen) & Fine-Tuning; Prompting &
      ja (Validierung) & -- \\
    \textcite{keraghel2024recent} & BERT, XLM-R, DeBERTa-v3, GLiNER;
      GPT-4o & 10 Datensätze, allgemein und Domänen &
      Fine-Tuning; Prompting mit Beispielentitäten & ja & -- \\
    \textcite{zaratiana2024gliner} & DeBERTa-v3 (50--300 Mio.) &
      Pile-NER; CrossNER, MIT & Pre-Training; Zero-Shot &
      berichtete Werte & -- \\
    \textcite{bogdanov2024nuner} & RoBERTa-base; UniNER-7B; GPT-3.5/4 &
      C4-Annotationen; 4 Datensätze & Pre-Training; Few-Shot-FT &
      ja & Kosten der Annotation (unter 500~\$); keine Latenz- oder Speichermessung \\
    \textcite{semin2026spanLabeling} & Qwen3, Mistral, Llama, Gemma
      (8--70 Mrd.) & UNER-Benchmark u.\,a. & Prompting & -- &
      Ausgabetokens \\
    \textcite{wang2025gptner} & GPT-3 & CoNLL, OntoNotes, ACE, GENIA &
      Few-Shot-Prompting (Retrieval mit feinjustiertem Tagger) &
      berichtete Werte & -- \\
    \textcite{xie2023zeroshot} & GPT-3.5 & 7 Benchmarks (EN, ZH) &
      Zero-Shot-Prompting & berichtete Werte & -- \\
    \textcite{ashok2023promptner} & GPT-4, GPT-3.5, T5 &
      CoNLL, GENIA, FewNERD, CrossNER & Few-Shot-Prompting &
      berichtete Werte & -- \\
    \textcite{li2023codeie} & Codex, GPT-3 & CoNLL, ACE;
      RE-Datensätze & Few-Shot-Prompting & berichtete Werte (UIE) & -- \\
    \textcite{zhou2024universalner} & LLaMA 7B/13B &
      43 Datensätze, 9 Domänen (inkl. MultiNERD-Variante) &
      Destillation; Fine-Tuning & berichtete Werte & -- \\
    \textcite{sainz2024gollie} & Code-LLaMA 7/13/34B &
      IE-Datensätze; Zero-Shot u.\,a. CrossNER, MultiNERD & QLoRA & -- &
      Trainingszeit, CO$_2$ je Größe (keine Inferenz) \\
    \textcite{ding2024gner} & Flan-T5 (bis 11 Mrd.), LLaMA-7B &
      Pile-NER; Zero-Shot CrossNER, MIT & Fine-Tuning (Label je Wort) &
      berichtete Werte (GLiNER) & Durchsatz; GPU-Stunden des Trainings \\
    \textcite{zhan2026generativeNER} & acht LLMs (3--9 Mrd.);
      BERT-MRC u.\,a. & CoNLL, OntoNotes, ACE2005, GENIA & LoRA &
      teilweise (BERT-MRC selbst trainiert) & Trainingskosten (keine Inferenz) \\
    \textcite{obeidat2025llm} & BERT-large, BiomedBERT,
      DeBERTa-v3-large; Mistral-7B, Llama-3.1-8B &
      5 biomedizinische Datensätze & Full-FT; QLoRA & ja &
      Inferenzzeit, GPU-Anzahl \\
    \textcite{lu2024llmstruggle} & 7B-LLMs; GPT-3.5/4; BioClinicalBERT &
      RareDis-v1 (klinisch) & Prompting; LoRA & ja & -- \\
    \textcite{dukic2024lookingright} & Llama-2-7B, Mistral-7B; RoBERTa &
      CoNLL, Rest14, ACE05 & QLoRA mit Maskenentfernung & ja & -- \\
    \textcite{lu2024padellm} & Llama-2-7B, Baichuan2-7B &
      CoNLL, ACE, GENIA; chin. Datensätze & Fine-Tuning & -- &
      Inferenzlatenz \\
    \midrule
    Diese Arbeit & DeBERTa-v3-base/-large; Qwen3.5-0.8B/\allowbreak-4B/\allowbreak-27B &
      MultiNERD (EN, 15 Typen); CrossNER-Transfer &
      Full-FT; Zero-Shot; QLoRA & ja (gemeinsame Pipeline) &
      Trainingsdauer, Latenz, GPU-Speicher \\
    \bottomrule
  \end{tabular}
\end{table}

Die erste Lücke betrifft Datenbasis und Pipeline, die ersten beiden der
vier Aspekte aus Abschnitt~\ref{sec:problem_ziel}.
Die direkten Vergleiche verwenden allgemeine Benchmarks wie CoNLL-2003
mit vier und OntoNotes mit 18 Typen aus mehreren Genres wie
Nachrichten, Rundfunk, Gesprächen und Webtexten~\cite{weischedel2013ontonotes}
oder Fachkorpora wie die biomedizinischen Datensätze bei
\textcite{obeidat2025llm}.
MultiCoNER~II bei \textcite{tan2023damo} ist mit 33 Typen zwar
feingranular, besteht aber aus kurzen, kontextarmen Sätzen und stellt
den Encodern nur ein per Prompting genutztes ChatGPT gegenüber.
Publizierte MultiNERD-Werte beruhen auf übernommenen Baselines oder
eigener Aufbereitung: \textcite{zhou2024universalner} übernehmen die
BERT-Werte aus InstructUIE auf dessen Datensatzvariante,
\textcite{zaratiana2024gliner} werten GLiNER auf der
InstructUIE-Aufbereitung mit 10\,000 Testbeispielen aus, und bei
\textcite{sainz2024gollie} bleibt das Verhältnis der 32\,908
Testbeispiele zu den offiziellen Splits offen.
Eine gemeinsame Auswertung beider Modellfamilien auf den offiziellen
englischen MultiNERD-Splits mit einer in derselben Pipeline trainierten
Encoder-Referenz fehlt damit; zudem beruhen die Encoder-Referenzwerte
der direkten Vergleiche teilweise auf übernommenen Baselines.
Die vorliegende Arbeit setzt hier an, indem sie DeBERTa-v3-base und
-large sowie Qwen3.5 in drei Größen auf den offiziellen englischen
MultiNERD-Splits in einer gemeinsamen Pipeline trainiert und auswertet,
mit derselben \texttt{seqeval}-Implementierung und derselben
Silberreferenz, und dabei Trainingsdauer, Inferenzlatenz und
GPU-Spitzenspeicher der Inferenz nach einem einheitlichen Protokoll
erfasst (Abschnitt~\ref{sec:effizienzmetriken}).
Mit der Indexing-Ausgabe wird zudem ein Format verwendet, das den
Eingabetext nicht reproduziert und dessen Formtreue eigens gemessen wird.
Diese Lücke betrifft vor allem FF1, den Vergleich von
Erkennungsleistung und Formtreue in einer gemeinsamen Pipeline, sowie
FF3, die Ableitung von Empfehlungen für die Auswahl geeigneter
NER-Systeme unter Speicher- und Latenzbeschränkungen
(Abschnitt~\ref{sec:problem_ziel}).

Die zweite Lücke betrifft Ressourcenerfassung und Modellgröße, den
dritten und vierten Aspekt.
Leistung und Ressourcen werden in den direkten Vergleichen nicht
gemeinsam erfasst: Inferenzzeiten ohne Speicher, Trainingskosten ohne
Inferenz oder gar keine Messung.
Zudem betrachten die direkten Vergleiche LLMs einer einzigen
Größenklasse; Größenstaffelungen finden sich nur innerhalb von
LLM-Familien, bei GoLLIE mit Trainingsressourcen je Größe, aber ohne
gemeinsamen Encodervergleich mit
Inferenzmessung~\cite{zhou2024universalner,sainz2024gollie}, während
Encoder-Studien die Größe nur innerhalb ihrer Klasse
variieren~\cite{keraghel2024recent}.
Keine der Arbeiten in Tabelle~\ref{tab:related-overview} untersucht eine
Größenstaffelung auf beiden Seiten mit gemeinsamer Leistungs- und
Ressourcenerfassung.
Die vorliegende Arbeit vergleicht drei Größen einer LLM-Familie zwischen
0{,}8 und 27~Mrd. Parametern in zwei Regimen, Zero-Shot und QLoRA, mit
zwei Encoder-Größen derselben Architektur und erfasst für jede
Konfiguration Leistung, Formtreue und Ressourcen.
Sie adressiert damit FF2, die Veränderung von Erkennungsleistung und
Ressourcenbedarf mit der Modellgröße, und liefert die Grundlage für die
größenbezogenen Empfehlungen in FF3.

Zwei Einschränkungen sind offen zu benennen.
Die erste betrifft eine mögliche Datenkontamination.
MultiNERD besteht aus Wikipedia- und Wikinews-Artikeln, die mit hoher
Wahrscheinlichkeit Teil der Pre-Training-Korpora der verwendeten LLMs
sind; die Zusammensetzung der Qwen3.5-Trainingsdaten ist nicht
vollständig offengelegt, sodass sich dies weder belegen noch ausschließen
lässt.
\textcite{zhan2026generativeNER} weisen auf bekannte Lecks von
Benchmark-Daten in LLM-Trainingsdaten hin und ziehen deshalb mit
CleanCoNLL einen Datensatz heran, dessen bereinigte Form kaum im
Pre-Training enthalten sein dürfte; ihr Befund, dass die Ergebnisse nicht
von den Typbezeichnungen abhängen, spricht gegen eine Memorisierung von
Entität-Typ-Paaren, gilt aber für CoNLL-2003 und CleanCoNLL.
Für MultiNERD bleibt die Frage offen.
Zu bedenken ist dabei, dass die Referenz ein automatisch erzeugter Silver
Standard ist: Ein Modell, das den Wikipedia-Text kennt, kennt damit noch
nicht die hier verwendeten Labels, auch wenn der Datensatz selbst seit
2022 öffentlich verfügbar ist.
Im QLoRA-Regime lernen alle Modelle das Typinventar ohnehin aus den
Trainingsdaten; im Zero-Shot-Regime könnte Vorwissen über die Artikel die
Ergebnisse dagegen begünstigen.
Die Arbeit kann diese Frage nicht klären; sie wird bei der Interpretation
der Zero-Shot-Ergebnisse berücksichtigt.

Die zweite Einschränkung betrifft die Vergleichsbedingungen.
Anders als eine idealisierte Studie hält die vorliegende Arbeit nicht
alle Bedingungen konstant: Die Konfigurationen liefen auf drei
GPU-Typen, die Hyperparameter wurden je Modellgröße gewählt, und Prompt,
Decoding und numerische Präzision unterscheiden sich zwischen Zero-Shot-
und QLoRA-Läufen (Abschnitt~\ref{sec:validitaet}).
Konstant sind Datenbasis, Splits, Evaluationsprotokoll und
Bibliotheksversionen.
Die Arbeit vergleicht deshalb konkrete Konfigurationen und bestimmt keine
isolierten kausalen Effekte einzelner Faktoren; Effizienzvergleiche
werden vorrangig innerhalb gleicher Hardware- und Präzisionsbedingungen
interpretiert.
Ähnliche Einschränkungen finden sich auch in den direkten Vergleichen
der Literatur: \textcite{obeidat2025llm} betreiben die LLMs auf zwei und
die Encoder auf einer GPU, \textcite{zhan2026generativeNER} berichten
Trainingskosten, aber keine Inferenzbedingungen; hier werden die
Abweichungen explizit ausgewiesen.
Der CrossNER-Ergänzungsversuch schließlich erweitert den Vergleich um
einen reinen Transfertest der jeweils größten Modelle beider Familien auf
die Grobtypen PER, ORG und LOC; er ist, wie in
Abschnitt~\ref{sec:llm_ner_related} dargelegt, nicht mit den Few-Shot-
und Richtlinienansätzen der Literatur vergleichbar und dient der
Einordnung, nicht dem Benchmark-Vergleich.

Aus der Literatur folgen drei Erwartungen als Lesehilfe, die
Kapitel~\ref{chap:fazit} je Forschungsfrage mit den Befunden abgleicht;
ein Hypothesentest ist damit nicht beabsichtigt.
E1 (zu FF1): Ein mit LoRA oder QLoRA angepasstes LLM erreicht das
Niveau eines vollständig feinjustierten Encoders oder übertrifft es
leicht~\cite{zhan2026generativeNER,obeidat2025llm}.
E2 (zu FF2 und FF3): Die Inferenz generativer LLMs ist um ein Vielfaches
langsamer als die eines Encoders; die Literatur nennt Faktoren von etwa
40 bis 220~\cite[Tab.~IV]{obeidat2025llm}.
E3 (zu FF1): Ohne Anpassung der Gewichte, also im Zero-Shot-Regime,
bleiben per Prompting genutzte LLMs deutlich hinter feinjustierten
Modellen zurück~\cite{keraghel2024recent,tan2023damo}.

% chapters/04_methodik.tex
%
% Methodikkapitel. Angaben zu Implementierung, Konfiguration und
% Laufumgebung sind gegen Quellcode und gespeicherte Laufmetadaten
% abgeglichen; die Zuordnung von Läufen zu Commits, Revisionen und
% Artefakten steht in anhang.tex (Label app:reproduzierbarkeit).
% Benötigte Pakete in der Präambel: xcolor, tabularx, listings, booktabs.

\chapter{Methodik}
\label{chap:methodik}

Anhang~\ref{app:reproduzierbarkeit} ordnet jedem hier beschriebenen Lauf
Konfiguration, Codestand,
Revisionen und Ausgabeverzeichnis zu;
Anhang~\ref{app:durchfuehrung} ergänzt Durchführungsdetails.

\section{Methodischer Ansatz}
\label{sec:methodischer-ansatz}

Die Arbeit verfolgt einen empirisch-vergleichenden Ansatz: Statt neue
Architekturen oder Trainingsverfahren zu entwickeln, stellt sie
etablierte NER-Verfahren unter möglichst einheitlichen Bedingungen
gegenüber.
Tragfähige Vergleiche setzen voraus, dass Datensatz, Datenaufteilung und
Evaluationsprotokoll über alle Konfigurationen konstant bleiben.

\paragraph{Drei Vergleichsregime.}
\begin{enumerate}
    \item \textbf{Encoder-Fine-Tuning:} Vortrainierte Transformer-Encoder
    werden zusammen mit einem linearen Klassifikationskopf vollständig
    feinjustiert und liefern direkt IOB2-Tag-Sequenzen
    \cite{devlin2019bert,keraghel2024recent}.

    \item \textbf{LLM-Zero-Shot:} Veröffentlichte, bereits vor- und
    nachtrainierte Decoder-LLMs werden ohne aufgabenspezifisches Training
    und ohne Demonstrationsbeispiele im Prompt eingesetzt und geben die
    Entitäten als JSON aus.
    Das Regime beschreibt die Leistung dieser Checkpoints unter den
    gewählten Prompt- und Decoding-Bedingungen
    \cite{brown2020language,zhao2023survey}.

    \item \textbf{LLM-Fine-Tuning mit LoRA (QLoRA-Regime):} Dieselben LLMs werden über
    trainierbare Low-Rank-Adapter angepasst, während die ursprünglichen
    Gewichte eingefroren bleiben \cite{hu2021lora}.
    Wegen der Modellgröße wird das Regime als QLoRA mit
    4-Bit-quantisiertem Basismodell umgesetzt \cite{dettmers2023qlora}
    (Abschnitt~\ref{sec:lora-konfiguration}) und im Folgenden
    QLoRA-Regime genannt.
\end{enumerate}

Encoder werden ausschließlich feinjustiert, da Token-Klassifikatoren mit
neu initialisiertem Kopf ohne Training keine NER-Vorhersagen erzeugen und
somit kein Zero-Shot-Pendant besitzen; Generalisten wie GLiNER
(Abschnitt~\ref{sec:encoder_ner_related}) bleiben als eigene Modellklasse außerhalb des
Vergleichs.

\paragraph{Experimentmatrix.}
Aus den Regimen ergeben sich acht Konfigurationen
(Tabelle~\ref{tab:experiment-matrix}).
Die fünf trainierten Konfigurationen werden mit den Seeds 42, 123 und 456
jeweils unabhängig trainiert und evaluiert (15 Trainingsläufe); die drei
Zero-Shot-Konfigurationen wurden je einmal mit Seed~42 ausgewertet.
Die Hauptauswertung umfasst damit 18 Testevaluationen.

Der Hauptstudie gingen explorative Vorläufe der fünf trainierten
Konfigurationen mit Seed~42 voraus, die der Stabilisierung der Pipeline
dienten und nicht in die Auswertung eingehen
(Tabelle~\ref{tab:datennutzung}, Tabelle~\ref{tab:app-vorlaeufe}).
Die Zero-Shot-Läufe stammen aus einem eindeutig identifizierbaren Commit
ohne lokale Änderungen (Anhang~\ref{app:reproduzierbarkeit}) und sind Teil
der Hauptauswertung.

\begin{table}[ht]
    \centering
    \small
    \caption[Untersuchte Konfigurationen]{Untersuchte Konfigurationen.
    Größe: relative Größenklasse innerhalb der Modellfamilie
    (Abschnitt~\ref{sec:groessenklassen}).
    Parameter: Textmodell ohne Adapter, bei Qwen3.5 ohne den nicht
    geladenen Vision-Encoder (Abschnitt~\ref{sec:llm-auswahl}).
    Trainierbar: bei LoRA die Adapterparameter.
    NF4: 4-Bit-quantisiertes Basismodell; die Adapter werden mit
    BF16-Mixed-Precision trainiert (Abschnitt~\ref{sec:lora-konfiguration}).}
    \label{tab:experiment-matrix}
    \setlength{\tabcolsep}{4pt}
    \begin{tabular}{clllrrr}
        \toprule
        \textbf{Größe} & \textbf{Modell} & \textbf{Regime} &
        \textbf{Präzision} & \textbf{Parameter} & \textbf{Trainierbar} &
        \textbf{Läufe} \\
        \midrule
        S & DeBERTa-v3-base  & Fine-Tuning & FP32 & 184\,M   & 184\,M   & 3 \\
        M & DeBERTa-v3-large & Fine-Tuning & FP32 & 434\,M   & 434\,M   & 3 \\
        S & Qwen3.5-0.8B     & Zero-Shot   & BF16 & 0,75\,B  & --       & 1 \\
        M & Qwen3.5-4B       & Zero-Shot   & BF16 & 4,21\,B  & --       & 1 \\
        L & Qwen3.5-27B      & Zero-Shot   & BF16 & 26,90\,B & --       & 1 \\
        S & Qwen3.5-0.8B     & QLoRA       & NF4  & 0,75\,B  & 6,4\,M   & 3 \\
        M & Qwen3.5-4B       & QLoRA       & NF4  & 4,21\,B  & 21,2\,M  & 3 \\
        L & Qwen3.5-27B      & QLoRA       & NF4  & 26,90\,B & 159,4\,M & 3 \\
        \bottomrule
    \end{tabular}
\end{table}

\paragraph{Variablen, Konstanten und Störgrößen.}
Die Konfigurationen unterscheiden sich in \emph{Modellfamilie und
Architektur} (Encoder-only gegenüber Decoder-only), \emph{Parameteranzahl}
und \emph{Anpassungsregime}; diese Merkmale werden nicht unabhängig
voneinander variiert. Der Seed dient bei den trainierten
Konfigurationen als Replikationsfaktor.
Abhängige Variablen sind
(1)~strikte Precision, Recall und F$_1$ auf Entitätsebene,
(2)~Trainingsdauer, Inferenzlatenz und GPU-Spitzenspeicher der Inferenz,
(3)~die Zahl gesamter und trainierbarer Parameter sowie
(4)~Kennzahlen zur Formtreue der LLM-Ausgaben.

Konstant sind die Datenbasis (der englische Teil von MultiNERD mit
seinen offiziellen Splits \cite{tedeschi2022multinerd}), das
Evaluationsprotokoll mit derselben Metrikimplementierung für beide
Pipelines (Abschnitt~\ref{sec:performanzmetriken}) und die Versionen der
zentralen Bibliotheken (Abschnitt~\ref{sec:rechenumgebung}).
Alle Konfigurationen werten den vollständigen Testsplit aus.
Nicht konstant sind dagegen der GPU-Typ, die je Modellgröße gewählten
Hyperparameter sowie Prompt, Decoding und numerische Präzision der
Zero-Shot- gegenüber den QLoRA-Läufen.
Abschnitt~\ref{sec:validitaet} fasst diese Abweichungen zusammen.

\paragraph{Aufbau der Pipelines.}
Die Encoder-Pipeline behandelt NER als Token-Klassifikation: Ein
DeBERTa-v3-Encoder mit linearem Klassifikationskopf sagt für jedes Wort
ein IOB2-Tag vorher
\cite{he2021debertav3,wolf2020transformers,devlin2019bert}.
Die LLM-Pipeline behandelt NER als strukturierte Generierung: Das Modell
erhält einen Prompt mit nummerierter Tokenliste und erzeugt ein
JSON-Array von Entitäten, das eine regelbasierte, deterministische
Nachverarbeitung in eine IOB2-Sequenz überführt
\cite{dagdelen2024structuredIE,semin2026spanLabeling}.
Erst diese Rückkonvertierung erlaubt die Evaluation beider Pipelines auf
derselben Repräsentation.
Der Gesamt-F$_1$ bewertet damit die vollständige Vorhersagepipeline
einschließlich Ausgabeformat und Validierung.

\section{Auswahl der Modelle}
\label{sec:modellauswahl}

\subsection{Encoder-Modelle (DeBERTa-v3)}
\label{sec:encoder-auswahl}

Auf der Encoder-Seite werden DeBERTa-v3-base und DeBERTa-v3-large
eingesetzt.
DeBERTa-Varianten zählen auf mehreren NER-Benchmarks zu den
leistungsfähigsten Encodern und dienen häufig als starke Baselines
\cite{keraghel2024recent}.
Base- und Large-Variante erlauben einen Größenvergleich innerhalb
derselben Architektur; einschließlich des Klassifikationskopfes umfassen
sie 183\,855\,391 bzw. 434\,043\,935 Parameter.
Beide Modelle stehen laut ihren Modellkarten unter der MIT-Lizenz
\cite{he2021debertav3card,he2021debertav3cardlarge} und sind über die
\emph{Transformers}-Bibliothek verfügbar \cite{wolf2020transformers}.

\subsection{LLM-Modelle (Qwen3.5)}
\label{sec:llm-auswahl}

Auf der LLM-Seite werden drei Größen der Qwen3.5-Familie verwendet:
Qwen3.5-0.8B, Qwen3.5-4B und Qwen3.5-27B.
Qwen3.5 baut auf der Architektur- und Trainingstradition von Qwen2.5 und
Qwen3 auf \cite{qwen2024qwen25,qwen2025qwen3,qwen2026qwen35blog}; die
Modelle sind laut ihren Modellkarten als offene Gewichte unter der
Apache-2.0-Lizenz verfügbar
\cite{qwen2026card08b,qwen2026card4b,qwen2026card27b}.
Für die Familie sprachen neben der Lizenz die fein abgestufte Größenleiter mit
gemeinsamem Tokenizer und Chat-Template, sodass Größeneffekte nicht mit
Architekturwechseln vermischt werden.
Die drei Stufen wurden so gewählt, dass sie den in der Literatur
untersuchten Bereich von 3 bis 9~Mrd. Parametern
(Abschnitt~\ref{sec:research_gap}) nach unten und oben verlassen und das
Training mit einer einzelnen GPU je Lauf möglich bleibt.
Verwendet werden die veröffentlichten, laut Modellkarten vor- und
nachtrainierten Checkpoints
\cite{qwen2026card08b,qwen2026card4b,qwen2026card27b}; der Begriff
\emph{Basismodell} bezeichnet im Folgenden lediglich die eingefrorenen
Ausgangsgewichte, auf die beim LoRA-Training Adapter aufgesetzt werden.

Laut den Konfigurationsdateien der verwendeten Revisionen sind alle drei
Größen dichte Decoder mit hybrider Attention
(Tabelle~\ref{tab:app-qwen-architektur}): In jedem Block aus vier Schichten
nutzen drei Schichten lineare Attention in Form von Gated DeltaNet, die
vierte vollständige Softmax-Attention mit Grouped-Query-Attention und
Ausgabe-Gating \cite{qwen2026card08b,qwen2026card4b,qwen2026card27b}.

Von den multimodal angelegten Checkpoints laden beide LLM-Regime (Zero-Shot
und QLoRA) mit der Klasse \texttt{AutoModelForCausalLM} nur das Textmodell; Vision-Encoder
und Module für Multi-Token-Prediction werden verworfen
(Tabelle~\ref{tab:app-qwen-architektur} in Anhang~\ref{app:durchfuehrung}).
Alle Parameterangaben beziehen sich daher auf das Textmodell, dessen Größe
von den nominellen Bezeichnungen abweicht
(Tabelle~\ref{tab:experiment-matrix}); die Stufen unterscheiden sich um
den Faktor 5,6 bzw. 6,4.

\subsection{Größenklassen}
\label{sec:groessenklassen}

Zur Strukturierung der Auswertung werden die Modelle in die Klassen
\emph{Small} (S), \emph{Medium} (M) und \emph{Large} (L) eingeteilt
(Tabelle~\ref{tab:experiment-matrix}).
Die Klassen bezeichnen die relative Position innerhalb der jeweiligen
Familie, nicht absolute Parametergrenzen; L ist nur auf der LLM-Seite
besetzt, und die Einteilung ist eine Festlegung dieser Arbeit für die
Analyse in Abschnitt~\ref{sec:einfluss-parameteranzahl}.

\section{Auswahl und Aufbereitung der Datensätze}
\label{sec:datensaetze}

\subsection{Auswahlkriterien}
\label{sec:datenauswahl}

Für MultiNERD sprechen vier Kriterien.
Erstens umfasst das Inventar 15 Entitätstypen
(Tabelle~\ref{tab:datensplits}) und geht damit über Person, Ort und
Organisation hinaus, sodass Unterschiede zwischen den Modellklassen auch
bei seltenen oder semantisch komplexen Typen sichtbar werden können.
Zweitens kombiniert der Datensatz mit Wikipedia- und Wikinews-Artikeln
zwei Textgenres \cite{tedeschi2022multinerd}.
Drittens stehen offizielle \texttt{train}-, \texttt{validation}- und
\texttt{test}-Splits zur Verfügung \cite{tedeschi2022multinerd}, die
künftigen Arbeiten einen festen Bezugspunkt bieten.
Veröffentlichte MultiNERD-Werte beruhen dagegen auf eigener Aufbereitung
oder abweichendem Testumfang; Werte auf den offiziellen Splits sind
nicht auffindbar (Abschnitte~\ref{sec:k3-multinerd}
und~\ref{sec:research_gap}).
Viertens ist der englische Teil mit rund 164\,000 Sätzen groß genug für
vollständiges Fine-Tuning der Encoder ebenso wie für LoRA-Training der
LLMs.
Der Datensatz ist laut Datensatzkarte unter CC~BY-NC-SA~4.0
veröffentlicht und auf nichtkommerzielle Forschung beschränkt
\cite{babelscape2023multinerdcard}.

MultiNERD ist automatisch annotiert \cite{tedeschi2022multinerd}; auch
die Referenz des Testsplits ist somit eine Silberreferenz.
Die von den Autoren manuell geprüfte Stichprobe von 1\,000 zufällig
gezogenen englischen Sätzen wird nicht verwendet.
Sie ist in den veröffentlichten Daten nicht enthalten, und aus welchem
Split ihre Sätze stammen, ist nicht angegeben, sodass Überschneidungen
mit den Trainingsdaten nicht auszuschließen sind.
Typbezogene Aussagen wären auf ihr zudem sehr unsicher
\cite[Tab.~5]{tedeschi2022multinerd}.
Nach dieser Stichprobe unterscheidet sich die Annotationsqualität stark
zwischen den Typen (Abschnitt~\ref{sec:datasets_corpora}).
Die trainierten Konfigurationen nutzen dieselben annotierten
Trainingssätze, alle Systeme werden gegen dieselbe Referenz bewertet;
die Zero-Shot-Konfigurationen erhalten kein Training.
Verlustfunktion, Anpassungsumfang (vollständiges Fine-Tuning gegenüber
Adaptern) und Checkpoint-Auswahl unterscheiden sich jedoch, und die
gemeinsame Silberreferenz garantiert nicht, dass ihre Fehler für beide
Modellfamilien neutral wirken; absolute Werte messen zudem die
Übereinstimmung mit einer automatisch erzeugten Referenz.

\subsection{Datenvorverarbeitung}
\label{sec:datenvorverarbeitung}

Vorverarbeitet wird nur, was die technische Integration erfordert; eine
Textbereinigung oder Normalisierung auf Datensatzebene findet nicht
statt.
Der Datensatz wird über die \emph{Datasets}-Bibliothek in einer festen
Revision (Tabelle~\ref{tab:app-revisionen}) aus dem
Hugging-Face-Repository \texttt{Babelscape/\allowbreak multinerd} geladen,
über das Sprachfeld auf \texttt{en} gefiltert und auf vorhandene Splits
und Spalten sowie exakte Splitgrößen geprüft.
Eine zeilenweise Prüfung beim Laden unterbleibt; die nachträgliche
Prüfung der Rohdaten ergab weder leere Sätze noch Längeninkonsistenzen
oder Label-IDs außerhalb des Inventars.
Die Label-IDs werden über eine feste Liste (\texttt{O}, \texttt{B-PER},
\texttt{I-PER}, \dots, \texttt{I-VEHI}) in der Typreihenfolge von
Tabelle~\ref{tab:datensplits} auf Labelnamen abgebildet.

\paragraph{Encoder-Pipeline.}
Die wortweise vorliegenden Sätze werden mit dem Tokenizer von DeBERTa-v3
in Subwörter zerlegt.
Sequenzen werden auf 256 Subwort-Tokens einschließlich Sondertokens
gekürzt.
Beim Label-Alignment erhält das erste Subwort jedes Wortes dessen Label;
Folge-Subwörter und die Sondertokens \texttt{[CLS]} und \texttt{[SEP]}
erhalten den Wert $-100$ und werden in Verlust und Evaluation ignoriert
\cite{devlin2019bert,keraghel2024recent}.
Zur Evaluation werden die Vorhersagen der ersten Subwörter auf die
Wortebene übertragen, sodass Vorhersage und Referenz genau ein Label
pro Wort umfassen.
Wörter jenseits dieser Längengrenze entfallen; im Testsplit betrifft
dies keinen Satz, sodass alle 347\,025 Testwörter evaluiert werden.
Für Trainings- und Validierungssplit wurde die Zahl der gekürzten Sätze
nicht protokolliert.

\paragraph{LLM-Pipeline.}
Jedes Beispiel wird in eine Chat-Struktur aus Systemnachricht (Aufgabe
und Typinventar) und Nutzernachricht (nullbasiert nummerierte Tokenliste)
überführt, im Training ergänzt um eine Assistentennachricht mit der
Zielantwort (Abschnitt~\ref{sec:prompt-design}); die Nummerierung der
Eingabetokens nach \textcite[Tab.~9]{ramponi2025fallacy} erlaubt die
Rückabbildung generierter Spans auf die Tokenfolge
\cite{semin2026spanLabeling}.
Die Zielantwort wird deterministisch aus der IOB2-Sequenz abgeleitet: Jede
Entität wird in Satzreihenfolge mit den vier Feldern aus
Abschnitt~\ref{sec:prompt-design} in ein JSON-Array geschrieben; Sätze
ohne Entitäten erhalten das leere Array.
Für die Seed-42-Läufe aller trainierten Konfigurationen wurde bestätigt,
dass die verwendeten Referenzsequenzen für sämtliche Testsätze exakt der
Originalannotation entsprechen, in der LLM-Pipeline einschließlich der
Hin- und Rückkonvertierung.

\subsection{Datenaufteilung}
\label{sec:datenaufteilung}

Die offiziellen Splits werden unverändert und ohne Teilstichproben
übernommen; die Datenaufteilung ist damit für alle Konfigurationen
identisch und mit Arbeiten vergleichbar, die dieselben Splits verwenden
\cite{tedeschi2022multinerd}.
Die Typverteilung ist stark unausgewogen (Tabelle~\ref{tab:datensplits}):
PER, ORG und LOC stellen 74\,\% der Trainings- und 80\,\% der
Testentitäten, BIO, INST, MYTH und VEHI jeweils weniger als 40
Testentitäten.
Zudem verschiebt sich die Verteilung: Der Anteil von LOC steigt vom
Trainings- zum Testsplit von 29,1\,\% auf 46,9\,\%, der von PER sinkt
von 31,2\,\% auf 20,5\,\%.

\begin{table}[ht]
    \centering
    \small
    \caption[Umfang und Typverteilung der englischen MultiNERD-Splits]{Umfang und Typverteilung der englischen MultiNERD-Splits.
    Tokens bezeichnen Wort-Tokens vor der Subwort-Tokenisierung;
    Entitäten wurden als Anzahl der \texttt{B}-Tags gezählt.}
    \label{tab:datensplits}
    \begin{tabular}{lrrr}
        \toprule
                                & \textbf{train} & \textbf{validation} & \textbf{test} \\
        \midrule
        Sätze                   & 131\,280    & 16\,410  & 16\,454 \\
        Tokens                  & 2\,842\,119 & 379\,011 & 347\,025 \\
        Entitäten               & 201\,769    & 26\,041  & 25\,645 \\
        \midrule
        PER (Person)            & 62\,987 & 7\,507 & 5\,265 \\
        ORG (Organisation)      & 27\,641 & 2\,737 & 3\,309 \\
        LOC (Ort)               & 58\,665 & 7\,850 & 12\,024 \\
        ANIM (Tier)             & 12\,736 & 1\,134 & 1\,604 \\
        BIO (biologische Entität) & 140   & 14     & 8 \\
        CEL (Himmelskörper)     & 2\,685  & 94     & 41 \\
        DIS (Krankheit)         & 8\,702  & 1\,695 & 759 \\
        EVE (Ereignis)          & 2\,525  & 299    & 352 \\
        FOOD (Lebensmittel)     & 8\,279  & 2\,120 & 566 \\
        INST (Instrument)       & 379     & 34     & 12 \\
        MEDIA (Medien, Werke)   & 6\,081  & 919    & 458 \\
        MYTH (mythologische Entität) & 569 & 56    & 32 \\
        PLANT (Pflanze)         & 7\,436  & 1\,180 & 894 \\
        TIME (Zeitangabe)       & 2\,540  & 322    & 289 \\
        VEHI (Fahrzeug)         & 404     & 80     & 32 \\
        \bottomrule
    \end{tabular}
\end{table}

Der \texttt{train}-Split dient vollständig der Optimierung der
Modellparameter bzw. Adapter, der \texttt{validation}-Split
ausschließlich der Checkpoint-Auswahl innerhalb eines Laufs
(Abschnitte~\ref{sec:encoder-pipeline} und~\ref{sec:llm-pipeline}); ein
systematischer Hyperparameter-Suchlauf findet nicht statt.
Der \texttt{test}-Split wird in der Hauptstudie je Lauf genau einmal
vollständig inferiert; ein separater Evaluationsschritt berechnet die
Metriken aus den gespeicherten Vorhersagen.

Einschränkend wurde der Testsplit bereits in den explorativen Vorläufen
ausgewertet; Tabelle~\ref{tab:datennutzung} ordnet jeder Entscheidung
Datengrundlage und Testbezug zu, einschließlich der nachträglichen
Sensitivitätsanalysen zum Ausgabelimit.
Für die trainierten Konfigurationen blieben die Hyperparameter zwischen
Vorlauf und Hauptstudie bis auf die Epochenzahl des 27B-Modells
unverändert. Die Läufe der Hauptstudie
wurden anschließend mit drei Seeds neu trainiert. Diese Wiederholungen
erfassen Trainingsvariabilität, stellen aber die Unabhängigkeit eines
bereits ausgewerteten Testsplits nicht wieder her.
Bei den Zero-Shot-Konfigurationen beruhen dagegen Decoding und
Prompt-Erweiterung auf Beobachtungen eines Vorlaufs auf dem Testsplit
(Abschnitt~\ref{sec:llm-pipeline}).
Die Zero-Shot-Ergebnisse werden deshalb als explorative Ergebnisse der
nach Testbeobachtungen gewählten Konfigurationen berichtet. Solche
Auswahlentscheidungen können zu einer optimistischen Verzerrung führen
\cite{cawley2010selectionbias}; deren Umfang lässt sich mit den vorhandenen Daten
nicht beziffern.

\begin{table}[ht]
  \centering
  \small
  \caption[Datengrundlage der Konfigurationsentscheidungen]{Datengrundlage
  der Konfigurationsentscheidungen und ihr Bezug zum MultiNERD-Testsplit.
  Direkt: Entscheidung anhand von Beobachtungen auf dem Testsplit;
  indirekt: im Vorlauf verwendet und nach dessen Testauswertung
  beibehalten; post hoc: nach Kenntnis der Testergebnisse ergänzt.}
  \label{tab:datennutzung}
  \begin{tabularx}{\linewidth}{@{}>{\raggedright\arraybackslash}p{3.0cm}%
      >{\raggedright\arraybackslash}X%
      >{\raggedright\arraybackslash}p{2.9cm}%
      >{\raggedright\arraybackslash}p{3.3cm}@{}}
    \toprule
    \textbf{Entscheidung} & \textbf{Datengrundlage} & \textbf{Testbezug} &
    \textbf{Folge} \\
    \midrule
    Thinking-Modus deaktiviert & generative Validierung eines
      abgebrochenen 4B-Vorlaufs (Parse-Fehlerraten 71 bis 83\,\%) &
      keiner & für alle Läufe übernommen \\
    \addlinespace
    Hyperparameter der trainierten Konfigurationen & im Vorlauf
      verwendet; nach dessen Testauswertung unverändert beibehalten &
      indirekt & keine Anpassung an Testwerte \\
    \addlinespace
    Epochenzahl 27B (2 auf 3) & Angleichung an die übrigen
      QLoRA-Konfigurationen; Testwert des Vorlaufs bekannt & indirekt &
      Hauptläufe mit drei Epochen \\
    \addlinespace
    Zero-Shot-Prompt (Glossar) und Sampling & Vorlauf von Qwen3.5-0.8B
      auf dem Testsplit (Greedy: fast nur leere Arrays) & direkt,
      leistungsbezogen & als explorative Ergebnisse berichtet;
      optimistische Verzerrung möglich \\
    \addlinespace
    Inferenzpräzision BF16 (Zero-Shot) & Durchsatzmessung auf einer
      Teilmenge des Testsplits (standardmäßig 96 Sätze) & direkt, nicht
      leistungsbezogen & BF16 statt NF4 des QLoRA-Regimes \\
    \addlinespace
    Ausgabelimit 256 Tokens & vor den Läufen ohne Längenanalyse
      festgelegt & keiner & Hauptauswertung bei 256 Tokens \\
    \addlinespace
    Analysen mit 1\,024 Tokens & Nachinferenz, Kontrollstichprobe und
      CrossNER-Neulauf nach Kenntnis
      der Testbefunde & post hoc & als Sensitivitätsanalysen berichtet \\
    \addlinespace
    Checkpoint-Auswahl & Validierungssplit (LLM: erste 200 Sätze) &
      keiner & -- \\
    \bottomrule
  \end{tabularx}
\end{table}

\section{Versuchsaufbau}
\label{sec:versuchsaufbau}

\subsection{Rechenumgebung und Reproduzierbarkeit}
\label{sec:rechenumgebung}

Alle Läufe wurden auf dem Slurm-Cluster des DFKI mit je einer GPU pro Job
ausgeführt.
Abhängig von Modellgröße und Verfügbarkeit der Partitionen kamen drei
GPU-Typen zum Einsatz (Tabelle~\ref{tab:hardware}); den angeforderten
Hauptspeicher nennt Tabelle~\ref{tab:app-einstellungen}.
Die MultiNERD-Läufe verwenden, soweit protokolliert, einheitliche
Softwareversionen (Anhang~\ref{app:durchfuehrung}).

\begin{table}[ht]
    \centering
    \small
    \caption[Eingesetzte GPU-Typen]{Eingesetzte GPU-Typen laut Laufmetadaten.
    Training und Inferenz einer MultiNERD-Konfiguration liefen jeweils
    auf demselben GPU-Typ; die Sensitivitätsanalysen mit 1\,024 Tokens
    (Abschnitte~\ref{sec:nachinferenz} und~\ref{sec:crossner}) liefen
    auf der H100, für das 4B- und das 27B-Modell also auf einem anderen
    GPU-Typ als ihre Hauptläufe.}
    \label{tab:hardware}
    \setlength{\tabcolsep}{4pt}
    \begin{tabular}{lll}
        \toprule
        \textbf{Modell} & \textbf{Regime} & \textbf{GPU} \\
        \midrule
        DeBERTa-v3-base, -large & Fine-Tuning & NVIDIA H100 (80\,GB HBM3) \\
    \addlinespace
        Qwen3.5-0.8B            & QLoRA       & NVIDIA H100 (80\,GB HBM3) \\
    \addlinespace
        Qwen3.5-4B              & QLoRA       & NVIDIA A100-SXM4 (80\,GB) \\
    \addlinespace
        Qwen3.5-27B             & QLoRA       & NVIDIA H100 NVL (94\,GB) \\
    \addlinespace
        Qwen3.5-0.8B, -4B       & Zero-Shot   & NVIDIA A100-SXM4 (80\,GB) \\
    \addlinespace
        Qwen3.5-27B             & Zero-Shot   & NVIDIA H100 NVL (94\,GB) \\
        \midrule
        DeBERTa-v3-large, Qwen3.5-27B & CrossNER-Inferenz & NVIDIA H100 (80\,GB HBM3) \\
    \addlinespace
        Qwen3.5-27B             & CrossNER (1\,024) & NVIDIA H100 (80\,GB HBM3) \\
    \addlinespace
        Qwen3.5, alle Größen    & Nachinferenz (1\,024) & NVIDIA H100 (80\,GB HBM3) \\
    \addlinespace
        Qwen3.5, alle Größen    & Kontrollstichprobe (1\,024) & NVIDIA H100 (80\,GB HBM3) \\
        \bottomrule
    \end{tabular}
\end{table}

Die Trainingsläufe der Hauptstudie laden Datensatz und Modelle über feste
Revisions-Hashes, setzen ihren Seed für Python, NumPy und PyTorch und
protokollieren je Lauf Git-Commit, Paketversionen, GPU- und Slurm-Angaben (Anhang~\ref{app:durchfuehrung}); die Zero-Shot-Läufe
luden Modelle und Datensatz dagegen ohne explizite Revisionsangabe und mit
reduzierten Metadaten.
Eine bitgenaue Reproduzierbarkeit wird für keinen der Läufe beansprucht;
die Zero-Shot-Ausgaben entstehen zudem durch Sampling, dessen
Generatorzustand bei einer Wiederaufnahme nicht wiederhergestellt wird.
Lauf, Commit und Ausgabeverzeichnis (Tabelle~\ref{tab:app-laeufe}),
vollständige Commits und Revisionen (Tabelle~\ref{tab:app-revisionen})
sowie nachrangige Einstellungen (Tabelle~\ref{tab:app-einstellungen})
stehen im Anhang.

\subsection{Encoder-basierte NER-Pipeline}
\label{sec:encoder-pipeline}

Die DeBERTa-v3-Checkpoints werden mit einem neu initialisierten
Klassifikationskopf aus Dropout (Rate 0,1) und linearer Schicht über 31
Klassen geladen und end-to-end auf dem vollständigen Trainingssplit
feinjustiert \cite{devlin2019bert,wolf2020transformers}.
Der Verlust ist die mittlere Token-Kreuzentropie über alle Positionen mit
Label (Abschnitt~\ref{sec:datenvorverarbeitung}), ohne Klassengewichte
oder Label Smoothing.
Optimiert wird mit AdamW (Standard des Trainers) und einem linearen
Lernratenplan mit Warm-up; Lernrate, Warm-up und Batchaufteilung je Modell
nennt Tabelle~\ref{tab:encoder-hyperparams}, die effektive Batchgröße
beträgt bei beiden Modellen 16.
Trainiert wird vollständig in FP32, da bei DeBERTa-v3 auf H100-GPUs mit
Mixed Precision numerische Instabilitäten (NaN-Werte) auftraten.

Nach jeder Epoche wird der vollständige Validierungssplit ausgewertet
(strikter, mikrogemittelter F$_1$-Wert;
Abschnitt~\ref{sec:performanzmetriken}) und ein Checkpoint gespeichert;
für die Testevaluation wird der Checkpoint mit dem höchsten
Validierungs-F$_1$ geladen.
Die konfigurierte Early-Stopping-Patience entspricht der maximalen
Epochenzahl, sodass das Kriterium nicht auslösen kann; alle sechs Läufe
durchliefen die volle Epochenzahl, die je Lauf gewählte Epoche nennt
Tabelle~\ref{tab:app-seeds} (Abschnitt~\ref{sec:erg-gesamt}).
In der Testinferenz wird jeder Satz einzeln verarbeitet und die Klasse
per argmax bestimmt.

\begin{table}[ht]
  \centering
  \small
  \caption[Hyperparameter der Encoder-Pipeline]{Hyperparameter der Encoder-Pipeline.
  Weitere Einstellungen (Adam-Parameter, Gradient Clipping,
  Checkpoint-Intervalle) enthält Tabelle~\ref{tab:app-einstellungen}.}
  \label{tab:encoder-hyperparams}
  \begin{tabular}{lll}
    \toprule
    \textbf{Hyperparameter} & \textbf{DeBERTa-v3-base} &
    \textbf{DeBERTa-v3-large} \\
    \midrule
    Lernrate & $2 \cdot 10^{-5}$ & $8 \cdot 10^{-6}$ \\
    Lernratenplan & linear, Warm-up 10\,\% & linear, Warm-up 6\,\% \\
    Mikro-Batch $\times$ Akkumulation & $16 \times 1$ & $8 \times 2$ \\
    Max.\ Epochen (Early-Stopping-Patience) & 5 (5) & 4 (4) \\
    Optimierer & \multicolumn{2}{l}{AdamW (fused), Weight Decay 0,01} \\
    Max.\ Sequenzlänge & \multicolumn{2}{l}{256 Subwort-Tokens} \\
    Numerische Präzision & \multicolumn{2}{l}{FP32} \\
    Seeds & \multicolumn{2}{l}{42, 123, 456} \\
    \bottomrule
  \end{tabular}
\end{table}

\subsection{LLM-basierte NER-Pipeline}
\label{sec:llm-pipeline}

\paragraph{Gemeinsame Einstellungen.}
Die Prompts werden mit dem Chat-Template des Tokenizers gerendert und um
den Beginn der Assistentenantwort ergänzt.
Der Thinking-Modus von Qwen3.5 ist in allen generativen Schritten
deaktiviert (Datengrundlage der Entscheidung in
Tabelle~\ref{tab:datennutzung} und Anhang~\ref{app:durchfuehrung}); das
Template setzt dazu einen leeren Denkblock vor die
Antwort, sodass das Modell unmittelbar mit der JSON-Ausgabe beginnt.
Die Antwort ist auf 256 neue Tokens begrenzt.
Dieser Wert wurde vor den Läufen ohne vorherige Längenanalyse
festgelegt; nachträglich zeigt sich, dass er die Referenzantworten des
Testsplits nahezu vollständig abdeckt: 55 der
16\,454 Referenzantworten überschreiten unter dem Qwen3.5-Tokenizer
256 Tokens und keine 1\,024.
Training und Validierung, an denen sich das Limit vorab hätte prüfen
lassen, zeigen dasselbe Bild: 1\,024 Tokens decken
alle Validierungsantworten und bis auf eine alle
Trainingsantworten ab (Längenverteilungen in
Anhang~\ref{app:durchfuehrung}).
Das Limit der Nachinferenz von 1\,024 Tokens
(Abschnitt~\ref{sec:nachinferenz}) ist damit ohne Rückgriff auf den
Testsplit begründbar.
Die Eingabe wird nicht gekürzt, und jeder Testsatz wird einzeln
verarbeitet.

\paragraph{Zero-Shot-Regime.}
Die Checkpoints werden unverändert in BF16 verwendet; alle drei Größen
erhalten denselben Prompt (Glossar-Profil,
Abschnitt~\ref{sec:prompt-design}) und dieselben Decoding-Parameter.
Abweichend vom QLoRA-Regime wird mit Sampling decodiert (Temperatur 1,0,
top-$k = 20$, top-$p = 1{,}0$, min-$p = 0$, Seed~42).
Die Einstellung übernimmt top-$k$ und min-$p$ der Qwen-Modellkarte
\cite{qwen2026card27b}; top-$p$ und Presence-Penalty weichen von deren
Empfehlungen für den Non-Thinking-Modus ab, die Temperatur entspricht
nur dem Profil für Reasoning-Aufgaben (Anhang~\ref{app:durchfuehrung}).
Diese Einstellungen wurden nach einem explorativen Vorlauf mit
Qwen3.5-0.8B festgelegt, in dem Greedy-Decoding für nahezu jeden Satz ein
leeres Array erzeugte (strikter F$_1$-Wert 0,5 Punkte; 16\,052 von 16\,454
Antworten waren leere Arrays).
Zugleich wurde der Prompt um ausgeschriebene Typbezeichnungen und eine
Definition benannter Entitäten erweitert, da die Abkürzungen des
Inventars ohne Training nicht selbsterklärend sind.
Decoding-Einstellung und Prompt-Änderung wurden damit gleichzeitig
eingeführt; ihre Einzelbeiträge sind nicht trennbar, und die Übertragung
der Sampling-Einstellung auf das 4B- und das 27B-Modell wurde nicht
eigens geprüft (Abschnitt~\ref{sec:limitationen}).
Für diese beiden Modelle erfolgte somit keine Konfigurationswahl anhand
eigener Testergebnisse.
Die Ergebnisse sind je eine Stichprobe aus der Ausgabeverteilung der
unangepassten Checkpoints unter genau diesen Prompt- und
Decoding-Bedingungen; ein isolierter Effekt der Anpassung ist wegen der
weiteren Änderungen gegenüber dem QLoRA-Regime nicht bestimmbar.

\paragraph{QLoRA-Regime.}
Die Anpassung erfolgt mit \texttt{peft} \cite{mangrulkar2022peft} und dem
\emph{SFTTrainer} aus \texttt{trl} \cite{vonwerra2026trldocs} als
Single-Turn-SFT: Jedes Trainingsbeispiel besteht aus Systemnachricht,
Nutzernachricht und Zielantwort im Standardprofil
(Abschnitt~\ref{sec:prompt-design}).
Multi-Turn-Ansätze wie ChatIE werden zugunsten eines einfachen, klar
evaluierbaren Setups nicht verwendet \cite{wei2023chatie}.
Entsprechend dem Standardverhalten von TRL~1.8.0 für Konversationsdaten
wird der Verlust über die gesamte Tokenfolge berechnet und nicht auf die
Assistentenantwort beschränkt (Option \texttt{assistant\_only\_loss},
Standardwert \texttt{False} \cite{vonwerra2026trldocs}); ein Teil des
Lernsignals entfällt damit auf die Reproduktion des konstanten Prompts.
Trainingssequenzen werden auf 1\,024 Tokens begrenzt und bei Überlänge
am Ende gekürzt, was auch die Zielantwort abschneiden kann; dies betrifft
48 der 131\,280 Trainings- und 11 der 16\,410 Validierungsbeispiele.
Sequence Packing wird nicht verwendet.

Zur Checkpoint-Auswahl erzeugt das Modell nach jeder Epoche per Greedy-Decoding
Vorhersagen für die ersten 200 Sätze des Validierungssplits, die wie in
der Testinferenz validiert und mit dem strikten F$_1$-Wert bewertet werden;
der zusätzlich auf dem vollständigen Validierungssplit protokollierte
Verlust dient nur der Überwachung.
Verwendet wird der Adapter der Epoche mit dem höchsten Wert, bei
Gleichstand der frühere; in acht der neun Läufe war dies die dritte, in
einem Lauf (4B, Seed~42) die zweite Epoche.
Die 200 Sätze weichen in den Typanteilen um höchstens 3,1 Prozentpunkte
vom Split ab, enthalten aber keine BIO- und INST-Entitäten
(Auszählung im Repository, Anhang~\ref{app:reproduzierbarkeit}).
Die nach F$_1$ gewählte Epoche stimmt nur in einem der neun LLM-Läufe
(4B, Seed~42) und in keinem der sechs Encoder-Läufe mit dem Minimum des
Validierungsverlusts überein.
Bei den LLMs liegt dieses Minimum in Epoche 1 oder 2.
Da 12 der 15 Läufe (8 von 9 LLM-, 4 von 6 Encoder-Läufen) die letzte
Epoche wählten, ist ein Untertraining nicht auszuschließen
(Abschnitt~\ref{sec:limitationen}).
Die Beschränkung auf 200 Sätze hält die generative Validierung je Epoche
im Bereich von Minuten (27B-Modell: rund 18 Minuten gegenüber rund
25 Stunden für den vollständigen Split), macht die Auswahl aber von einem
kleinen, nicht zufällig gezogenen Ausschnitt abhängig.
Ob eine größere oder zufällige Stichprobe dieselben Checkpoints gewählt
hätte, wurde nicht geprüft (Abschnitt~\ref{sec:limitationen}).
Für die Testinferenz wird der Adapter ohne Verschmelzung (Merge) auf das
4-Bit-quantisierte Basismodell aufgesetzt; die Ausgaben entstehen per
Greedy-Decoding.

\subsection{Prompt-Design und Ausgabeformat}
\label{sec:prompt-design}

Jeder Prompt besteht aus einer Systemnachricht mit Aufgabe, Typinventar
und Formatvorgaben sowie einer Nutzernachricht mit der nummerierten
Tokenliste im Format \texttt{Index: Token}.
Erwartet wird ein JSON-Array, dessen Objekte genau vier Felder
enthalten: \texttt{start\_token} (Index des ersten Tokens, inklusiv),
\texttt{end\_token} (Index nach dem letzten Token, exklusiv),
\texttt{text} (der abgedeckte Text) und \texttt{type} (eine der 15
Typabkürzungen); ohne Entitäten ist das leere Array zurückzugeben.
Eine Entität, die nur das Token mit Index~4 umfasst, wird also durch
\texttt{start\_token}~=~4 und \texttt{end\_token}~=~5 beschrieben.
Die Einordnung des Formats in die \emph{Indexing}-Familie nach
\textcite[Abschn.~3.2]{semin2026spanLabeling} gibt
Abschnitt~\ref{sec:k3-ausgabeformate}.

Es werden zwei Prompt-Profile verwendet, die jeweils wortgleich für alle
drei Modellgrößen gelten.
Das Standardprofil (\texttt{task\_v1}) nennt das Typinventar nur in Form
der Abkürzungen und wird für Training, generative Validierung und
Testinferenz im QLoRA-Regime verwendet (Beispielprompt und Zielantwort:
Listing~\ref{lst:prompt-task} in Anhang~\ref{app:durchfuehrung}).
Das Glossar-Profil (\texttt{multinerd\_\allowbreak glossary\_\allowbreak v2})
wird ausschließlich im Zero-Shot-Regime eingesetzt; es ergänzt die
Systemnachricht um eine Definition benannter Entitäten, die
ausgeschriebenen Typbezeichnungen und die Aufforderung, die gesamte
Tokenliste zu prüfen (Listing~\ref{lst:prompt-glossar} im Anhang).
Keines der Profile enthält Demonstrationsbeispiele.

Auf \emph{Constrained Decoding} wird bewusst verzichtet.
Verfahren der geführten Generierung begrenzen die Ausgabe auf die von
einer formalen Grammatik zugelassenen Formen
\cite{geng2023grammar}, können aber die Modellleistung beeinträchtigen
\cite{tam2024speakfreely}.
Freie Generierung mit nachgelagerter Validierung lässt dem Modell
dagegen das volle Ausdrucksspektrum, macht die Formtreue messbar und
entspricht Einsatzszenarien, in denen nur die Textausgabe einer API
zugänglich ist \cite{dagdelen2024structuredIE}.

\subsection{LoRA-Konfiguration und Hyperparameter}
\label{sec:lora-konfiguration}

LoRA wird gewählt, weil LLMs ohne aufgabenspezifisches Fine-Tuning bei
NER auf Tokenebene nur eingeschränkt konkurrenzfähig sind
\cite{lu2024llmstruggle}, Adapter den Speicher- und Rechenbedarf
gegenüber vollständigem Fine-Tuning erheblich senken
\cite{hu2021lora,han2024parameter,obeidat2025llm} und in den Vergleichen
von \textcite[Anh.~D, Tab.~11]{zhan2026generativeNER} zudem 0,14 bis 0,48
F$_1$-Punkte vor vollständigem Fine-Tuning liegen
(Abschnitte~\ref{sec:lora} und~\ref{sec:peft_ner_related}).
Der Rang beträgt $r = 16$ für die 0.8B- und 4B-Modelle und $r = 32$ für
das 27B-Modell, die Skalierung jeweils $\alpha = 2r$
(Tabelle~\ref{tab:lora-hyperparams}).
Bias-Terme werden nicht trainiert, Varianten wie rsLoRA oder DoRA nicht
eingesetzt.
Der höhere Rang des 27B-Modells folgt der Heuristik, größeren Modellen
mehr Adapterkapazität zuzuweisen; alternative Ränge wurden nicht
untersucht.

Die Adapter werden über die Namen von sieben Projektionstypen platziert:
\texttt{q\_proj}, \texttt{k\_proj}, \texttt{v\_proj} und
\texttt{o\_proj} in der Attention sowie \texttt{gate\_proj},
\texttt{up\_proj} und \texttt{down\_proj} im Feed-Forward-Netz.
Die Einbeziehung der Feed-Forward-Schichten folgt Befunden, nach denen
Adapter in unterschiedlichen Linearschichten komplementäre Kapazität
bereitstellen \cite{dettmers2023qlora,han2024parameter,wang2024parameter}.
Wegen der hybriden Architektur erfasst diese namensbasierte Auswahl
jedoch nicht alle Linearschichten: Die Attention-Projektionen existieren
nur in den Schichten mit voller Attention, während die anders benannten
Projektionen der Gated-DeltaNet-Schichten ohne Adapter bleiben.
Adaptiert werden somit die Feed-Forward-Netze aller Schichten und die
Attention jeder vierten Schicht (Modulzahlen in
Tabelle~\ref{tab:lora-hyperparams}).

Die Basismodelle werden mit \texttt{bitsandbytes}
\cite{dettmers2022llmint8} im 4-Bit-Format NormalFloat (NF4) geladen;
gerechnet wird in BF16, trainiert mit BF16-Mixed-Precision und Gradient
Checkpointing.
Nicht quantisierte Komponenten wie Embeddings und
Normalisierungsschichten bleiben in höherer Präzision.
Die gespeicherten Adapter der 0.8B- und 4B-Läufe liegen in BF16 vor, die
der 27B-Läufe nach der Wiederaufnahme aus Checkpoints
(Abschnitt~\ref{sec:effizienzmetriken}) in FP32
(Anhang~\ref{app:durchfuehrung}).
In zwei Punkten weicht die Umsetzung vom QLoRA-Standardrezept
\cite{dettmers2023qlora} ab: Double Quantization war vorgesehen, wurde
aber wegen eines fehlerhaft benannten Konfigurationsparameters nicht
aktiviert.
Nach \textcite[Abschn.~3]{dettmers2023qlora} spart sie im Mittel rund
0,37~Bit je Parameter ohne Leistungseinbuße, beim 27B-Modell also
höchstens rund 1,2~GiB des gemessenen Inferenzspeichers.
Zudem wird statt
eines Paged Optimizers der Standardoptimierer des Trainers verwendet
(AdamW, fusionierte PyTorch-Implementierung).
Die Konfiguration ist für alle drei Größen identisch, sodass sich
Größeneffekte nicht mit Unterschieden im Präzisionsregime vermischen.
Die Ergebnisse beziehen sich auf diese tatsächlich ausgeführte Variante.
Der Einfluss von Double Quantization, anderer Adapterplatzierung oder
eines auf die Antwort beschränkten Verlusts wurde nicht isoliert gemessen.

\begin{table}[ht]
  \centering
  \small
  \caption[LoRA-/QLoRA-Hyperparameter]{LoRA-/QLoRA-Hyperparameter.
  Der Anteil trainierbarer Parameter bezieht sich auf die Gesamtzahl aus
  Basismodell und Adaptern (0,759, 4,227 bzw. 27,055~Mrd.).
  Weitere Einstellungen enthält Tabelle~\ref{tab:app-einstellungen}.}
  \label{tab:lora-hyperparams}
  \begin{tabular}{lrrr}
    \toprule
    \textbf{Hyperparameter} & \textbf{0.8B} & \textbf{4B} & \textbf{27B} \\
    \midrule
    Rang $r$ / Skalierung $\alpha$ & 16 / 32 & 16 / 32 & 32 / 64 \\
    Adaptierte Module (FFN/Attention) & 72/24 & 96/32 & 192/64 \\
    Parameter Basismodell & 0,752\,B & 4,206\,B & 26,896\,B \\
    Trainierbare Parameter (Adapter) & 6,39\,M & 21,23\,M & 159,38\,M \\
    Anteil trainierbarer Parameter & 0,84\,\% & 0,50\,\% & 0,59\,\% \\
    Lernrate & $3 \cdot 10^{-4}$ & $2 \cdot 10^{-4}$ & $1 \cdot 10^{-4}$ \\
    Warm-up-Anteil (Kosinus-Schedule) & 5\,\% & 5\,\% & 3\,\% \\
    Mikro-Batch $\times$ Akkumulation & $8 \times 2$ & $4 \times 4$ & $1 \times 16$ \\
    \midrule
    \multicolumn{4}{@{}l}{\emph{Für alle Größen identisch:}} \\
    LoRA-Dropout & \multicolumn{3}{l}{0,05} \\
    Quantisierung / Rechenpräzision & \multicolumn{3}{l}{NF4 (ohne Double Quantization) / BF16} \\
    Optimierer & \multicolumn{3}{l}{AdamW (fused), Weight Decay 0,01} \\
    Effektive Batchgröße & \multicolumn{3}{l}{16} \\
    Epochen / Optimierungsschritte & \multicolumn{3}{l}{3 / 24\,615} \\
    Max.\ Sequenzlänge (Training) & \multicolumn{3}{l}{1\,024 Tokens} \\
    Seeds & \multicolumn{3}{l}{42, 123, 456} \\
    \bottomrule
  \end{tabular}
\end{table}

\subsection{Nachverarbeitung und Validierung der Ausgaben}
\label{sec:postprocessing}

Die generierten Antworten werden in beiden LLM-Regimen mit derselben
deterministischen, regelbasierten Kette in IOB2-Sequenzen überführt:

\begin{enumerate}
    \item \textbf{Bereinigung und Parsing:} Die Antwort wird von
    umgebendem Leerraum und Denkblöcken befreit und als JSON-Array
    geparst, bei Bedarf über zwei Rückfallstrategien
    (Anhang~\ref{app:durchfuehrung}).
    Liefert keine Strategie ein JSON-Array, gilt die Antwort als nicht
    parsebar.

    \item \textbf{Validierung:} Jedes Array-Element wird einzeln geprüft
    und verworfen, sobald eine der vier Regeln verletzt ist:
    \begin{itemize}
        \item \emph{Schema:} Das Element muss ein Objekt mit allen vier
        Feldern sein, wobei zusätzliche Felder zulässig sind und ignoriert
        werden; \texttt{text} und \texttt{type} müssen nicht-leere
        Zeichenketten sein.
        \item \emph{Indizes:} Start- und Endindex $s$ und $e$ müssen
        ganze Zahlen mit $0 \leq s < e \leq n$ sein ($n$: Tokenanzahl).
        \item \emph{Typ:} Der Typ muss nach Entfernen umgebenden
        Leerraums und ohne Beachtung der Groß- und Kleinschreibung im
        Inventar enthalten sein.
        \item \emph{Exakter Text:} \texttt{text} muss nach Entfernen
        umgebenden Leerraums exakt dem mit Leerzeichen verbundenen Text
        der indizierten Tokens entsprechen.
    \end{itemize}
    Gültige Typen werden auf die kanonische Schreibweise normalisiert;
    unbekannte Typen werden nicht auf Inventartypen abgebildet, sondern
    verworfen.

    \item \textbf{Konfliktauflösung:} Gültige Spans werden nach Start-,
    dann Endindex und schließlich Reihenfolge in der Antwort sortiert und
    nur übernommen, wenn sie keine bereits belegten Tokens überdecken.
    Damit setzt sich der früher beginnende, bei gleichem Beginn der
    kürzere Span durch; von exakten Duplikaten bleibt das erste erhalten.

    \item \textbf{IOB2-Konvertierung:} Ausgehend von einer reinen
    \texttt{O}-Sequenz erhält das erste Token jedes Spans das Label
    \texttt{B-}\emph{Typ}, jedes weitere \texttt{I-}\emph{Typ}.
\end{enumerate}

Nicht parsebare Antworten werden wie leere Vorhersagen behandelt und
ergeben eine reine \texttt{O}-Sequenz; bei parsebaren Antworten bleiben
gültige Elemente erhalten.
Verworfene Elemente zählen nicht als Vorhersage: Entsprach ein
verworfenes Element einer Referenzentität, bleibt diese unerkannt und
mindert den Recall; andernfalls entfällt lediglich eine falsch positive
Vorhersage.
Die Kette toleriert umgebenden Text, Markdown-Formatierung, Leerraum,
Typschreibweise und zusätzliche Felder, repariert aber keine Inhalte;
insbesondere wird syntaktisch fehlerhaftes JSON nicht instand gesetzt.

\section{Evaluationskonzept}
\label{sec:evaluationskonzept}

\subsection{Metriken der Erkennungsleistung}
\label{sec:performanzmetriken}

Zentrale Metriken sind strikte Precision, Recall und F$_1$ auf
Entitätsebene nach IOB2-Schema \cite{manning2008introduction,sang2003conll,ramshaw1995text}.
Eine Entität gilt nur als korrekt erkannt, wenn Spangrenzen und Typ exakt
mit der Referenz übereinstimmen, was dem \emph{strict}-Modus der
SemEval-2013-Auswertung entspricht \cite{segura2013semeval}.
Berechnet werden die Werte mit \texttt{seqeval}
\cite{nakayama2018seqeval} bei explizit gesetztem striktem Modus und
IOB2-Schema, sodass formal unzulässige Tag-Folgen keine Entität bilden
(Abschnitt~\ref{sec:evaluation}).
Primäre Kennzahl ist der mikrogemittelte F$_1$-Wert
\cite{manning2008introduction,keraghel2024recent}; fehlen Vorhersagen
oder Referenzentitäten, wird der betreffende Quotient als 0 definiert.
Beide Pipelines werden auf Wortebene gegen identische Referenzsequenzen
ausgewertet
(Abschnitt~\ref{sec:datenvorverarbeitung}).

Für trainierte Konfigurationen werden Mittelwert und
Stichprobenstandardabweichung über die drei Seeds berichtet, für die
Zero-Shot-Konfigurationen der jeweilige Einzellauf. Die
Standardabweichung beschreibt die beobachtete Streuung der Trainingsläufe
und dient nicht als Schwelle für die Bewertung von Unterschieden.
Die Unsicherheit durch die Auswahl der Testsätze wird mit einem
gepaarten Satz-Bootstrap beziffert, der dem Bootstrap nach
\textcite{efron1993bootstrap} folgt und wie bei
\textcite{koehn2004bootstrap} dieselben Satzstichproben für alle
verglichenen Systeme zieht.
Aus den 16\,454 Testsätzen werden 2\,000 Stichproben mit Zurücklegen
gezogen; für jede Stichprobe werden Mikro- und Makro-F$_1$ aller Läufe
sowie die Differenzen zwischen Konfigurationen neu berechnet.
Pro Stichprobe werden die Laufwerte über die drei festen Seeds
gemittelt; berichtet werden 95-\%-Perzentilintervalle (2,5- und
97,5-Perzentil) dieser Mittelwerte und ihrer Differenzen.
Bei Zero-Shot geht der jeweilige Einzellauf ein und wird gegen jeden der
drei Seeds der Vergleichskonfiguration gepaart.
Ein Differenzintervall, das null ausschließt, spricht unter diesen
Annahmen für eine gerichtete Differenz. Enthält es null, wird kein
eindeutiger Unterschied festgestellt; Gleichwertigkeit ist damit nicht
nachgewiesen \cite{lakens2017equivalence}. Typbezogene Vergleiche ohne eigene Intervalle werden
deskriptiv behandelt.

Das Bootstrap ist auf die vorhandenen trainierten Modelle und
gespeicherten Vorhersagen bedingt. Es erfasst weder die Unsicherheit
durch neue Trainingsläufe noch die Sampling-Streuung einer erneuten
Zero-Shot-Inferenz. Die Sätze werden als unabhängige Einheiten behandelt;
Abhängigkeiten innerhalb von Quelldokumenten werden nicht modelliert.
Als Hauptvergleiche gelten die aus FF1 und FF2 abgeleiteten
Gegenüberstellungen der Konfigurationen in
Tabelle~\ref{tab:experiment-matrix} bei 256 Ausgabetokens
(DeBERTa-v3-large gegen Qwen3.5-27B im QLoRA-Regime sowie die
Größenschritte innerhalb der Familien); die Vergleiche der nachträglichen Sensitivitätsanalysen
(1\,024 Tokens, CrossNER-Domänen) sind explorativ.
Die Intervalle sind nicht für Mehrfachvergleiche korrigiert. Sie beseitigen auch
keine mögliche Verzerrung durch die frühere Testnutzung.
Für die Fehleranalyse werden Precision, Recall, F$_1$ und Support zusätzlich
je Entitätstyp berechnet; sie sind für die Trainingsläufe der Hauptstudie
gespeichert und für alle Konfigurationen aus den Testvorhersagen
reproduzierbar.
Wegen der unausgewogenen Typverteilung (Abschnitt~\ref{sec:datenaufteilung},
Tabelle~\ref{tab:datensplits}) sind typbezogene Werte seltener Klassen mit
erheblicher Unsicherheit behaftet.

Ergänzend wird der makrogemittelte F$_1$-Wert als ungewichtetes Mittel der
typbezogenen F$_1$-Werte über alle 15 Typen berichtet
(Abschnitt~\ref{sec:evaluation}); er reagiert entsprechend stark auf die
seltenen Typen.
Für den Vergleich der Fehlerprofile wird jede Referenzentität genau einer
Klasse zugeordnet: korrekt erkannt (Grenzen und Typ stimmen überein),
\emph{Typfehler} (gleiche Grenzen, anderer Typ), \emph{Grenzfehler}
(überlappender Span, gleicher Typ), \emph{Grenz- und Typfehler}
(überlappender Span, anderer Typ) oder \emph{Auslassung} (keine
überlappende Vorhersage).
Vorhergesagte Spans ohne Überlappung mit einer Referenzentität zählen
gesondert als \emph{Zusätze}.
Die Zuordnung folgt der Logik der weniger strengen Evaluationsmodi von
\textcite{segura2013semeval}, dient aber nur der Beschreibung; bewertet
wird in allen Fällen strikt.

\subsection{Effizienzmetriken}
\label{sec:effizienzmetriken}

Für praxisnahe NER-Systeme ist neben der Erkennungsleistung die Effizienz
von Training und Inferenz zentral \cite{han2024parameter,obeidat2025llm}.

\paragraph{Trainingsdauer.}
Berichtet wird für alle Konfigurationen die Wall-Clock-Dauer der
Slurm-Trainingsjobs laut Abrechnung (\texttt{sacct}); sie umfasst Laden
von Modell und Daten, Training, Validierung nach jeder Epoche (bei den
LLMs inklusive generativer Validierung), Speichern von Checkpoints und
Export.
Die vom Trainer zusätzlich protokollierte Dauer ohne Laden und Export
liegt bei den ununterbrochenen Läufen darunter: bei den Encodern um 30
bis 49~s, bei den 0.8B- und 4B-Läufen um 32 bis 141~s (Maximum beim
4B-Modell mit Seed~42).
Die drei 27B-Läufe überschritten das Zeitlimit der Partition von drei
Tagen und wurden vom Scheduler dreimal neu eingereiht
(Anhang~\ref{app:durchfuehrung}); ihre Trainingsdauer ist die Summe der
vier Zuteilungen einschließlich des erneuten Ladens und mit den übrigen
Läufen nur näherungsweise vergleichbar.
Als Plausibilitätsprüfung dient die Hochrechnung der Schrittrate der
letzten Zuteilung auf alle Schritte.
Die Trainingsdauer hängt zudem von der Auslastung des Knotens ab: Die
drei 27B-Läufe liefen gleichzeitig auf demselben Knoten und benötigten
rund 38~s je Optimierungsschritt, der explorative Vorlauf mit zwei
Epochen lief allein auf dem Knoten und benötigte 15{,}5~s je Schritt.
Sie beschreibt daher den tatsächlich angefallenen Aufwand dieser Jobs;
ein von Auslastung und Neustarts unabhängiger Trainingsbedarf der
Modellkonfiguration wird daraus nicht abgeleitet.

\paragraph{Inferenzlatenz und Durchsatz.}
Gemessen wird je Testsatz bei Batchgröße~1 die Dauer des Modellaufrufs
(Vorwärtsdurchlauf bzw. Generierung), mit GPU-Synchronisation vor und
nach dem Aufruf und nach einem nicht gewerteten Aufwärmdurchlauf.
Nicht enthalten sind Laden, Tokenisierung, Datentransfer, Detokenisierung,
Validierung, Metrikberechnung und Speichern.
Berichtet werden Mittelwert, Median und 95.~Perzentil in Millisekunden
(für die Zero-Shot-Läufe wurde der Median nicht protokolliert),
für die LLMs zusätzlich der Durchsatz in Testsätzen je Sekunde reiner
Generierungszeit.
Die Werte beschreiben die Latenz einzelner Anfragen, nicht den mit
Batching erreichbaren Durchsatz; bei den LLMs hängen sie zudem von der
Länge der erzeugten Antwort ab.
Die Gesamtdauer der Testinferenz stammt aus der Phasenhistorie der
Laufverwaltung, bei den Zero-Shot-Läufen ohne Phasenhistorie aus der
mittleren Latenz hochgerechnet.

\paragraph{GPU-Speicher.}
Erfasst wird der maximal von PyTorch allokierte GPU-Speicher während der
Testinferenz in MiB; die Ergebnistabellen rechnen durch Division durch
1\,024 in GiB um.
Der Zähler wird nach dem Laden des Modells zurückgesetzt, sodass der Wert
die Modellgewichte und den zusätzlichen Inferenzbedarf umfasst;
reservierter Speicher und der über \texttt{nvidia-smi} sichtbare
Verbrauch werden nicht erfasst.
Der GPU-Speicherbedarf während des Trainings wurde nicht protokolliert.

\paragraph{Parameterzahlen.}
Berichtet werden Gesamtzahl und Zahl der trainierbaren Parameter
(Tabelle~\ref{tab:experiment-matrix}); ihr Verhältnis quantifiziert die Reduktion der aktualisierten Parameter durch
das Adaptertraining \cite{hu2021lora,dettmers2023qlora}.
Für die QLoRA-Modelle nennt Tabelle~\ref{tab:lora-hyperparams} Basismodell und Adapter getrennt,
da eine einfache Summierung der Tensorelemente die in 4~Bit gepackten
Gewichte nur zur Hälfte erfasst.

\paragraph{Vergleichbarkeit.}
Da die Konfigurationen auf unterschiedlichen GPU-Typen liefen
(Tabelle~\ref{tab:hardware}), enthalten Unterschiede in der Laufzeit
auch Hardwareeffekte (Beispiel in Anhang~\ref{app:durchfuehrung}).
Zero-Shot- und QLoRA-Inferenz unterscheiden sich zudem in Präzision und
Decoding.
Effizienzvergleiche werden daher vorrangig innerhalb gleicher Hardware-
und Präzisionsbedingungen interpretiert.

\paragraph{Ressourcenschwellen für FF3.}
Speicher- und Latenzbudgets wurden vor den Läufen nicht festgelegt.
FF3 wird deshalb über Schwellen beantwortet, die nachträglich aus den
gemessenen Werten abgeleitet sind: Mittlere Latenz und Inferenzspeicher
jeder Konfiguration geben an, ab welcher Grenze sie zulässig ist, und für
jede Kombination solcher Grenzen wird die zulässige Konfiguration mit dem
höchsten Mikro-F$_1$ bestimmt (Abschnitt~\ref{sec:disk-effizienz}).
Die Schwellen beschreiben die gemessenen Bedingungen (Batchgröße~1,
jeweiliger GPU-Typ), keine Anforderungen einer konkreten Anwendung.

\subsection{Formtreue der Ausgaben}
\label{sec:robustheit-ausgaben}

Während Encoder per Konstruktion genau ein Label pro Wort aus dem festen
Inventar liefern, können LLMs syntaktisch oder inhaltlich unzulässige
Ausgaben erzeugen.
Diese erfasst Tabelle~\ref{tab:formtreue} über Kennzahlen der
Validierungskette (Abschnitt~\ref{sec:postprocessing}) in vier
Fehlerklassen: \emph{JSON-Validität}
\cite{dagdelen2024structuredIE}, \emph{Schema-Treue},
\emph{Span-Alignment} \cite{semin2026spanLabeling} und
\emph{Konsistenz}.
Typen außerhalb des Inventars~\cite{tedeschi2022multinerd} gelten dabei
als halluzinierte Kategorien.
Für Encoder entfallen diese Kennzahlen; der Unterschied ist Teil des
Vergleichs, da LLMs Ausgabeformat und Typinventar erst über Prompt und
Fine-Tuning erlernen müssen \cite{dagdelen2024structuredIE,wei2023chatie}.

\begin{table}[ht]
  \centering
  \small
  \caption[Kennzahlen zur Formtreue der LLM-Ausgaben]{Kennzahlen zur Formtreue der LLM-Ausgaben in der Reihenfolge
  der Prüfung.
  Die Parse-Fehlerrate bezieht sich auf die Zahl der Testantworten, alle
  übrigen Werte sind absolute Häufigkeiten über den Testsplit.}
  \label{tab:formtreue}
  \begin{tabularx}{\linewidth}{@{}>{\raggedright\arraybackslash}p{3.3cm}%
      >{\raggedright\arraybackslash}X%
      >{\raggedright\arraybackslash}p{2.4cm}@{}}
    \toprule
    \textbf{Kennzahl} & \textbf{Definition} & \textbf{Fehlerklasse} \\
    \midrule
    Parse-Fehlerrate & Anteil der Antworten, aus denen keine Strategie ein
      JSON-Array gewinnt & JSON-Validität \\
    \addlinespace
    Parse-Strategien & Antworten, die direkt, nach Entfernen eines
      Markdown-Codeblocks bzw.\ per Extraktion aus umgebendem Text geparst
      wurden & JSON-Validität \\
    \addlinespace
    Falsche Struktur & gültiges JSON, das kein Array ist (je
      Parse-Versuch) & Schema-Treue \\
    \addlinespace
    Nicht-Objekte & Array-Elemente, die kein JSON-Objekt sind &
      Schema-Treue \\
    \addlinespace
    Fehlende Felder & fehlendes oder leeres Pflichtfeld & Schema-Treue \\
    \addlinespace
    Unzulässige Indizes & nicht ganzzahlig oder Verletzung von
      $0 \leq s < e \leq n$ & Span-Alignment \\
    \addlinespace
    Unbekannte Typen & Typ außerhalb des 15-Typen-Inventars &
      Schema-Treue \\
    \addlinespace
    Text-Abweichungen & \texttt{text} passt nicht zu den indizierten
      Tokens & Span-Alignment \\
    \addlinespace
    Überlappungen & wegen Überlappung mit einem übernommenen Span
      verworfen & Konsistenz \\
    \bottomrule
  \end{tabularx}
\end{table}

Bei der Interpretation sind drei Eigenschaften zu beachten.
Erstens sind die Fehlerzähler absolute Häufigkeiten; da alle
Konfigurationen denselben Testsplit verarbeiten, sind sie zwischen
Konfigurationen vergleichbar, erlauben aber keine elementbezogenen Raten,
weil die Gesamtzahl erzeugter Elemente nicht protokolliert wird.
Zweitens wird jedes verworfene Element nur der ersten verletzten Regel
zugeordnet, sodass die Prüfreihenfolge die Verteilung auf die Kategorien
bestimmt.
Drittens gilt eine Antwort bereits als geparst, wenn sie ein JSON-Array
enthält, auch wenn danach alle Elemente verworfen werden; die
Parse-Fehlerrate ist daher eine Untergrenze für den Anteil formal
fehlerhafter Antworten.
Ob eine Antwort am Ausgabelimit abbrach, wird nicht protokolliert
(Abschnitt~\ref{sec:nachinferenz}).

\subsection{Analyse des Einflusses der Parameteranzahl}
\label{sec:einfluss-parameteranzahl}

Der Einfluss der Parameteranzahl auf Erkennungsleistung, Effizienz und
Formtreue
wird getrennt nach Regime über die Größenklassen
(Abschnitt~\ref{sec:groessenklassen}) deskriptiv betrachtet.
Mit zwei Encoder- und drei LLM-Größen, bei denen Hyperparameter und
Hardware mitvariieren, lassen sich Größenunterschiede zwar gegen die
Variabilität zwischen den Seeds abgrenzen, Skalierungsgesetze im Sinne
angepasster Potenzgesetze \cite{kaplan2020scaling,hoffmann2022training} aber
nicht bestimmen.
Leitfrage ist das Verhältnis von Leistungsgewinn und Ressourcenbedarf
größerer LLMs gegenüber kompakten Encodern
\cite{keraghel2024recent,obeidat2025llm}.

\subsection{Nachinferenz am Ausgabelimit}
\label{sec:nachinferenz}

Das Limit von 256 Ausgabetokens ist eine Einstellung der Pipeline, die bei
langen Sätzen mit vielen Entitäten greift
(Abschnitt~\ref{sec:llm-pipeline}).
Um seinen Einfluss auf MultiNERD zu messen statt hochzurechnen, wurden für
alle neun QLoRA-Läufe genau die Testsätze, deren gespeicherte Antwort nicht
parsebar war, mit demselben Adapter, demselben Prompt und Greedy-Decoding,
aber einem Limit von 1\,024 Tokens erneut generiert; alle übrigen
Antworten wurden unverändert übernommen.
Die Auswahlregel \enquote{nicht parsebar} ist ein Stellvertreter für
\enquote{abgeschnitten}: Jede nicht parsebare Antwort der QLoRA-Läufe
endet ohne schließende Klammer, und eine vollständige, aber aus anderem
Grund nicht parsebare Antwort würde unter Greedy-Decoding identisch
reproduziert; abgeschnittene, aber zufällig parsebare Antworten lassen
sich in den gespeicherten Daten dagegen nicht erkennen.
Bei identischen Logits und Greedy-Decoding bliebe eine Antwort, die vor
dem Limit mit dem Ende-Token abschließt, unverändert; da die
ursprünglichen Läufe den Stoppgrund nicht protokollierten, ist diese
Bedingung für die übernommenen Antworten nicht vollständig prüfbar.
Zwei Kontrollen prüfen die Annahme.
Für jeden nachgenerierten Satz wird festgehalten, ob die ersten
256 erzeugten Tokens mit der gespeicherten Antwort übereinstimmen.
Zusätzlich wurden nachträglich je Lauf 1\,000 zufällig gezogene
Testsätze mit parsebarer gespeicherter Antwort mit demselben Adapter,
demselben Prompt und Greedy-Decoding bei einem Limit von 1\,024 Tokens
neu generiert und zeichenweise mit der gespeicherten Antwort verglichen
(\emph{Kontrollstichprobe}, Anhang~\ref{app:durchfuehrung}).
Nachinferenz und Kontrollstichprobe liefen für alle drei Größen auf der
H100, für das 4B- und das 27B-Modell also auf einem anderen GPU-Typ als
die ursprüngliche Testinferenz (Tabelle~\ref{tab:hardware}); da auch die
Encoder auf der H100 liefen, liefert die Kontrollstichprobe zugleich
Latenzen auf einheitlicher Hardware.
Die zusammengeführten Vorhersagen werden mit derselben strikten Metrik über
alle 16\,454 Sätze ausgewertet und je Seed den Werten bei 256 Tokens
gegenübergestellt.
Die Hauptauswertung beruht weiterhin auf dem Limit von 256 Tokens; die
zusammengeführte Nachinferenz wird als Sensitivitätsanalyse und nicht als
vollständiger Neulauf aller Testsätze berichtet. Beide Auswertungen gehen
in die Schlussfolgerungen ein
(Abschnitt~\ref{sec:erg-formtreue}).

\section{CrossNER-Ergänzungsversuch}
\label{sec:crossner}

Ergänzend wird untersucht, wie sich die auf MultiNERD angepassten Modelle
auf domänenspezifische Texte übertragen lassen.
Grundlage ist CrossNER mit den Domänen Künstliche Intelligenz, Literatur,
Musik, Politik und Naturwissenschaften und jeweils domänenspezifischen
Entitätstypen \cite{liu2021crossner}.
Die Daten stehen laut der Lizenzdatei des Repositorys
\texttt{zliucr/CrossNER} (Commit in Tabelle~\ref{tab:app-revisionen})
unter der MIT-Lizenz.
Der Versuch ist ein reiner Transfertest: Auf CrossNER findet weder
Training noch Checkpoint-Auswahl statt.

\paragraph{Daten und Typabbildung.}
Verwendet werden ausschließlich die Testsplits in einer festgelegten, über SHA-256-Prüfsummen verifizierten Revision: 431 (KI), 416
(Literatur), 465 (Musik), 651 (Politik) und 543 Sätze
(Naturwissenschaften), insgesamt 2\,506.
Ausgewertet werden nur die gemeinsamen Grobtypen PER, ORG und LOC, denen
die CrossNER-Typen gemäß Tabelle~\ref{tab:crossner-mapping} zugeordnet
werden; alle übrigen Referenztypen gelten als \texttt{O}.
Vorhersagen anderer MultiNERD-Typen werden ebenfalls auf \texttt{O}
abgebildet, Spangrenzen bleiben unverändert.
Vorhersagen der drei Grobtypen für Spans, die in CrossNER einem nicht
abgebildeten Typ wie \texttt{misc} angehören, zählen als falsch positiv.

\begin{table}[ht]
  \centering
  \small
  \caption[Abbildung der CrossNER-Typen auf die gemeinsamen Grobtypen]{Abbildung der CrossNER-Typen auf die gemeinsamen Grobtypen;
  alle übrigen CrossNER-Typen (z.\,B. \texttt{misc}, \texttt{event},
  \texttt{product}, \texttt{song}) werden als \texttt{O} behandelt.}
  \label{tab:crossner-mapping}
  \begin{tabularx}{\linewidth}{@{}lX@{}}
    \toprule
    \textbf{Grobtyp} & \textbf{CrossNER-Typen} \\
    \midrule
    PER & \texttt{person}, \texttt{politician}, \texttt{scientist},
      \texttt{writer}, \texttt{musicalartist}, \texttt{researcher} \\
    \addlinespace
    ORG & \texttt{organisation}, \texttt{politicalparty},
      \texttt{university}, \texttt{band} \\
    \addlinespace
    LOC & \texttt{location}, \texttt{country} \\
    \bottomrule
  \end{tabularx}
\end{table}

\paragraph{Modelle und Inferenz.}
Untersucht werden die jeweils größten Modelle beider Familien,
DeBERTa-v3-large und Qwen3.5-27B (QLoRA), mit den Checkpoints aller drei
Seeds; sechs eingefrorene Modelle auf fünf Domänen ergeben
30 Evaluationen.
Die Inferenz entspricht der MultiNERD-Testinferenz: Der Encoder
verarbeitet jeden Satz ungekürzt (Sätze über 256 Subwort-Tokens würden
zum Abbruch führen), das LLM erhält unverändert das Standardprofil mit
allen 15 Typen und erzeugt seine Ausgaben per Greedy-Decoding, ohne Thinking, mit
4-Bit-Basismodell und höchstens 256 Ausgabetokens.
Alle CrossNER-Evaluationen liefen auf NVIDIA H100 (80\,GB HBM3), für
Qwen3.5-27B also auf einem anderen GPU-Typ als die MultiNERD-Testinferenz
(Tabelle~\ref{tab:hardware}); Latenzwerte sind daher nur eingeschränkt
zwischen beiden Datensätzen vergleichbar.

\paragraph{Auswertung.}
Je Domäne werden strikte Precision, Recall und F$_1$ berechnet und je Seed
zu einem gepoolten Mikrowert über alle Domänen sowie zu einem gleich
gewichteten Domänenmittel aggregiert; berichtet werden Mittelwert und
Stichprobenstandardabweichung über die Seeds.
Für die Differenz zwischen Qwen3.5-27B und DeBERTa-v3-large wird zusätzlich
ein gepaartes Satz-Bootstrap mit 2\,000 Replikaten berechnet, das je
Domäne geschichtet so viele Sätze mit Zurücklegen zieht, wie der
jeweilige Testsplit enthält; dieselbe Ziehung gilt für beide Systeme,
alle Seeds und beide Ausgabelimits.
Die 95-\%-Perzentilintervalle unterliegen denselben Einschränkungen wie
das MultiNERD-Bootstrap (Abschnitt~\ref{sec:performanzmetriken}); die
domänenweisen Vergleiche sind explorativ und nicht für Mehrfachvergleiche
korrigiert.
Bezugsgröße ist der F$_1$-Wert derselben Modelle auf dem MultiNERD-Testsplit,
ebenfalls auf PER, ORG und LOC beschränkt.
Die Differenz $\Delta$F$_1$ beschreibt den beobachteten Unterschied zwischen
beiden Datensätzen; sie umfasst neben dem Domänenwechsel auch
Unterschiede in Annotationsrichtlinien, Typdefinitionen und
Textmerkmalen und ist daher nicht als isolierter Effekt des
Domänenwechsels zu interpretieren.
Zusätzlich werden für das LLM die Parse-Fehlerrate sowie für beide
Modelle Latenz und GPU-Spitzenspeicher erfasst.

\paragraph{Überschneidungen.}
21 CrossNER-Testsätze (KI: 2, Musik: 4, Politik: 3,
Naturwissenschaften: 12) stimmen in ihrer vollständigen Tokenfolge exakt
mit Trainingssätzen von MultiNERD überein; Überschneidungen mit
Validierungs- oder Testsplit bestehen nicht, auf Beinahe-Duplikate wurde
nicht geprüft.
Die Hauptauswertung umfasst alle 2\,506 Sätze; eine Sensitivitätsanalyse
wiederholt die Aggregation ohne die 21 Sätze.

\paragraph{Ausgabelimit.}
Die Wirkung des Limits von 256 Ausgabetokens (Abschnitt~\ref{sec:nachinferenz})
auf Recall und F$_1$ des LLM prüft Abschnitt~\ref{sec:erg-crossner}.
Eine Vorstudie auf den jeweils ersten 25 Sätzen der Domänen Musik und
Politik, einer nicht repräsentativen Stichprobe, gab den Anlass für
einen vollständigen Neulauf mit 1\,024 Ausgabetokens
(Anhang~\ref{app:durchfuehrung}).
Der Neulauf des 27B-Modells umfasst alle drei Seeds und alle 2\,506
Sätze der fünf Domänen; Prompt, Parser, Typabbildung und Greedy-Decoding
blieben unverändert, nur das Limit wurde angehoben.
Die Auswertung stellt beide Limits auf denselben Sätzen je Seed gepaart
gegenüber und wiederholt die Sensitivitätsanalyse ohne die
21 überschneidenden Sätze (Abschnitt~\ref{sec:erg-crossner},
Tabelle~\ref{tab:limit}).

\section{Validität und bekannte Abweichungen}
\label{sec:validitaet}

Tabelle~\ref{tab:abweichungen} fasst die beschriebenen Abweichungen vom
idealisierten Versuchsdesign mit Verweis auf die jeweiligen Abschnitte
zusammen; ihre Folgen für die Aussagekraft der Ergebnisse diskutiert
Abschnitt~\ref{sec:limitationen}.

\begin{table}[ht]
  \centering
  \small
  \caption[Bekannte Abweichungen vom idealisierten Versuchsdesign]{Bekannte Abweichungen vom idealisierten Versuchsdesign.}
  \label{tab:abweichungen}
  \begin{tabularx}{\linewidth}{@{}>{\raggedright\arraybackslash}p{3.3cm}%
      >{\raggedright\arraybackslash}X%
      >{\raggedright\arraybackslash}p{2cm}@{}}
    \toprule
    \textbf{Aspekt} & \textbf{Abweichung} & \textbf{Abschnitt} \\
    \midrule
    Referenz & Alle MultiNERD-Splits automatisch annotiert
      (Silberreferenz) & \ref{sec:datenauswahl} \\
    \addlinespace
    Hardware & Drei GPU-Typen; CrossNER-Inferenz von Qwen3.5-27B auf
      anderem GPU-Typ als MultiNERD & \ref{sec:rechenumgebung},
      \ref{sec:crossner} \\
    \addlinespace
    Zero-Shot-Setup & Glossar-Prompt, Sampling und BF16 statt
      Standardprofil, Greedy-Decoding und NF4; ein nicht
      deterministischer Lauf je Größe & \ref{sec:llm-pipeline} \\
    \addlinespace
    Testnutzung & Testsplit bereits in Vorläufen ausgewertet;
      Zero-Shot-Prompt und Decoding nach Testbeobachtungen gewählt;
      Präzision anhand von Durchsatzmessungen & \ref{sec:datenaufteilung},
      Tab.~\ref{tab:datennutzung} \\
    \addlinespace
    Herkunft & Zero-Shot-Läufe ohne explizite Revisionsangabe und mit reduzierten
      Metadaten; keine bitgenaue Reproduzierbarkeit &
      \ref{sec:rechenumgebung} \\
    \addlinespace
    LoRA-Platzierung & Keine Adapter in den Gated-DeltaNet-Schichten &
      \ref{sec:lora-konfiguration} \\
    \addlinespace
    QLoRA-Umsetzung & Ohne Double Quantization und Paged Optimizer &
      \ref{sec:lora-konfiguration} \\
    \addlinespace
    SFT-Verlust & Verlust über die gesamte Sequenz einschließlich Prompt
      & \ref{sec:llm-pipeline} \\
    \addlinespace
    Adapter-Präzision & Gespeicherte 27B-Adapter nach der Wiederaufnahme in
      FP32 statt BF16 & \ref{sec:lora-konfiguration} \\
    \addlinespace
    Hyperparameter & Größenabhängig gewählt, kein Suchlauf (Encoder und
      LLM) & \ref{sec:encoder-pipeline}, \ref{sec:lora-konfiguration} \\
    \addlinespace
    Checkpoint-Auswahl & LLM-Adapter nach den ersten 200
      Validierungssätzen, Encoder nach dem vollständigen Split &
      \ref{sec:llm-pipeline}, \ref{sec:encoder-pipeline} \\
    \addlinespace
    Längengrenzen & 256 Subwort-Tokens (Encoder), 1\,024 Tokens (SFT),
      256 Ausgabetokens (LLM-Inferenz, a priori ohne Längenanalyse
      gesetzt; Sensitivitätsanalysen mit 1\,024 Tokens) & \ref{sec:datenvorverarbeitung},
      \ref{sec:llm-pipeline}, \ref{sec:crossner} \\
    \addlinespace
    Effizienzmessung & 27B-Training in vier Slurm-Zuteilungen, Dauer aus
      der Slurm-Abrechnung; drei 27B-Seeds gleichzeitig auf einem Knoten;
      kein Trainingsspeicher protokolliert &
      \ref{sec:effizienzmetriken} \\
    \bottomrule
  \end{tabularx}
\end{table}

% chapters/05_ergebnisse_und_diskussion.tex
%
% Erstentwurf vom 27.10.2026 (KI-gestützt erstellt, siehe Erklärung).
% Struktur, Tabellenskelette, Konventionen und Deutungsrahmen stehen;
% alle Messwerte und alle Aussagen, deren Richtung von den Messwerten
% abhängt, sind mit \PH{...} rot markiert und aus den Laufartefakten
% (Anhang~A.1: evaluation_metrics.yaml, inference_metrics.yaml,
% test_predictions.json) zu füllen. Die "% TODO"-Kommentare je Abschnitt
% nennen die dafür nötigen Auswertungen.
%
% Konventionen dieses Kapitels:
%  - Alle F1-, Precision- und Recall-Werte in Punkten (0--100).
%  - Trainierte Konfigurationen: Mittelwert +- Stichprobenstandardabweichung
%    über die Seeds 42, 123, 456; Zero-Shot: Einzellauf (Seed 42).
%  - Abschnitte 5.1 bis 5.6 berichten, 5.7 interpretiert.
%  - Bereits bekannte Werte (Support, Parameterzahlen, GPU-Typen,
%    gewählte Epochen, 8,7 ms Encoder-Latenz) sind eingetragen.
%
% Hinweis für Kapitel 4: Dort steht in 4.4.3 "strikter F1-Wert unter 0,01"
% (Skala 0--1). Für Konsistenz mit diesem Kapitel auf "unter 1" ändern.
\providecommand{\PH}[1]{\textcolor{red}{\textbf{[#1]}}}

\chapter{Ergebnisse und Diskussion}
\label{chap:ergebnisse}

Dieses Kapitel berichtet die Ergebnisse der in Kapitel~\ref{chap:methodik}
beschriebenen Experimente und diskutiert sie anschließend.
Abschnitt~\ref{sec:erg-leistung} stellt die Erkennungsleistung aller
Konfigurationen insgesamt und je Entitätstyp dar,
Abschnitt~\ref{sec:erg-formtreue} die Formtreue der LLM-Ausgaben; beide
zusammen beantworten FF1.
Abschnitt~\ref{sec:erg-groesse} betrachtet den Einfluss der Modellgröße,
Abschnitt~\ref{sec:erg-effizienz} den Ressourcenbedarf (FF2).
Abschnitt~\ref{sec:erg-fehlertypen} vergleicht die Fehlerprofile der
Modellklassen, Abschnitt~\ref{sec:erg-crossner} berichtet den
CrossNER-Ergänzungsversuch.
Abschnitt~\ref{sec:erg-diskussion} ordnet die Befunde in die verwandten
Arbeiten ein, leitet daraus die Empfehlungen zu FF3 ab und benennt,
welche Unterschiede angesichts der Streuung belastbar sind.
Die Abschnitte~\ref{sec:erg-leistung} bis~\ref{sec:erg-crossner}
beschränken sich auf die Darstellung der Messwerte; Interpretationen
folgen in Abschnitt~\ref{sec:erg-diskussion}.

Für alle Kennzahlen gelten die Festlegungen aus
Abschnitt~\ref{sec:evaluationskonzept}.
Precision, Recall und F$_1$ sind strikte Entity-Level-Werte nach dem
IOB2-Schema, berechnet mit \texttt{seqeval} auf dem vollständigen
englischen MultiNERD-Testsplit mit 16\,454 Sätzen und 25\,645
Referenzentitäten (Tabelle~\ref{tab:datensplits}); sie werden in Punkten
auf einer Skala von 0 bis 100 angegeben.
Für die fünf trainierten Konfigurationen werden Mittelwert und
Stichprobenstandardabweichung über die Seeds 42, 123 und 456 berichtet,
in den Tabellen als Mittelwert~$\pm$~Standardabweichung.
Die drei Zero-Shot-Konfigurationen beruhen auf je einem Lauf mit
Sampling (Abschnitt~\ref{sec:llm-pipeline}); ihre Werte sind daher
Einzelwerte ohne Streuungsangabe.
Ein Unterschied zwischen zwei trainierten Konfigurationen wird im
Folgenden nur dann als Unterschied bezeichnet, wenn er größer ist als
die Summe der beiden Standardabweichungen; kleinere Differenzen werden
als nicht unterscheidbar bezeichnet.
Dieses Kriterium ist in Abschnitt~\ref{sec:performanzmetriken}
festgelegt; es ersetzt keinen statistischen Test, macht die Lesart aber
nachvollziehbar.

\section{Erkennungsleistung auf MultiNERD}
\label{sec:erg-leistung}

Dieser Abschnitt berichtet die Erkennungsleistung zunächst über alle
Typen hinweg (Abschnitt~\ref{sec:erg-gesamt}) und anschließend je
Entitätstyp (Abschnitt~\ref{sec:erg-typen}).

\subsection{Gesamtergebnisse}
\label{sec:erg-gesamt}

% TODO: Werte aus evaluation_metrics.yaml der sechs DeBERTa-Läufe:
%  - Precision, Recall, Micro-F1 je Seed -> Mittelwert und Stichproben-SD
%  - Makro-F1 aus den typbezogenen F1-Werten (ungewichtetes Mittel über 15 Typen)
%  - Seed-Einzelwerte in tab:app-seeds (Anhang A.2)

Tabelle~\ref{tab:haupt} fasst die Erkennungsleistung aller acht
Konfigurationen zusammen; die Seed-Einzelwerte stehen in
Anhang~\ref{app:ergebnistabellen}.
Die beiden Encoder stehen in den ersten beiden Zeilen.

\begin{table}[htbp]
  \centering
  \small
  \caption{Erkennungsleistung aller Konfigurationen auf dem englischen
  MultiNERD-Testsplit (strikte Entity-Level-Werte, IOB2, in Punkten).
  Trainierte Konfigurationen: Mittelwert~$\pm$~Stichprobenstandardabweichung
  über drei Seeds; Zero-Shot: Einzellauf mit Sampling.
  Makro-F$_1$: ungewichtetes Mittel der typbezogenen F$_1$-Werte über
  15 Typen.
  Seed-Einzelwerte in Anhang~\ref{app:ergebnistabellen}.}
  \label{tab:haupt}
  \setlength{\tabcolsep}{4pt}
  \begin{tabular}{llcccc}
    \toprule
    \textbf{Modell} & \textbf{Regime} & \textbf{Precision} &
    \textbf{Recall} & \textbf{Micro-F$_1$} & \textbf{Makro-F$_1$} \\
    \midrule
    DeBERTa-v3-base  & Fine-Tuning & \PH{xx,x}\,$\pm$\,\PH{x,x} & \PH{xx,x}\,$\pm$\,\PH{x,x} & \PH{xx,x}\,$\pm$\,\PH{x,x} & \PH{xx,x} \\
    DeBERTa-v3-large & Fine-Tuning & \PH{xx,x}\,$\pm$\,\PH{x,x} & \PH{xx,x}\,$\pm$\,\PH{x,x} & \PH{xx,x}\,$\pm$\,\PH{x,x} & \PH{xx,x} \\
    \midrule
    Qwen3.5-0.8B & Zero-Shot & \PH{xx,x} & \PH{xx,x} & \PH{xx,x} & \PH{xx,x} \\
    Qwen3.5-4B   & Zero-Shot & \PH{xx,x} & \PH{xx,x} & \PH{xx,x} & \PH{xx,x} \\
    Qwen3.5-27B  & Zero-Shot & \PH{xx,x} & \PH{xx,x} & \PH{xx,x} & \PH{xx,x} \\
    \midrule
    Qwen3.5-0.8B & QLoRA & \PH{xx,x}\,$\pm$\,\PH{x,x} & \PH{xx,x}\,$\pm$\,\PH{x,x} & \PH{xx,x}\,$\pm$\,\PH{x,x} & \PH{xx,x} \\
    Qwen3.5-4B   & QLoRA & \PH{xx,x}\,$\pm$\,\PH{x,x} & \PH{xx,x}\,$\pm$\,\PH{x,x} & \PH{xx,x}\,$\pm$\,\PH{x,x} & \PH{xx,x} \\
    Qwen3.5-27B  & QLoRA & \PH{xx,x}\,$\pm$\,\PH{x,x} & \PH{xx,x}\,$\pm$\,\PH{x,x} & \PH{xx,x}\,$\pm$\,\PH{x,x} & \PH{xx,x} \\
    \bottomrule
  \end{tabular}
\end{table}

DeBERTa-v3-base erreicht einen Micro-F$_1$ von \PH{xx,x}, DeBERTa-v3-large
einen von \PH{xx,x}.
Die Differenz von \PH{x,x} Punkten ist \PH{größer / kleiner} als die Summe
der beiden Standardabweichungen (\PH{x,x} und \PH{x,x}); die größere
Variante ist damit \PH{unterscheidbar besser / nicht unterscheidbar}.
Precision und Recall liegen bei beiden Modellen \PH{nahe beieinander /
auseinander}: \PH{Precision xx,x gegenüber Recall xx,x (base) bzw.
xx,x gegenüber xx,x (large)}.
Der Makro-F$_1$ liegt mit \PH{xx,x} bzw. \PH{xx,x} um \PH{x} Punkte unter
dem Micro-F$_1$, was auf \PH{deutlich schwächere Werte bei den seltenen
oder verrauschten Typen} zurückgeht (Tabelle~\ref{tab:typen} in
Abschnitt~\ref{sec:erg-typen}).

Die Streuung über die Seeds ist mit Standardabweichungen von \PH{x,x}
(base) und \PH{x,x} (large) \PH{gering}.
Die Checkpoint-Auswahl fiel bei DeBERTa-v3-base auf die vierte
(Seed~42) bzw. fünfte Epoche (Seeds 123 und 456), bei DeBERTa-v3-large
auf die dritte bzw. vierte Epoche (Abschnitt~\ref{sec:encoder-pipeline});
die Validierungs-F$_1$-Werte der ausgewählten Checkpoints liegen bei
\PH{xx,x} bzw. \PH{xx,x} und damit \PH{über / nahe an} den Testwerten.
% TODO: Validierungs-F1 der gewählten Checkpoints aus results.yaml;
% Abweichung Validierung/Test kommentieren (Verteilungsverschiebung LOC
% 29,1 % -> 46,9 %, Abschnitt 4.3.3).


% TODO: Werte aus den drei Zero-Shot-Läufen (inference_metrics.yaml bzw.
% aus test_predictions.json nachgerechnet) und den neun QLoRA-Läufen
% (evaluation_metrics.yaml). Parse-Fehlerraten gehören in tab:formtreue-erg (Abschnitt 5.2).

Die Zeilen drei bis acht in Tabelle~\ref{tab:haupt} zeigen die
LLM-Konfigurationen.
Im Zero-Shot-Regime erreichen die drei Qwen3.5-Modelle Micro-F$_1$-Werte
von \PH{xx,x} (0.8B), \PH{xx,x} (4B) und \PH{xx,x} (27B).
\PH{Beschreibung des Profils: Precision gegenüber Recall je Größe, etwa
ob die Modelle viele Entitäten übersehen (niedriger Recall) oder
Nicht-Entitäten markieren (niedrige Precision).}
Die Werte sind Einzelwerte aus gesampelten Ausgaben und beruhen auf einem
Prompt und Decoding-Einstellungen, die nach Beobachtungen auf dem
Testsplit gewählt wurden (Abschnitt~\ref{sec:datenaufteilung}); sie sind
daher als obere Abschätzung dessen zu lesen, was die unveränderten
Checkpoints mit diesem Format leisten.

Nach der Anpassung mit QLoRA steigt der Micro-F$_1$ auf \PH{xx,x}
(0.8B), \PH{xx,x} (4B) und \PH{xx,x} (27B).
Der Zugewinn gegenüber dem Zero-Shot-Regime beträgt damit \PH{xx}, \PH{xx}
bzw. \PH{xx} Punkte und ist \PH{bei allen drei Größen deutlich größer als
die Streuung über die Seeds}.
Die Standardabweichungen über die Seeds liegen bei \PH{x,x}, \PH{x,x} und
\PH{x,x}.
Der Makro-F$_1$ der QLoRA-Modelle beträgt \PH{xx,x}, \PH{xx,x} und
\PH{xx,x}.
Ausgewählt wurde in acht der neun Läufe der Adapter der dritten Epoche,
beim 4B-Modell mit Seed~42 der der zweiten
(Abschnitt~\ref{sec:llm-pipeline}); die generative Validierung auf 200
Sätzen ergab für die gewählten Adapter F$_1$-Werte von \PH{xx,x},
\PH{xx,x} und \PH{xx,x}.
% TODO: Validierungs-F1 (200 Sätze, greedy) aus results.yaml; prüfen, ob
% der Verlauf über die Epochen (Epoche 1/2/3) bereits sättigt.


Der beste Encoder, \PH{DeBERTa-v3-large}, erreicht einen Micro-F$_1$ von
\PH{xx,x}; das beste LLM, \PH{Qwen3.5-27B im QLoRA-Regime}, einen von
\PH{xx,x}.
Der Abstand beträgt \PH{x,x} Punkte zugunsten \PH{des Encoders / des LLM}
und ist \PH{größer / kleiner} als die Summe der Standardabweichungen.
\PH{Einordnung der kleineren Konfigurationen: Liegt Qwen3.5-4B (QLoRA)
über oder unter DeBERTa-v3-base? Wo steht Qwen3.5-0.8B (QLoRA) relativ zu
den Encodern?}
Alle Zero-Shot-Konfigurationen liegen \PH{xx bis xx} Punkte unter dem
schwächsten trainierten Modell.

Die beiden Modellklassen unterscheiden sich im Profil aus Precision und
Recall: \PH{Encoder: Precision xx,x / Recall xx,x; bestes LLM: Precision
xx,x / Recall xx,x. Beschreiben, ob der Vorsprung einer Klasse vor allem
aus der Precision oder aus dem Recall stammt.}

\subsection{Ergebnisse nach Entitätstyp}
\label{sec:erg-typen}


% TODO: typbezogene P/R/F1 je Konfiguration (evaluation_metrics.yaml für
% die trainierten Läufe; Zero-Shot aus test_predictions.json nachrechnen).
% Mittelwerte über Seeds; vollständige Tabelle mit P/R/F1 in Anhang A.2.

\begin{table}[htbp]
  \centering
  \footnotesize
  \caption{Typbezogene F$_1$-Werte (strikt, in Punkten) ausgewählter
  Konfigurationen auf dem Testsplit; trainierte Konfigurationen als
  Mittelwert über drei Seeds, Zero-Shot als Einzellauf.
  Support: Zahl der Referenzentitäten im Testsplit
  (Tabelle~\ref{tab:datensplits}).
  Die Typen sind nach der Zuverlässigkeit ihrer Annotation gruppiert
  (Abschnitt~\ref{sec:datenauswahl}).
  Precision und Recall je Typ in Anhang~\ref{app:ergebnistabellen}.}
  \label{tab:typen}
  \setlength{\tabcolsep}{3pt}
  \begin{tabular}{lrcccccc}
    \toprule
    & & \multicolumn{2}{c}{\textbf{DeBERTa-v3}} &
      \multicolumn{3}{c}{\textbf{Qwen3.5 QLoRA}} &
      \textbf{Qwen3.5} \\
    \cmidrule(lr){3-4}\cmidrule(lr){5-7}\cmidrule(lr){8-8}
    \textbf{Typ} & \textbf{Support} & \textbf{base} & \textbf{large} &
      \textbf{0.8B} & \textbf{4B} & \textbf{27B} & \textbf{27B Zero-Shot} \\
    \midrule
    \multicolumn{8}{l}{\emph{Zuverlässig annotierte Typen}} \\
    PER   & 5\,265  & \PH{xx,x} & \PH{xx,x} & \PH{xx,x} & \PH{xx,x} & \PH{xx,x} & \PH{xx,x} \\
    ORG   & 3\,309  & \PH{xx,x} & \PH{xx,x} & \PH{xx,x} & \PH{xx,x} & \PH{xx,x} & \PH{xx,x} \\
    LOC   & 12\,024 & \PH{xx,x} & \PH{xx,x} & \PH{xx,x} & \PH{xx,x} & \PH{xx,x} & \PH{xx,x} \\
    CEL   & 41      & \PH{xx,x} & \PH{xx,x} & \PH{xx,x} & \PH{xx,x} & \PH{xx,x} & \PH{xx,x} \\
    DIS   & 759     & \PH{xx,x} & \PH{xx,x} & \PH{xx,x} & \PH{xx,x} & \PH{xx,x} & \PH{xx,x} \\
    EVE   & 352     & \PH{xx,x} & \PH{xx,x} & \PH{xx,x} & \PH{xx,x} & \PH{xx,x} & \PH{xx,x} \\
    MEDIA & 458     & \PH{xx,x} & \PH{xx,x} & \PH{xx,x} & \PH{xx,x} & \PH{xx,x} & \PH{xx,x} \\
    \midrule
    \multicolumn{8}{l}{\emph{Häufig verwechselte oder unvollständig verlinkte Typen}} \\
    ANIM  & 1\,604  & \PH{xx,x} & \PH{xx,x} & \PH{xx,x} & \PH{xx,x} & \PH{xx,x} & \PH{xx,x} \\
    BIO   & 8       & \PH{xx,x} & \PH{xx,x} & \PH{xx,x} & \PH{xx,x} & \PH{xx,x} & \PH{xx,x} \\
    FOOD  & 566     & \PH{xx,x} & \PH{xx,x} & \PH{xx,x} & \PH{xx,x} & \PH{xx,x} & \PH{xx,x} \\
    PLANT & 894     & \PH{xx,x} & \PH{xx,x} & \PH{xx,x} & \PH{xx,x} & \PH{xx,x} & \PH{xx,x} \\
    \midrule
    \multicolumn{8}{l}{\emph{Übrige Typen}} \\
    INST  & 12      & \PH{xx,x} & \PH{xx,x} & \PH{xx,x} & \PH{xx,x} & \PH{xx,x} & \PH{xx,x} \\
    MYTH  & 32      & \PH{xx,x} & \PH{xx,x} & \PH{xx,x} & \PH{xx,x} & \PH{xx,x} & \PH{xx,x} \\
    TIME  & 289     & \PH{xx,x} & \PH{xx,x} & \PH{xx,x} & \PH{xx,x} & \PH{xx,x} & \PH{xx,x} \\
    VEHI  & 32      & \PH{xx,x} & \PH{xx,x} & \PH{xx,x} & \PH{xx,x} & \PH{xx,x} & \PH{xx,x} \\
    \midrule
    Micro-F$_1$ (alle) & 25\,645 & \PH{xx,x} & \PH{xx,x} & \PH{xx,x} & \PH{xx,x} & \PH{xx,x} & \PH{xx,x} \\
    Makro-F$_1$ & -- & \PH{xx,x} & \PH{xx,x} & \PH{xx,x} & \PH{xx,x} & \PH{xx,x} & \PH{xx,x} \\
    \bottomrule
  \end{tabular}
\end{table}

Tabelle~\ref{tab:typen} schlüsselt die F$_1$-Werte nach Entitätstyp auf;
Abbildung~\ref{fig:typen} stellt die drei Kernkonfigurationen gegenüber.
Bei den drei häufigen Typen PER, ORG und LOC, die 80\,\% der
Testentitäten stellen, liegen alle trainierten Konfigurationen
\PH{über xx Punkten}; die Unterschiede zwischen den Modellklassen betragen
hier \PH{höchstens x,x Punkte}.
Die Typen ANIM, BIO, FOOD und PLANT, deren Silberannotation nach
\textcite{tedeschi2022multinerd} häufig verwechselt oder unvollständig
verlinkt ist, erreichen bei allen Modellen \PH{deutlich niedrigere} Werte
zwischen \PH{xx} und \PH{xx} Punkten.
\PH{Benennen, bei welchen Typen die Modellklassen am stärksten
auseinanderliegen, etwa ob das LLM bei DIS, EVE oder MEDIA vorn liegt
und der Encoder bei LOC oder ANIM.}
Für BIO (8 Referenzentitäten), INST (12), MYTH (32) und VEHI (32) sind
die Werte mit erheblicher Unsicherheit behaftet, da wenige Fehler den
F$_1$ um zweistellige Beträge verändern; sie werden nicht einzeln
interpretiert.
Im Zero-Shot-Regime \PH{beschreiben: welche Typen das 27B-Modell ohne
Training erkennt (vermutlich PER, LOC, ORG, DIS) und welche nahezu nicht
(vermutlich MEDIA, EVE, TIME mit abweichenden Typgrenzen)}.


\begin{figure}[htbp]
  \centering
  \fbox{\parbox[c][6cm][c]{0.9\linewidth}{\centering
    \PH{Abbildung: F$_1$ je Entitätstyp (y) für DeBERTa-v3-large,
    Qwen3.5-27B (QLoRA) und Qwen3.5-27B (Zero-Shot); Typen auf der
    x-Achse nach Support sortiert, Gruppen aus Tabelle~\ref{tab:typen}
    farblich abgesetzt; ein Punkt je Konfiguration und Typ}}}
  \caption{Typbezogene F$_1$-Werte der drei Kernkonfigurationen auf dem
  Testsplit (Werte aus Tabelle~\ref{tab:typen}).}
  \label{fig:typen}
\end{figure}

\section{Formtreue der LLM-Ausgaben}
\label{sec:erg-formtreue}

% TODO: Zähler aus inference_metrics.yaml je LLM-Lauf (Kennzahlen nach
% Tab. 4.8); QLoRA als Mittel über drei Seeds, gerundet; zusätzlich die
% Zahl nicht parsebarer Antworten ohne schließende Klammer (Indikator für
% Abbruch am 256-Token-Limit) aus test_predictions.json.

Tabelle~\ref{tab:formtreue-erg} berichtet die Kennzahlen zur Formtreue
aus Tabelle~\ref{tab:formtreue} für die sechs LLM-Konfigurationen.

\begin{table}[htbp]
  \centering
  \footnotesize
  \caption{Formtreue der LLM-Ausgaben auf dem Testsplit (16\,454
  Antworten).
  Parse-Fehlerrate in Prozent der Antworten; übrige Werte absolute
  Häufigkeiten (QLoRA: Mittel über drei Seeds, gerundet).
  Jedes verworfene Element zählt nur bei der ersten verletzten Regel
  (Abschnitt~\ref{sec:robustheit-ausgaben}).
  Leere Vorhersagen: Sätze ohne übernommenen Span.}
  \label{tab:formtreue-erg}
  \setlength{\tabcolsep}{4pt}
  \begin{tabular}{>{\raggedright\arraybackslash}p{4.4cm}rrrrrr}
    \toprule
    & \multicolumn{3}{c}{\textbf{Zero-Shot}} &
      \multicolumn{3}{c}{\textbf{QLoRA}} \\
    \cmidrule(lr){2-4}\cmidrule(lr){5-7}
    \textbf{Kennzahl} & \textbf{0.8B} & \textbf{4B} & \textbf{27B} &
      \textbf{0.8B} & \textbf{4B} & \textbf{27B} \\
    \midrule
    Parse-Fehlerrate [\%]          & \PH{x,x} & \PH{x,x} & \PH{x,x} & \PH{x,x} & \PH{x,x} & \PH{x,x} \\
    \quad davon ohne schließende Klammer & \PH{x} & \PH{x} & \PH{x} & \PH{x} & \PH{x} & \PH{x} \\
    Parse direkt / Codeblock / Extraktion & \PH{x/x/x} & \PH{x/x/x} & \PH{x/x/x} & \PH{x/x/x} & \PH{x/x/x} & \PH{x/x/x} \\
    Falsche Struktur (kein Array)  & \PH{x} & \PH{x} & \PH{x} & \PH{x} & \PH{x} & \PH{x} \\
    Nicht-Objekte                  & \PH{x} & \PH{x} & \PH{x} & \PH{x} & \PH{x} & \PH{x} \\
    Fehlende Felder                & \PH{x} & \PH{x} & \PH{x} & \PH{x} & \PH{x} & \PH{x} \\
    Unzulässige Indizes            & \PH{x} & \PH{x} & \PH{x} & \PH{x} & \PH{x} & \PH{x} \\
    Unbekannte Typen               & \PH{x} & \PH{x} & \PH{x} & \PH{x} & \PH{x} & \PH{x} \\
    Text-Abweichungen              & \PH{x} & \PH{x} & \PH{x} & \PH{x} & \PH{x} & \PH{x} \\
    Überlappungen                  & \PH{x} & \PH{x} & \PH{x} & \PH{x} & \PH{x} & \PH{x} \\
    \midrule
    Leere Vorhersagen & \PH{x} & \PH{x} & \PH{x} & \PH{x} & \PH{x} & \PH{x} \\
    \bottomrule
  \end{tabular}
\end{table}

Im QLoRA-Regime ist die Parse-Fehlerrate mit \PH{x,x}, \PH{x,x} und
\PH{x,x}\,\% \PH{gering}; \PH{nahezu alle / x von x} nicht parsebaren
Antworten enden ohne schließende Klammer und gehen damit mit hoher
Wahrscheinlichkeit auf das Limit von 256 Ausgabetokens zurück.
Unbekannte Typen und Text-Abweichungen treten nach dem Training
\PH{praktisch nicht mehr / in x Fällen} auf; die häufigste verbleibende
Fehlerklasse ist \PH{Überlappungen / unzulässige Indizes}.
Im Zero-Shot-Regime zeigt sich ein anderes Bild: \PH{Beschreiben je
Größe: Parse-Fehlerrate, Anteil der Extraktion aus umgebendem Text,
halluzinierte Typen (Beispiele nennen, etwa ausgeschriebene Typnamen
oder neue Kategorien), Text-Abweichungen als Indiz für fehlerhaft
gezählte Indizes, leere Arrays.}
Da die Zähler absolute Häufigkeiten ohne Bezug zur Zahl erzeugter
Elemente sind, lassen sie sich zwischen Konfigurationen vergleichen,
aber nicht als Fehlerraten je Element lesen
(Abschnitt~\ref{sec:robustheit-ausgaben}).

\section{Einfluss der Modellgröße auf die NER-Leistung}
\label{sec:erg-groesse}

% TODO: Differenzen zwischen den Größenklassen aus tab:haupt berechnen;
% Abbildung fig:groesse (F1 gegen Parameterzahl, log-x, Fehlerbalken = SD) aus
% den Mittelwerten erzeugen (pgfplots oder extern als PDF).

Tabelle~\ref{tab:groesse} und Abbildung~\ref{fig:groesse} stellen die
Erkennungsleistung über die Größenklassen aus
Abschnitt~\ref{sec:groessenklassen} getrennt nach Familie und Regime dar.

\begin{table}[htbp]
  \centering
  \small
  \caption{Micro-F$_1$ (in Punkten) nach Größenklasse und Regime.
  S/M/L bezeichnen die relative Größe innerhalb der Familie
  (Tabelle~\ref{tab:experiment-matrix}); $\Delta$ ist die Differenz zur
  nächstkleineren Klasse.
  Trainierte Konfigurationen als Mittelwert~$\pm$~Standardabweichung über
  drei Seeds.}
  \label{tab:groesse}
  \setlength{\tabcolsep}{3pt}
  \begin{tabular}{llccccc}
    \toprule
    \textbf{Familie} & \textbf{Regime} & \textbf{S} & \textbf{M} &
    \textbf{L} & $\Delta$(M$-$S) & $\Delta$(L$-$M) \\
    \midrule
    DeBERTa-v3 & Fine-Tuning & \PH{xx,x}\,$\pm$\,\PH{x,x} & \PH{xx,x}\,$\pm$\,\PH{x,x} & -- & \PH{+x,x} & -- \\
    Qwen3.5    & Zero-Shot   & \PH{xx,x} & \PH{xx,x} & \PH{xx,x} & \PH{+x,x} & \PH{+x,x} \\
    Qwen3.5    & QLoRA       & \PH{xx,x}\,$\pm$\,\PH{x,x} & \PH{xx,x}\,$\pm$\,\PH{x,x} & \PH{xx,x}\,$\pm$\,\PH{x,x} & \PH{+x,x} & \PH{+x,x} \\
    \bottomrule
  \end{tabular}
\end{table}

\begin{figure}[htbp]
  \centering
  \fbox{\parbox[c][6cm][c]{0.9\linewidth}{\centering
    \PH{Abbildung: Micro-F$_1$ (y) gegen Parameterzahl des Textmodells (x,
    logarithmisch) für DeBERTa-v3 (Fine-Tuning), Qwen3.5 (Zero-Shot) und
    Qwen3.5 (QLoRA); eine Linie je Familie und Regime, Fehlerbalken =
    Standardabweichung über Seeds}}}
  \caption{Erkennungsleistung in Abhängigkeit von der Parameterzahl,
  getrennt nach Modellfamilie und Regime.
  Parameterzahlen nach Tabelle~\ref{tab:experiment-matrix}
  (DeBERTa-v3: 184 und 434~Mio.; Qwen3.5: 0{,}75, 4{,}21 und 26{,}90~Mrd.).}
  \label{fig:groesse}
\end{figure}

Innerhalb der Encoder-Familie bringt der Schritt von 184 auf 434~Mio.
Parameter \PH{+x,x} Punkte.
Innerhalb der LLM-Familie wächst der Micro-F$_1$ im QLoRA-Regime von
\PH{xx,x} (0.8B) über \PH{xx,x} (4B) auf \PH{xx,x} (27B); der Zugewinn von
0.8B auf 4B beträgt \PH{x,x} Punkte, der von 4B auf 27B \PH{x,x} Punkte.
\PH{Einordnen: Sättigt der Zugewinn ab 4B? Ist der Schritt 4B zu 27B
überhaupt unterscheidbar (Summe der SD)?}
Im Zero-Shot-Regime fällt der Größeneffekt \PH{deutlich größer / ähnlich}
aus: \PH{xx,x}, \PH{xx,x} und \PH{xx,x} Punkte.
\PH{Falls das 0.8B-Modell im Zero-Shot nahezu versagt: das als Befund
festhalten (Formtreue in Abschnitt~\ref{sec:erg-formtreue} dazu heranziehen).}
Über die Familien hinweg erreicht \PH{das kleinste LLM im QLoRA-Regime
(0,75 Mrd.) / DeBERTa-v3-base (184 Mio.)} \PH{xx,x} Punkte und damit
\PH{mehr / weniger} als \PH{DeBERTa-v3-base / Qwen3.5-0.8B}, obwohl es
\PH{viermal so viele / ein Viertel so viele} Parameter besitzt.
Mit zwei Encoder- und drei LLM-Größen bleibt diese Betrachtung
deskriptiv (Abschnitt~\ref{sec:einfluss-parameteranzahl}); zudem
variieren Hyperparameter und GPU-Typ mit der Größe.

\section{Effizienzanalyse}
\label{sec:erg-effizienz}

% TODO: inference_metrics.yaml je Lauf: Latenz (Mittel, Median, P95),
% Durchsatz (LLMs), max. allokierter GPU-Speicher (MiB -> GB);
% results.yaml: Trainingsdauer (27B nur Untergrenze, Fußnote).
% Mittelwerte über Seeds; Abbildung fig:tradeoff (F1 gegen Latenz, F1
% gegen Speicher) aus tab:haupt und tab:effizienz erzeugen.

Tabelle~\ref{tab:effizienz} stellt den Ressourcenbedarf der acht
Konfigurationen nach dem Protokoll aus
Abschnitt~\ref{sec:effizienzmetriken} zusammen.
Die Latenz bezeichnet die mittlere Dauer des Modellaufrufs je Testsatz
bei Batchgröße~1; der GPU-Speicher ist der maximal von PyTorch
allokierte Speicher während der Testinferenz.
Da die Konfigurationen auf drei GPU-Typen liefen
(Tabelle~\ref{tab:hardware}), ist der GPU-Typ in der Tabelle
mitgeführt; Vergleiche werden vorrangig innerhalb gleicher Hardware
gezogen.

\begin{table}[htbp]
  \centering
  \footnotesize
  \caption{Ressourcenbedarf der Konfigurationen.
  Training: Wall-Clock-Dauer des Trainingsaufrufs in Stunden, Mittel über
  drei Seeds; Latenz: Dauer des Modellaufrufs je Testsatz bei
  Batchgröße~1 in Millisekunden (Mittelwert, Median, 95.~Perzentil);
  Durchsatz: Testsätze je Sekunde reiner Generierungszeit; Speicher:
  maximal allokierter GPU-Speicher während der Testinferenz in GB.
  Parameter: Textmodell gesamt / trainierbar
  (Tabelle~\ref{tab:experiment-matrix}).
  GPU: H100 = H100 80\,GB HBM3, A100 = A100-SXM4 80\,GB, NVL = H100 NVL
  94\,GB.}
  \label{tab:effizienz}
  \setlength{\tabcolsep}{3pt}
  \begin{tabular}{llrrrrrrr}
    \toprule
    \textbf{Modell} & \textbf{Regime} & \textbf{GPU} &
    \textbf{Training} & \multicolumn{3}{c}{\textbf{Latenz [ms]}} &
    \textbf{Durchsatz} & \textbf{Speicher} \\
    \cmidrule(lr){5-7}
    & & & \textbf{[h]} & \textbf{Mittel} & \textbf{Median} & \textbf{P95} &
    \textbf{[Sätze/s]} & \textbf{[GB]} \\
    \midrule
    DeBERTa-v3-base  & Fine-Tuning & H100 & \PH{x,x} & 8{,}7 & \PH{x,x} & \PH{x,x} & -- & \PH{x,x} \\
    DeBERTa-v3-large & Fine-Tuning & H100 & \PH{x,x} & \PH{xx,x} & \PH{x,x} & \PH{x,x} & -- & \PH{x,x} \\
    \midrule
    Qwen3.5-0.8B & Zero-Shot & A100 & -- & \PH{xxx} & \PH{xxx} & \PH{xxx} & \PH{x,x} & \PH{x,x} \\
    Qwen3.5-4B   & Zero-Shot & A100 & -- & \PH{xxx} & \PH{xxx} & \PH{xxx} & \PH{x,x} & \PH{xx,x} \\
    Qwen3.5-27B  & Zero-Shot & NVL  & -- & \PH{xxxx} & \PH{xxxx} & \PH{xxxx} & \PH{x,x} & \PH{xx,x} \\
    \midrule
    Qwen3.5-0.8B & QLoRA & H100 & \PH{x,x} & \PH{xxx} & \PH{xxx} & \PH{xxx} & \PH{x,x} & \PH{x,x} \\
    Qwen3.5-4B   & QLoRA & A100 & \PH{x,x} & \PH{xxx} & \PH{xxx} & \PH{xxx} & \PH{x,x} & \PH{x,x} \\
    Qwen3.5-27B  & QLoRA & NVL  & $\geq$\,\PH{xx,x}$^{a}$ & \PH{xxxx} & \PH{xxxx} & \PH{xxxx} & \PH{x,x} & \PH{xx,x} \\
    \bottomrule
    \multicolumn{9}{@{}l}{\scriptsize $^{a}$ Nur der Abschnitt nach der
      Wiederaufnahme protokolliert; Untergrenze
      (Abschnitt~\ref{sec:effizienzmetriken}).} \\
  \end{tabular}
\end{table}

\paragraph{Trainingsdauer.}
Das vollständige Fine-Tuning der Encoder dauerte \PH{x,x}~h (base,
fünf Epochen) bzw. \PH{x,x}~h (large, vier Epochen).
Das QLoRA-Training über drei Epochen benötigte \PH{x,x}~h für das
0.8B-Modell und \PH{x,x}~h für das 4B-Modell; für das 27B-Modell ist nur
eine Untergrenze von \PH{xx}~h bekannt.
\PH{Einordnen: Verhältnis der Trainingsdauern LLM zu Encoder; Hinweis,
dass die LLM-Läufe generative Validierung einschließen und das 4B-Modell
auf einer A100 lief.}

\paragraph{Inferenzlatenz.}
Die Encoder verarbeiten einen Testsatz in \PH{8{,}7} ms (base) bzw.
\PH{xx,x} ms (large).
Die LLMs benötigen im QLoRA-Regime \PH{xxx} ms (0.8B), \PH{xxx} ms (4B) und
\PH{xxxx} ms (27B) je Satz, also das \PH{xx}-, \PH{xx}- bzw. \PH{xxx}-Fache
des schnelleren Encoders.
\PH{Falls das 0.8B-Modell und DeBERTa-v3-base auf derselben H100 liefen:
diesen Vergleich als hardwaregleich hervorheben.}
Das 95.~Perzentil liegt bei den LLMs \PH{um das x-Fache} über dem Median,
was auf die Abhängigkeit der Generierungsdauer von der Zahl der
ausgegebenen Entitäten zurückgeht (Abschnitt~\ref{sec:effizienzmetriken}).
Die Zero-Shot-Läufe sind \PH{langsamer / schneller} als die QLoRA-Läufe
derselben Größe; sie unterscheiden sich in Präzision (BF16 statt NF4),
Decoding (Sampling statt Greedy) und beim 0.8B-Modell im GPU-Typ, sodass
dieser Unterschied nicht einem Faktor zugeordnet werden kann.

\paragraph{GPU-Speicher.}
Der Spitzenspeicher der Inferenz beträgt \PH{x,x} GB (DeBERTa-v3-base),
\PH{x,x} GB (DeBERTa-v3-large), \PH{x,x} GB (Qwen3.5-0.8B, QLoRA),
\PH{x,x} GB (Qwen3.5-4B, QLoRA) und \PH{xx,x} GB (Qwen3.5-27B, QLoRA).
Im Zero-Shot-Regime mit BF16-Gewichten liegt der Bedarf bei \PH{x,x},
\PH{xx,x} und \PH{xx,x} GB; der Unterschied zum QLoRA-Regime entspricht
im Wesentlichen der 4-Bit-Quantisierung des Basismodells.
\PH{Einordnen: Welche Konfigurationen passen auf eine GPU mit 24 GB bzw.
mit 16 GB? Das ist die Grundlage für FF3.}

\paragraph{Trade-off.}
Abbildung~\ref{fig:tradeoff} stellt den Micro-F$_1$ der Latenz und dem
GPU-Speicher gegenüber.
\PH{Beschreiben: Liegt eine Konfiguration auf der Pareto-Front (kein
anderes Modell ist zugleich besser und sparsamer)? Typischerweise: die
beiden Encoder und das beste LLM; einordnen, ob das 4B- oder das
27B-Modell dominiert wird.}

\begin{figure}[htbp]
  \centering
  \fbox{\parbox[c][6cm][c]{0.9\linewidth}{\centering
    \PH{Abbildung mit zwei Teilbildern: links Micro-F$_1$ (y) gegen
    mittlere Inferenzlatenz in ms (x, logarithmisch); rechts Micro-F$_1$
    gegen GPU-Spitzenspeicher in GB (x, logarithmisch). Ein Punkt je
    Konfiguration, Form nach Regime, Farbe nach Familie, GPU-Typ als
    Beschriftung; Pareto-Front als Linie}}}
  \caption{Erkennungsleistung gegenüber Inferenzlatenz (links) und
  GPU-Spitzenspeicher (rechts).
  Latenz und Speicher nach Tabelle~\ref{tab:effizienz}; Werte verschiedener
  GPU-Typen sind nur eingeschränkt vergleichbar.}
  \label{fig:tradeoff}
\end{figure}

\section{Fehleranalyse}
\label{sec:erg-fehlertypen}

% TODO: aus test_predictions.json je Lauf: Referenz- und vorhergesagte
% Spans aus den BIO-Sequenzen extrahieren (seqeval-Logik, strikt, IOB2)
% und jeden Referenzspan einer Klasse zuordnen:
%   korrekt (Grenzen und Typ gleich), Typfehler (Grenzen gleich, Typ
%   anders), Grenzfehler (überlappend, Typ gleich), Grenz- und Typfehler
%   (überlappend, Typ anders), Auslassung (keine Überlappung).
% Vorhergesagte Spans ohne Überlappung = Zusätze. Mittel über Seeds.
% Beispiele je Fehlerklasse für tab:app-beispiele (Anhang A.2) auswählen.

Um die Fehlerprofile der Modellklassen zu vergleichen, wird jede
Referenzentität nach dem in Abschnitt~\ref{sec:performanzmetriken}
festgelegten Schema genau einer Klasse zugeordnet: korrekt erkannt,
Typfehler, Grenzfehler, Grenz- und Typfehler oder Auslassung;
vorhergesagte Spans ohne Überlappung mit einer Referenzentität zählen als
Zusätze.
Tabelle~\ref{tab:fehlertypen} zeigt die Verteilung.

\begin{table}[htbp]
  \centering
  \footnotesize
  \caption{Fehlerprofile auf dem Testsplit: Verteilung der 25\,645
  Referenzentitäten auf Fehlerklassen (in Prozent) sowie Zahl der
  Zusätze (vorhergesagte Spans ohne Überlappung mit einer Referenz).
  Trainierte Konfigurationen als Mittel über drei Seeds.}
  \label{tab:fehlertypen}
  \setlength{\tabcolsep}{3pt}
  \begin{tabular}{llrrrrrr}
    \toprule
    \textbf{Modell} & \textbf{Regime} & \textbf{Korrekt} &
    \textbf{\shortstack[c]{Typ-\\fehler}} & \textbf{\shortstack[c]{Grenz-\\fehler}} &
    \textbf{\shortstack[c]{Grenz- und\\Typfehler}} &
    \textbf{\shortstack[c]{Aus-\\lassung}} & \textbf{Zusätze} \\
    \midrule
    DeBERTa-v3-base  & Fine-Tuning & \PH{xx,x} & \PH{x,x} & \PH{x,x} & \PH{x,x} & \PH{x,x} & \PH{xxx} \\
    DeBERTa-v3-large & Fine-Tuning & \PH{xx,x} & \PH{x,x} & \PH{x,x} & \PH{x,x} & \PH{x,x} & \PH{xxx} \\
    Qwen3.5-0.8B & QLoRA & \PH{xx,x} & \PH{x,x} & \PH{x,x} & \PH{x,x} & \PH{x,x} & \PH{xxx} \\
    Qwen3.5-4B   & QLoRA & \PH{xx,x} & \PH{x,x} & \PH{x,x} & \PH{x,x} & \PH{x,x} & \PH{xxx} \\
    Qwen3.5-27B  & QLoRA & \PH{xx,x} & \PH{x,x} & \PH{x,x} & \PH{x,x} & \PH{x,x} & \PH{xxx} \\
    Qwen3.5-27B  & Zero-Shot & \PH{xx,x} & \PH{x,x} & \PH{x,x} & \PH{x,x} & \PH{x,x} & \PH{xxxx} \\
    \bottomrule
  \end{tabular}
\end{table}

\PH{Beschreiben: Welche Fehlerklasse dominiert bei den Encodern, welche
bei den LLMs? Zhan et al. (2026) fanden auf CoNLL-2003, dass der Encoder
vor allem Entitäten übersieht und das LLM vor allem falsche Typen
vergibt; prüfen, ob sich das hier wiederholt. Grenzfehler bei
mehrteiligen Namen (ORG, MEDIA) und Typfehler zwischen ANIM/PLANT/FOOD
sowie LOC/ORG gesondert nennen.}
Beispiele typischer Fehler jeder Klasse zeigt
Tabelle~\ref{tab:app-beispiele} im Anhang.

\section{Transfer auf CrossNER}
\label{sec:erg-crossner}

% TODO: results/crossner/<lauf>/<domäne>/: P/R/F1 je Domäne und Seed für
% DeBERTa-v3-large und Qwen3.5-27B (QLoRA); gepoolter Mikrowert und
% Domänenmittel je Seed -> Mittel +- SD; MultiNERD-Referenz (nur
% PER/ORG/LOC) aus den MultiNERD-Testvorhersagen derselben Checkpoints;
% Sensitivität ohne die 21 überschneidenden Sätze; Parse-Fehlerrate,
% Latenz, Speicher auf CrossNER (beide auf H100 80 GB).

Tabelle~\ref{tab:crossner-erg} berichtet den Ergänzungsversuch aus
Abschnitt~\ref{sec:crossner}: die jeweils größten Modelle beider
Familien, auf MultiNERD trainiert, ausgewertet auf den Testsplits der
fünf CrossNER-Domänen und beschränkt auf die Grobtypen PER, ORG und LOC.

\begin{table}[htbp]
  \centering
  \footnotesize
  \caption{CrossNER-Transfertest (nur PER, ORG, LOC; strikte Werte in
  Punkten; Mittelwert~$\pm$~Standardabweichung über die drei
  MultiNERD-Checkpoints je Modell).
  Gepoolt: Mikrowert über alle 2\,506 Sätze; Domänenmittel: ungewichtetes
  Mittel der fünf Domänen-F$_1$.
  MultiNERD-Referenz: dieselben Checkpoints auf dem MultiNERD-Testsplit,
  ebenfalls auf PER, ORG und LOC beschränkt; $\Delta$F$_1$ = gepoolter
  CrossNER-Wert minus Referenz.
  Sensitivität: gepoolter Wert ohne die 21 Sätze, die mit
  MultiNERD-Trainingssätzen übereinstimmen.}
  \label{tab:crossner-erg}
  \setlength{\tabcolsep}{3pt}
  \begin{tabular}{>{\raggedright\arraybackslash}p{3.2cm}rcccccc}
    \toprule
    & & \multicolumn{3}{c}{\textbf{DeBERTa-v3-large}} &
      \multicolumn{3}{c}{\textbf{Qwen3.5-27B (QLoRA)}} \\
    \cmidrule(lr){3-5}\cmidrule(lr){6-8}
    \textbf{Domäne} & \textbf{Sätze} & \textbf{P} & \textbf{R} & \textbf{F$_1$} &
      \textbf{P} & \textbf{R} & \textbf{F$_1$} \\
    \midrule
    Künstliche Intelligenz & 431 & \PH{xx,x} & \PH{xx,x} & \PH{xx,x}\,$\pm$\,\PH{x,x} & \PH{xx,x} & \PH{xx,x} & \PH{xx,x}\,$\pm$\,\PH{x,x} \\
    Literatur              & 416 & \PH{xx,x} & \PH{xx,x} & \PH{xx,x}\,$\pm$\,\PH{x,x} & \PH{xx,x} & \PH{xx,x} & \PH{xx,x}\,$\pm$\,\PH{x,x} \\
    Musik                  & 465 & \PH{xx,x} & \PH{xx,x} & \PH{xx,x}\,$\pm$\,\PH{x,x} & \PH{xx,x} & \PH{xx,x} & \PH{xx,x}\,$\pm$\,\PH{x,x} \\
    Politik                & 651 & \PH{xx,x} & \PH{xx,x} & \PH{xx,x}\,$\pm$\,\PH{x,x} & \PH{xx,x} & \PH{xx,x} & \PH{xx,x}\,$\pm$\,\PH{x,x} \\
    Naturwissenschaften    & 543 & \PH{xx,x} & \PH{xx,x} & \PH{xx,x}\,$\pm$\,\PH{x,x} & \PH{xx,x} & \PH{xx,x} & \PH{xx,x}\,$\pm$\,\PH{x,x} \\
    \midrule
    Gepoolt                & 2\,506 & \PH{xx,x} & \PH{xx,x} & \PH{xx,x}\,$\pm$\,\PH{x,x} & \PH{xx,x} & \PH{xx,x} & \PH{xx,x}\,$\pm$\,\PH{x,x} \\
    Domänenmittel          & -- & -- & -- & \PH{xx,x}\,$\pm$\,\PH{x,x} & -- & -- & \PH{xx,x}\,$\pm$\,\PH{x,x} \\
    Sensitivität (ohne 21 Sätze) & 2\,485 & -- & -- & \PH{xx,x} & -- & -- & \PH{xx,x} \\
    \midrule
    MultiNERD-Referenz (PER/ORG/LOC) & 16\,454 & \PH{xx,x} & \PH{xx,x} & \PH{xx,x} & \PH{xx,x} & \PH{xx,x} & \PH{xx,x} \\
    $\Delta$F$_1$ (gepoolt $-$ Referenz) & & & & \PH{$-$xx,x} & & & \PH{$-$xx,x} \\
    \bottomrule
  \end{tabular}
\end{table}

Gegenüber der MultiNERD-Referenz fällt der F$_1$ auf CrossNER bei
DeBERTa-v3-large um \PH{xx,x} Punkte und bei Qwen3.5-27B um \PH{xx,x}
Punkte.
\PH{Beschreiben: Welches Modell verliert weniger? Ist der Unterschied
zwischen den Modellen größer als die Streuung? Welche Domäne ist am
schwersten (vermutlich Naturwissenschaften oder KI wegen vieler
nicht abgebildeter Typen, die als Zusätze zählen) und welche am
leichtesten (vermutlich Politik mit vielen PER/ORG/LOC)? Precision
gegenüber Recall: Verlieren die Modelle vor allem Recall, weil
CrossNER-Typen wie politician auf PER abgebildet werden, oder
Precision, weil misc-Spans als PER/ORG/LOC vorhergesagt werden?}
Der Ausschluss der 21 mit MultiNERD-Trainingssätzen übereinstimmenden
Sätze verändert die gepoolten Werte um \PH{höchstens x,x} Punkte.
Beim LLM waren \PH{x,x}\,\% der Antworten nicht parsebar, \PH{nahezu alle}
davon ohne schließende Klammer; da CrossNER-Sätze \PH{länger sind und
mehr Entitäten enthalten als MultiNERD-Sätze}, wirkt das Limit von 256
Ausgabetokens hier stärker (Abschnitt~\ref{sec:crossner}), und Recall
und F$_1$ des LLM sind entsprechend unterschätzt.
Die Latenz auf CrossNER betrug \PH{xx,x} ms (DeBERTa-v3-large) und
\PH{xxxx} ms (Qwen3.5-27B) je Satz auf derselben H100.
% TODO: Satz zur 1024-Token-Sensitivitätsanalyse nur aufnehmen, wenn sie
% durchgeführt wurde; andernfalls den entsprechenden Absatz in 4.6
% streichen.

\section{Diskussion der Ergebnisse}
\label{sec:erg-diskussion}

Die folgenden Abschnitte ordnen die Befunde entlang der
Forschungsfragen aus Abschnitt~\ref{sec:problem_ziel} ein und
vergleichen sie mit den verwandten Arbeiten aus
Kapitel~\ref{chap:verwandte_arbeiten}.
Abschnitt~\ref{sec:disk-validitaet} benennt, welche Unterschiede
angesichts der Streuung und der bekannten Abweichungen
(Tabelle~\ref{tab:abweichungen}) belastbar sind.

\subsection{Erkennungsleistung im Vergleich zur Literatur}
\label{sec:disk-leistung}

% TODO: Richtung des Hauptbefunds eintragen und den passenden Absatz
% ausformulieren; den nicht zutreffenden Absatz streichen.

Der zentrale Befund zu FF1 lautet: \PH{Das beste LLM (Qwen3.5-27B, QLoRA)
liegt um x,x Punkte vor / hinter DeBERTa-v3-large; der Abstand ist
unterscheidbar / nicht unterscheidbar.}

\PH{Falls das LLM vorn liegt:} Das Ergebnis entspricht dem Muster, das
\textcite{obeidat2025llm} auf vier von fünf biomedizinischen Datensätzen
berichten, dort mit 2 bis 8 Punkten Vorsprung, und dem von
\textcite{zhan2026generativeNER} auf CoNLL-2003 mit 93{,}85 gegenüber
93{,}04 für BERT-MRC.
Der hier gemessene Vorsprung von \PH{x,x} Punkten liegt \PH{am unteren
Rand / innerhalb / oberhalb} dieser Spanne.
Anders als bei Obeidat et al., deren Vorsprung vor allem aus der
Precision stammt, geht er hier \PH{auf Precision / Recall / beides}
zurück.
\PH{Falls der Encoder vorn liegt oder kein Unterschied besteht:} Das
Ergebnis weicht von \textcite{obeidat2025llm} ab, deckt sich aber mit dem
Befund von \textcite{zhan2026generativeNER}, dass LoRA-angepasste LLMs
spezialisierte Modelle auf flachen Datensätzen nur knapp erreichen, und
mit der Beobachtung von \textcite{lu2024llmstruggle}, dass ein
LoRA-angepasstes 7B-Modell hinter einem feinjustierten Encoder bleibt.
Zwei Eigenschaften des Setups sprechen dafür, dass die Encoder hier
begünstigt sind: Die Encoder wurden vollständig feinjustiert, die LLMs
nur über Adapter mit 0{,}5 bis 0{,}8\,\% trainierbarer Parameter
(Tabelle~\ref{tab:lora-hyperparams}); und der Trainingssplit ist mit
131\,280 Sätzen groß genug, um den Vorteil des Vorwissens großer Modelle
zu verkleinern, den \textcite{bogdanov2024nuner} und
\textcite{zaratiana2024gliner} vor allem bei wenigen Beispielen finden.

Unabhängig von der Richtung ist der Abstand zwischen den trainierten
Klassen \PH{klein} im Verhältnis zum Abstand zwischen den Regimen:
Die Anpassung mit QLoRA hebt den F$_1$ um \PH{xx bis xx} Punkte über das
Zero-Shot-Regime.
\PH{Das bestätigt (prüfen: Zero-Shot deutlich unter allen trainierten
Konfigurationen) für ein feingranulares Inventar mit 15 Typen, was
\textcite{keraghel2024recent} mit GPT-4o (65{,}17 gegenüber 90{,}98 auf
CoNLL-2003) und \textcite{tan2023damo} mit ChatGPT auf MultiCoNER II
(14{,}78 gegenüber 92{,}77) für geprompte Modelle berichten: Ohne
aufgabenspezifische Anpassung bleiben LLMs weit hinter überwachten
Systemen zurück.}
Die hier gemessenen Zero-Shot-Werte von \PH{xx bis xx} Punkten liegen
\PH{deutlich unter} dem Zero-Shot-Wert von 76{,}0, den
\textcite{sainz2024gollie} für GoLLIE-7B auf einer MultiNERD-Fassung
berichten; GoLLIE ist jedoch auf Annotationsrichtlinien vieler
Datensätze trainiert und hatte acht der 15 Typen bereits gesehen, sodass
der Vergleich nur die Größenordnung einordnet.
Die trainierten Konfigurationen erreichen Werte \PH{in der Größenordnung}
der 93{,}73 (UniNER-7B) und 91{,}25 (BERT-base), die
\textcite{zhou2024universalner} auf einer MultiNERD-Variante mit
eigener Datenaufbereitung berichten; auch dort ist der Abstand zwischen
LLM und Encoder klein.
Da alle Werte dieser Arbeit die Übereinstimmung mit einer
Silberreferenz messen (Abschnitt~\ref{sec:datenauswahl}), bleibt offen,
ob die Unterschiede bei den verrauschten Typen ANIM, BIO, FOOD und PLANT
Modellfehler oder Annotationsfehler abbilden; \textcite{tedeschi2022multinerd}
berichten für diese Typen auf ihrer manuell geprüften Stichprobe
F$_1$-Werte der Annotation zwischen 0{,}52 und 0{,}71.

\subsection{Ausgabeformat und Formtreue}
\label{sec:disk-format}

Die zweite Hälfte von FF1 betrifft die Formtreue.
Im QLoRA-Regime halten die Modelle das JSON-Format mit Wortindizes
\PH{nahezu vollständig} ein: Parse-Fehlerraten von \PH{unter x}\,\%,
\PH{keine} unbekannten Typen und \PH{nur x} Text-Abweichungen auf 16\,454
Antworten.
Das stützt die in Abschnitt~\ref{sec:llm_ner_related} formulierte
Erwartung, dass Indizes, die explizit in der Eingabe stehen, den von
\textcite{zhan2026generativeNER} für Zeichenoffsets beobachteten
Einbruch vermeiden; der in Tabelle~\ref{tab:haupt} erreichte F$_1$ liegt
\PH{weit über} den rund 54 Punkten, die dort für JSON mit Zeichenoffsets
berichtet werden, und \PH{nahe an} den Inline-Formaten.
\textcite{semin2026spanLabeling} hatten die Wortindex-Variante nicht
evaluiert; die vorliegenden Ergebnisse zeigen, dass sie nach einem
Fine-Tuning \PH{tragfähig} ist.
Zugleich bestätigt sich der von \textcite{dagdelen2024structuredIE}
berichtete Befund, dass feinjustierte Modelle JSON fast immer parsebar
ausgeben, auch ohne erzwungene Formatierung.
Die verbleibenden Parse-Fehler gehen \PH{fast ausschließlich} auf das
Ausgabelimit zurück, nicht auf das Format; sie sind eine Eigenschaft der
Pipeline-Einstellung und würden mit einem höheren Limit
\PH{voraussichtlich} entfallen.

Im Zero-Shot-Regime ist die Formtreue \PH{deutlich schwächer}:
\PH{zusammenfassen, welche Fehlerklassen dominieren}.
Das zeigt, dass Formatvorgabe und Typinventar ohne Training zum
Flaschenhals werden, wie es \textcite{lu2024llmstruggle} für BIO-Tags
beschreiben; die Nachverarbeitung aus Abschnitt~\ref{sec:postprocessing}
rettet gültige Elemente, repariert aber keine fehlerhaften Indizes.
Ob Constrained Decoding diese Fehler ohne Leistungsverlust beseitigen
würde, bleibt offen; \textcite{tam2024speakfreely} berichten
Leistungseinbußen durch Formatzwang, \textcite{semin2026spanLabeling}
dagegen Gewinne durch ihr LogitMatch-Verfahren.
Für die Effizienz ist das Format ein Faktor: Weil die Ausgabelänge nur
von der Zahl der Entitäten abhängt, nicht von der Satzlänge, bleiben die
Latenzen \PH{unter / bei} denen, die \textcite{obeidat2025llm} für
BIO-Ausgaben berichten (40- bis 220-fache Encoder-Latenz); hier beträgt
der Faktor \PH{xx bis xxx}.

\subsection{Modellgröße}
\label{sec:disk-groesse}

Zu FF2 zeigen die Ergebnisse \PH{kleine / deutliche} Größeneffekte
innerhalb der Familien.
\PH{Bei den Encodern: base gegenüber large.}
Das deckt sich mit \textcite{keraghel2024recent}, bei denen
DeBERTa-v3-base und -large je nach Datensatz die Plätze tauschen.
\PH{Bei den LLMs im QLoRA-Regime: Zugewinn 0.8B zu 4B und 4B zu 27B.}
Das Muster entspricht den Beobachtungen bei GoLLIE und UniversalNER,
wo größere Backbones nur wenig hinzugewinnen und der Gewinn durch die
Anpassung den Gewinn durch die Größe übersteigt
\cite{sainz2024gollie,zhou2024universalner}.
Im Zero-Shot-Regime ist der Größeneffekt \PH{größer}, was zu der
Beobachtung passt, dass Fähigkeiten wie das Befolgen komplexer
Formatvorgaben erst ab einer gewissen Größe zuverlässig auftreten
\cite{wei2022emergent}.
Der Vergleich über die Familien hinweg \PH{zeigt, dass die Parameterzahl
allein die Leistung nicht erklärt (prüfen: gilt, wenn DeBERTa-v3-base
Qwen3.5-0.8B erreicht oder übertrifft)}: \PH{DeBERTa-v3-base mit 184 Mio.
Parametern erreicht xx,x Punkte, Qwen3.5-0.8B mit dem Vierfachen im
QLoRA-Regime xx,x.}
Architektur, Vortrainingsziel und Anpassungsverfahren überlagern den
Größenunterschied, wie \textcite{zaratiana2024gliner} und
\textcite{bogdanov2024nuner} für Encoder-Generalisten zeigen.
Die Skalierungsgesetze aus Abschnitt~\ref{sec:model_size} sind auf
diese Befunde nicht übertragbar: Sie beschreiben den Vortrainingsverlust
eigens trainierter Modellreihen, hier wurden veröffentlichte
Checkpoints mit größenabhängigen Hyperparametern und auf verschiedenen
GPUs angepasst.

\subsection{Effizienz und praktische Empfehlungen}
\label{sec:disk-effizienz}

FF3 fragt nach dem günstigsten Verhältnis von Erkennungsleistung und
Ressourcenbedarf unter Speicher- und Latenzbeschränkungen.
Die Effizienzwerte lassen sich zu drei Aussagen verdichten.
Erstens liegen die Encoder bei Latenz und Speicher um \PH{ein bis zwei}
Größenordnungen unter den LLMs: \PH{8,7 ms gegenüber xxx bis xxxx ms;
x,x GB gegenüber x,x bis xx,x GB}.
Zweitens kostet der Leistungsunterschied zwischen den Klassen
\PH{x,x} Punkte, der Effizienzunterschied dagegen den Faktor \PH{xx}
bei der Latenz; dieses Verhältnis entspricht der Schlussfolgerung von
\textcite{obeidat2025llm}, dass Encoder bei kleinen
Leistungsunterschieden oder bei Latenzanforderungen die geeignetere
Wahl bleiben.
Drittens ist das Adaptertraining selbst sparsam: Die QLoRA-Läufe
trainieren 6{,}4 bis 159~Mio. Parameter und passen mit NF4-Basismodell
auf eine GPU, das 27B-Modell benötigt dafür aber \PH{xx} GB während der
Inferenz und \PH{mindestens xx} Stunden Training.

Daraus ergeben sich Empfehlungen, die in Abschnitt~\ref{sec:fazit-zusammenfassung}
zusammengefasst werden:
\PH{(1) Bei Latenzanforderungen im Millisekundenbereich oder GPUs unter
x GB: DeBERTa-v3-base bzw. -large, mit einem Verlust von höchstens x,x
Punkten gegenüber dem besten LLM.
(2) Wenn der Leistungsunterschied von x,x Punkten den Mehraufwand
rechtfertigt und xx GB GPU-Speicher sowie Latenzen von rund x Sekunden
je Satz akzeptabel sind: Qwen3.5-27B mit QLoRA.
(3) Qwen3.5-4B mit QLoRA als Mittelweg, falls es DeBERTa-v3-large
erreicht; andernfalls dominiert.
(4) Zero-Shot nur, wenn keine Trainingsdaten vorliegen; bei 15 Typen
mit xx Punkten nicht konkurrenzfähig.}
Diese Empfehlungen gelten für die hier gemessenen Bedingungen:
Batchgröße~1, ein Satz je Anfrage, die genannten GPU-Typen.
Mit Batching steigt der Durchsatz der LLMs stärker als der der Encoder,
sodass sich die Faktoren bei Offline-Verarbeitung verschieben; dieser
Fall wurde nicht gemessen (Abschnitt~\ref{sec:abgrenzung}).

\subsection{Belastbarkeit der Befunde}
\label{sec:disk-validitaet}

Welche der berichteten Unterschiede tragen, bestimmt sich aus der
Streuung und aus den Abweichungen in Tabelle~\ref{tab:abweichungen}.
Die Standardabweichungen über die Seeds liegen bei \PH{x,x bis x,x}
Punkten; Unterschiede unterhalb von \PH{etwa x} Punkten werden daher
nicht interpretiert.
Davon betroffen sind \PH{der Vergleich DeBERTa-v3-base gegen -large /
der Schritt von 4B auf 27B / der Abstand zwischen DeBERTa-v3-large und
Qwen3.5-27B}.
Belastbar sind dagegen \PH{der Abstand zwischen Zero-Shot- und
QLoRA-Regime, der Abstand zwischen den trainierten Modellen und dem
kleinsten LLM sowie die Effizienzfaktoren}, die um ein Vielfaches über
der Streuung liegen.
Für die Effizienzwerte gilt das nur innerhalb gleicher Hardware; der
hardwaregleiche Vergleich beschränkt sich auf \PH{DeBERTa-v3 und
Qwen3.5-0.8B (QLoRA) auf der H100}.

Die Einschränkungen des Versuchsdesigns, darunter die Silberreferenz,
der einzelne Zero-Shot-Lauf je Größe, die drei GPU-Typen und die
asymmetrische Checkpoint-Auswahl, sind in Tabelle~\ref{tab:abweichungen}
zusammengefasst und werden in Abschnitt~\ref{sec:limitationen} diskutiert.
Keine von ihnen kehrt die Richtung der belastbaren Befunde um; sie
begrenzen aber, wie genau die kleinen Unterschiede zwischen den
trainierten Konfigurationen zu beziffern sind.

% chapters/06_fazit_und_ausblick.tex
%
% Erstentwurf vom 27.10.2026 (KI-gestützt erstellt, siehe Erklärung).
% 6.1 folgt den Forschungsfragen FF1 bis FF3 aus Abschnitt 1.2; die
% Antworten enthalten \PH{...}-Platzhalter, die nach Kapitel 5 zu füllen
% sind. 6.2 und 6.3 sind aus Kapitel 4 und 5 ableitbar und ausformuliert.
\providecommand{\PH}[1]{\textcolor{red}{\textbf{[#1]}}}

\chapter{Fazit und Ausblick}
\label{chap:fazit}

Diese Arbeit hat zwei Klassen vortrainierter Sprachmodelle für
feingranulare Named Entity Recognition in einer gemeinsamen Pipeline
verglichen: DeBERTa-v3-base und -large mit vollständigem Fine-Tuning
sowie Qwen3.5-0.8B, -4B und -27B im Zero-Shot-Regime und nach einer
Anpassung mit QLoRA.
Datenbasis waren die offiziellen englischen Splits von MultiNERD mit
15 Entitätstypen; bewertet wurde mit strikten Entity-Level-Metriken,
ergänzt um Trainingsdauer, Inferenzlatenz, GPU-Spitzenspeicher und die
Formtreue der generierten Ausgaben.
Abschnitt~\ref{sec:fazit-zusammenfassung} beantwortet die drei
Forschungsfragen aus Abschnitt~\ref{sec:problem_ziel},
Abschnitt~\ref{sec:limitationen} benennt die Grenzen der Untersuchung,
Abschnitt~\ref{sec:ausblick} skizziert weiterführende Arbeiten.

\section{Zusammenfassung der zentralen Ergebnisse}
\label{sec:fazit-zusammenfassung}

% TODO: Zahlen aus tab:haupt, tab:effizienz und tab:formtreue-erg
% übernehmen; die Richtung der Aussagen an Abschnitt 5.7 angleichen.
% Keine neuen Ergebnisse einführen, die nicht in Kapitel 5 stehen.
% Nach dem Füllen jeden FF-Absatz auf sechs bis acht Sätze kürzen; Typ-
% und Formtreue-Details bleiben in Kapitel 5.

\paragraph{FF1: Erkennungsleistung und Formtreue.}
Auf dem englischen MultiNERD-Testsplit erreicht das beste LLM,
Qwen3.5-27B im QLoRA-Regime, einen Micro-F$_1$ von \PH{xx,x} Punkten,
der beste Encoder, DeBERTa-v3-large, einen von \PH{xx,x}.
Der Abstand von \PH{x,x} Punkten ist \PH{größer / kleiner} als die
Streuung über die Seeds; \PH{das LLM / der Encoder / keine der beiden
Klassen} ist damit unterscheidbar besser.
Die kleineren LLMs erreichen nach der Anpassung \PH{xx,x} (4B) und
\PH{xx,x} Punkte (0.8B) und liegen damit \PH{über / unter / auf dem
Niveau von} DeBERTa-v3-base mit \PH{xx,x}.
Ohne Training fallen die LLMs auf \PH{xx} bis \PH{xx} Punkte zurück;
\PH{die Anpassung mit QLoRA ist damit der mit Abstand wirksamste Faktor
im Vergleich, nicht die Modellklasse und nicht die Modellgröße}.
Bei den drei häufigen Typen PER, ORG und LOC unterscheiden sich die
trainierten Konfigurationen um \PH{höchstens x,x} Punkte; die
Unterschiede konzentrieren sich auf \PH{die verrauschten Typen ANIM, FOOD
und PLANT sowie auf DIS, EVE und MEDIA}.
Die Formtreue ist nach dem Fine-Tuning \PH{hoch}: Weniger als
\PH{x}\,\% der Antworten sind nicht parsebar, unbekannte Typen und
Text-Abweichungen treten \PH{praktisch nicht mehr} auf, und die
verbleibenden Parse-Fehler gehen \PH{nahezu vollständig} auf das Limit
von 256 Ausgabetokens zurück.
Im Zero-Shot-Regime ist die Formtreue \PH{deutlich schwächer} und
\PH{beim 0.8B-Modell der begrenzende Faktor}.
Das JSON-Format mit expliziten Wortindizes hat sich als \PH{tragfähig}
erwiesen.

\paragraph{FF2: Modellgröße und Ressourcen.}
Innerhalb der Encoder-Familie bringt der Schritt von 184 auf 434~Mio.
Parameter \PH{x,x} Punkte, innerhalb der LLM-Familie im QLoRA-Regime der
Schritt von 0{,}75 auf 4{,}21~Mrd. \PH{x,x} Punkte und der von 4{,}21 auf
26{,}90~Mrd. weitere \PH{x,x} Punkte.
\PH{Die Zugewinne nehmen mit der Größe ab / Der Schritt auf 27B ist nicht
unterscheidbar.}
Dem stehen Ressourcen gegenüber, die mit der Größe \PH{nahezu linear}
wachsen: Die Inferenzlatenz je Satz steigt von \PH{8,7}~ms
(DeBERTa-v3-base) über \PH{xx,x}~ms (DeBERTa-v3-large) und \PH{xxx}~ms
(Qwen3.5-0.8B) auf \PH{xxxx}~ms (Qwen3.5-27B), der GPU-Spitzenspeicher
der Inferenz von \PH{x,x} auf \PH{xx,x}~GB, die Trainingsdauer von
\PH{x,x} auf mindestens \PH{xx} Stunden.
Im Zero-Shot-Regime ist der Größeneffekt \PH{größer}: \PH{Erst ab 4B /
27B} erzeugen die Modelle in nennenswertem Umfang gültige Ausgaben.
Über die Familien hinweg erklärt die Parameterzahl die Leistung nicht:
DeBERTa-v3-base erreicht mit 184~Mio. Parametern \PH{mehr / ebenso viel}
wie Qwen3.5-0.8B mit dem Vierfachen.

\paragraph{FF3: Empfehlungen.}
\PH{Unter Latenz- oder Speicherbeschränkungen ist DeBERTa-v3 die
effizientere Wahl}: \PH{DeBERTa-v3-large} erreicht \PH{xx,x} Punkte bei \PH{xx}~ms je
Satz und \PH{x,x}~GB GPU-Speicher und verliert gegenüber dem besten LLM
\PH{höchstens x,x} Punkte, bei einem \PH{xx}-fach geringeren
Latenz- und einem \PH{x}-fach geringeren Speicherbedarf.
Ein LLM rechtfertigt den Mehraufwand nur, wenn \PH{der Leistungsvorsprung
von x,x Punkten für die Anwendung zählt und} Latenzen von \PH{rund x}
Sekunden je Satz sowie \PH{xx}~GB GPU-Speicher akzeptabel sind; dann ist
Qwen3.5-27B mit QLoRA die beste Konfiguration, \PH{während das 4B-Modell
einen Mittelweg bietet / während das 4B-Modell von DeBERTa-v3-large
dominiert wird}.
Das Zero-Shot-Regime ist bei 15 Typen \PH{nicht konkurrenzfähig} und nur
dann sinnvoll, wenn keine annotierten Daten vorliegen.
Diese Empfehlungen gelten für die Verarbeitung einzelner Sätze auf den
verwendeten GPU-Typen; für Offline-Verarbeitung mit Batching verschieben
sich die Faktoren zugunsten der LLMs, ohne dass dies hier gemessen wurde.

\section{Limitationen der Arbeit}
\label{sec:limitationen}

Die Aussagekraft der Ergebnisse ist durch das Versuchsdesign und die
Datenbasis in mehreren Punkten begrenzt; Tabelle~\ref{tab:abweichungen}
fasst die bekannten Abweichungen vom idealisierten Versuchsdesign
zusammen.

\paragraph{Datenbasis.}
Alle Werte messen die Übereinstimmung mit einer automatisch erzeugten
Referenz.
Die MultiNERD-Splits sind per Distant Supervision annotiert, und die
Originalarbeit beziffert die Qualität der Annotation für ANIM, BIO, FOOD
und PLANT mit F$_1$-Werten zwischen 0{,}52 und 0{,}71 gegenüber einer
manuellen Prüfung \cite{tedeschi2022multinerd}.
Ein Teil der gemessenen Fehler bei diesen Typen kann daher auf die
Referenz zurückgehen, und ein Modell, das die Silberannotation
\enquote{zu gut} reproduziert, lernt auch deren Fehler.
Die Untersuchung beschränkt sich zudem auf einen Datensatz, eine Sprache
und ein Genre; Wikipedia- und Wikinews-Texte sind mit hoher
Wahrscheinlichkeit Teil der Vortrainingsdaten der LLMs, und eine
Kontamination wurde nicht empirisch geprüft
(Abschnitt~\ref{sec:research_gap}).
Der CrossNER-Transfertest mildert die Beschränkung auf ein Genre nur
teilweise, da er auf drei Grobtypen reduziert ist und 21 seiner Sätze
mit MultiNERD-Trainingssätzen übereinstimmen.

\paragraph{Vergleichsbedingungen.}
Die Arbeit vergleicht Konfigurationen, nicht isolierte Faktoren.
Die Encoder wurden vollständig feinjustiert, die LLMs nur über Adapter
mit 0{,}5 bis 0{,}8\,\% trainierbarer Parameter; das entspricht der
praxisüblichen Anpassung je Modellklasse, bedeutet aber, dass ein
vollständiges Fine-Tuning der LLMs oder höhere LoRA-Ränge zu anderen
Ergebnissen führen könnten.
Die Adapter decken wegen der hybriden Architektur nur die
Feed-Forward-Netze aller Schichten und die Attention jeder vierten
Schicht ab, und die QLoRA-Umsetzung weicht in zwei Punkten vom
Standardrezept ab (Abschnitt~\ref{sec:lora-konfiguration}).
Hyperparameter wurden je Modellgröße nach üblichen Werten gesetzt und
nicht gesucht; für keine Konfiguration liegt daher ein Nachweis vor,
dass sie ihr Optimum erreicht hat.
Die Checkpoint-Auswahl ist asymmetrisch: Die Encoder wurden auf dem
vollständigen Validierungssplit ausgewählt, die LLM-Adapter anhand von
200 Validierungssätzen.
Der Trainingsverlust der LLMs umfasst den konstanten Prompt, sodass ein
Teil des Lernsignals auf dessen Reproduktion entfällt.

\paragraph{Zero-Shot-Regime.}
Die Zero-Shot-Ergebnisse beruhen auf je einem Lauf mit Sampling bei
Temperatur 1{,}0 und sind damit nicht deterministisch; eine Streuung
über Wiederholungen wurde nicht gemessen.
Prompt und Decoding-Einstellungen wurden nach Beobachtungen auf dem
Testsplit festgelegt, und die Zero-Shot-Läufe verwenden einen anderen
Prompt, eine andere Präzision und ein anderes Decoding als die
QLoRA-Läufe.
Die Zero-Shot-Werte sind deshalb als optimistische Einzelwerte für genau
dieses Format zu lesen; mit Demonstrationsbeispielen oder Constrained
Decoding wären höhere Werte zu erwarten.

\paragraph{Effizienzmessung.}
Die Konfigurationen liefen auf drei GPU-Typen; Latenz- und
Speicherwerte verschiedener Konfigurationen enthalten daher
Hardwareeffekte, die den gemessenen Unterschied zwischen A100 und H100
beim Encoder (17{,}6 gegenüber 8{,}7~ms) erreichen können.
Gemessen wurde die Latenz bei Batchgröße~1 und einem Satz je Anfrage;
der mit Batching erreichbare Durchsatz, der Speicherbedarf während des
Trainings sowie Energieverbrauch und Kosten wurden nicht erfasst.
Die Trainingsdauer des 27B-Modells ist nur als Untergrenze bekannt.
Das Limit von 256 Ausgabetokens führt bei langen Sätzen zu
abgeschnittenen, nicht parsebaren Antworten und unterschätzt Recall und
F$_1$ der LLMs\PH{, auf CrossNER stärker als auf MultiNERD}.

\paragraph{Statistische Absicherung.}
Jede trainierte Konfiguration wurde mit drei Seeds trainiert; auf
inferenzstatistische Tests wurde verzichtet.
Unterschiede in der Größenordnung der Standardabweichungen, darunter
\PH{der Vergleich der beiden Encoder und der Schritt von 4B auf 27B},
lassen sich mit diesem Umfang nicht belegen.
Der Testsplit wurde außerdem bereits in den explorativen Vorläufen
ausgewertet; für die trainierten Konfigurationen blieben die
Hyperparameter danach bis auf die Epochenzahl des 27B-Modells
unverändert, eine strenge Trennung von Entwicklung und Test ist aber
nicht gegeben.

\section{Ausblick}
\label{sec:ausblick}

Aus den Ergebnissen und ihren Grenzen ergeben sich mehrere
Anschlussfragen.

\paragraph{Gold-Referenz.}
Die wichtigste Ergänzung wäre eine manuell geprüfte Teststichprobe,
etwa im Umfang der 1\,000 Sätze, die \textcite{tedeschi2022multinerd}
für ihre Qualitätsmessung verwendet haben.
Damit ließe sich trennen, ob die Unterschiede zwischen den Modellklassen
bei ANIM, FOOD und PLANT Modellfehler oder Annotationsfehler sind, und
ob ein Modell, das von der Silberreferenz abweicht, in Wahrheit die
bessere Annotation liefert.

\paragraph{Anpassungsverfahren.}
Der Vergleich zwischen vollständigem Fine-Tuning und Adaptern ließe sich
auf der LLM-Seite schließen, indem das 0.8B-Modell vollständig
feinjustiert wird; umgekehrt könnten die Encoder mit LoRA angepasst
werden, wie es \textcite{hu2021lora} für RoBERTa und DeBERTa zeigen.
Auf der Adapterseite bleiben höhere Ränge, Adapter in den
Gated-DeltaNet-Schichten, ein auf die Antwort beschränkter Verlust und
der Verzicht auf Quantisierung offen.
Eine dritte Klasse, Encoder-Generalisten wie GLiNER
\cite{zaratiana2024gliner} oder mit LLM-Annotationen vortrainierte
Encoder wie NuNER \cite{bogdanov2024nuner}, könnte den Vergleich um
Modelle ergänzen, die Vorwissen aus LLMs mit der Effizienz von Encodern
verbinden.

\paragraph{Zero-Shot und Ausgabeformat.}
Für das Zero-Shot-Regime wären mehrere Samples je Satz oder ein
Greedy-Kontrolllauf bei allen Größen nötig, um die Streuung zu
beziffern; Few-Shot-Prompting mit Demonstrationsbeispielen und
Constrained Decoding, etwa das LogitMatch-Verfahren von
\textcite{semin2026spanLabeling}, sind die naheliegenden Erweiterungen.
Für das Ausgabeformat steht ein Vergleich der Wortindex-Variante mit
Inline- und Matching-Formaten unter gleichen Trainingsbedingungen aus,
der zugleich den Einfluss der Ausgabelänge auf die Latenz beziffern
würde \cite{zhan2026generativeNER,lu2024padellm}.
Ein höheres Ausgabelimit würde die Abbrüche bei langen Sätzen
beseitigen.

\paragraph{Effizienz.}
Eine Messung auf einheitlicher Hardware, mit Batching und mit
Energieverbrauch würde die Empfehlungen aus FF3 auf
Offline-Verarbeitung ausweiten.
Für die Inferenz der LLMs bieten sich zudem spezialisierte
Inferenzsysteme und das Zusammenführen der Adapter mit dem Basismodell
an, die hier nicht verwendet wurden.

\paragraph{Übertragbarkeit.}
MultiNERD umfasst zehn Sprachen mit demselben Typinventar.
Die Pipeline dieser Arbeit ließe sich ohne Änderung des Formats auf die
übrigen Sprachen anwenden und würde zeigen, ob der hier für Englisch
gemessene Abstand zwischen Encodern und LLMs in Sprachen mit weniger
Vortrainingsdaten anders ausfällt.
Der CrossNER-Versuch legt nahe, den Domänenwechsel mit dem vollständigen
Typinventar und mit wenigen Zieldomänenbeispielen zu untersuchen, wie es
\textcite{ashok2023promptner} und \textcite{sainz2024gollie} für
LLM-Ansätze tun.


% =========================
% Literaturverzeichnis
% =========================
\printbibliography[
  heading=bibintoc,
  title={Literaturverzeichnis}
]

% =========================
% Anhang
% =========================
\appendix
% chapters/anhang.tex
%
% A.1 dokumentiert die Reproduzierbarkeit der in Kapitel 4 beschriebenen
% Läufe (Labels app:reproduzierbarkeit, lst:prompt-glossar), A.2 ergänzende
% Ergebnistabellen (Tabellenkörper aus thesis/generated/).
% Benötigte Pakete in der Präambel: xcolor, tabularx, listings, booktabs.

\chapter{Anhang}
\label{chap:anhang}

\section{Reproduzierbarkeit der Läufe}
\label{app:reproduzierbarkeit}

Dieser Abschnitt ordnet den in Kapitel~\ref{chap:methodik} beschriebenen
Läufen Konfigurationsdateien, Codestände, Revisionen und
Ausgabeverzeichnisse zu und dokumentiert nachrangige Einstellungen.
Quellcode, Konfigurationsdateien, Auswertungsskripte und die aus den
Laufartefakten erzeugten Tabellen liegen im Repository
\url{https://github.com/luajosman/bachelor-thesis-llm-vs-encoder-ner}
(Commits nach Tabelle~\ref{tab:app-revisionen}; der Stand der
Zero-Shot-Läufe ist als Tag \texttt{zeroshot-runs-137bffd} markiert).
Konfigurationsdateien liegen dort unter \texttt{ba-ner/configs/}, die
Auswertung unter \texttt{thesis/scripts/aggregate\_results.py} mit den
Kennzahlen in \texttt{thesis/generated/summary.json}.
Die Laufartefakte selbst (Vorhersagen, Metriken, Adapter; rund 2~GB ohne
Modellgewichte) liegen relativ zum Ergebnisverzeichnis unter
\texttt{results/multinerd/} bzw. \texttt{results/crossner/} und werden
den Prüfenden auf Anfrage bereitgestellt.

\paragraph{Läufe und Artefakte.}
Tabelle~\ref{tab:app-laeufe} nennt für jede Konfiguration die
Konfigurationsdatei und das Ausgabeverzeichnis des Laufs mit Seed~42.
Für die Seeds 123 und 456 ersetzen die Konfigurationsdateien das Suffix
\texttt{\_canonical} durch \texttt{\_seed123} bzw. \texttt{\_seed456};
beim 27B-Modell wird es an \texttt{\_3ep} angehängt.
Die Ausgabeverzeichnisse erhalten das Suffix \texttt{-seed123} bzw.
\texttt{-seed456}.
Die Ausgabeverzeichnisse enthalten folgende Artefakte:
{\raggedright
\begin{itemize}
    \item alle Läufe: Testvorhersagen
    (\texttt{test\_\allowbreak predictions.json}) und Inferenzmetriken mit
    eingebetteten Laufmetadaten (\texttt{inference\_\allowbreak metrics.yaml});
    \item zusätzlich bei den trainierten Konfigurationen:
    Trainingsergebnisse (\texttt{results.yaml}), Trainingsverlauf
    (\texttt{training\_\allowbreak history.json}), ausgewählter Checkpoint
    (\texttt{best\_\allowbreak model} bzw.
    \texttt{best\_\allowbreak lora\_\allowbreak adapter}; Adapter 12,8, 42,5
    bzw. 637,6~MB), aufgelöste Konfiguration
    (\texttt{config\_\allowbreak resolved.yaml}), wissenschaftlicher
    Vertrag mit Code- und Revisions-Hashes
    (\texttt{scientific\_\allowbreak contract.yaml}), Paketversionen
    (\texttt{environment/}), Metadaten der Laufverwaltung
    (\texttt{run\_\allowbreak metadata.yaml}, \texttt{status.json};
    bei den 27B-Läufen zusätzlich \texttt{recovery-*.yaml}) sowie
    Evaluationsergebnisse einschließlich typbezogener Werte
    (\texttt{evaluation\_\allowbreak metrics.yaml}).
\end{itemize}\par}
Die CrossNER-Evaluationen verwenden die Checkpoints der Hauptstudie und
legen ihre Ergebnisse unter \texttt{results/crossner/} im Verzeichnis des
jeweiligen Laufs mit einem Unterordner je Domäne ab.

\begin{table}[ht]
  \centering
  \footnotesize
  \caption[Konfigurationsdateien, Ausgabeverzeichnisse und Codestände der Läufe]{Konfigurationsdateien, Ausgabeverzeichnisse (jeweils Seed~42)
  und Codestände der ausgewerteten Läufe; vollständige Commit-Hashes in
  Tabelle~\ref{tab:app-revisionen}.}
  \label{tab:app-laeufe}
  \setlength{\tabcolsep}{3pt}
  \begin{tabularx}{\linewidth}{@{}>{\raggedright\arraybackslash}p{3.3cm}%
      >{\raggedright\arraybackslash}p{4.35cm}>{\raggedright\arraybackslash}X%
      l@{}}
    \toprule
    \textbf{Konfiguration} & \textbf{Konfigurationsdatei} &
    \textbf{Ausgabeverzeichnis} & \textbf{Commit} \\
    \midrule
    DeBERTa-v3-base       & \texttt{deberta\_base\_canonical}  & \texttt{deberta-v3-base-canonical}  & \texttt{179eed8} \\
    DeBERTa-v3-large      & \texttt{deberta\_large\_canonical} & \texttt{deberta-v3-large-canonical} & \texttt{179eed8} \\
    Qwen3.5-0.8B QLoRA    & \texttt{qwen35\_08b\_canonical}    & \texttt{qwen35-08b-qlora-canonical} & \texttt{179eed8} \\
    Qwen3.5-4B QLoRA      & \texttt{qwen35\_4b\_canonical}     & \texttt{qwen35-4b-qlora-canonical}  & \texttt{179eed8} \\
    Qwen3.5-27B QLoRA     & \texttt{qwen35\_27b\_3ep}          & \texttt{qwen35-27b-qlora-3ep}       & \texttt{179eed8}$^{a}$ \\
    Qwen3.5-0.8B Zero-Shot & \texttt{qwen35\_08b\_zeroshot}$^{b}$ & \texttt{qwen35-08b-zeroshot}     & \texttt{137bffd} \\
    Qwen3.5-4B Zero-Shot  & \texttt{qwen35\_4b\_zeroshot}$^{b}$  & \texttt{qwen35-4b-zeroshot}      & \texttt{137bffd} \\
    Qwen3.5-27B Zero-Shot & \texttt{qwen35\_27b\_zeroshot}$^{b}$ & \texttt{qwen35-27b-zeroshot}     & \texttt{137bffd} \\
    CrossNER           & wie Hauptstudie                    & \texttt{results/crossner/}\dots      & \texttt{b620041} \\
    CrossNER, 1\,024 Tokens (Vorstudie) & wie Hauptstudie & \texttt{results/\allowbreak crossner\_output\_limit\_1024/\allowbreak smoke/}\dots & \texttt{ec49279} \\
    CrossNER, 1\,024 Tokens (alle Sätze, drei Seeds) & wie Hauptstudie & \texttt{results/\allowbreak crossner\_output\_limit\_1024/}\dots & \texttt{ff91342} \\
    Nachinferenz MultiNERD, 1\,024 Tokens (neun QLoRA-Läufe) & wie Hauptstudie & \texttt{results/\allowbreak multinerd\_output\_limit\_1024/}\dots & \texttt{ff91342} \\
    Greedy-Kontrolle Qwen3.5-0.8B & \texttt{qwen35\_08b\_zeroshot} & \texttt{qwen35-08b-zeroshot-greedy-v1} & \texttt{4679bf1}$^{c}$ \\
    Vorlauf DeBERTa-v3-base (Latenz auf A100) & \texttt{deberta\_base} & \texttt{deberta-v3-base} & \texttt{cdf10db}$^{c}$ \\
    \bottomrule
    \multicolumn{4}{@{}p{\linewidth}@{}}{\scriptsize $^{a}$ Start mit \texttt{179eed8};
      Fortsetzung nach den drei Unterbrechungen und Testinferenz mit
      \texttt{c01bcbe}.} \\
    \multicolumn{4}{@{}p{\linewidth}@{}}{\scriptsize $^{b}$ Dateistand von Commit
      \texttt{137bffd} (Tag \texttt{zeroshot-runs-137bffd}); der aktuelle
      Stand im Repository (Greedy-Decoding) weicht ab.} \\
    \multicolumn{4}{@{}p{\linewidth}@{}}{\scriptsize $^{c}$ Arbeitskopie mit lokalen,
      nicht erhaltenen Änderungen; Kontroll- bzw. Vorlauf außerhalb der
      Hauptauswertung.} \\
  \end{tabularx}
\end{table}

\paragraph{Ausführungszeiten und Cluster-Jobs.}
Die Trainings- und Inferenzläufe der Hauptstudie liefen im August 2026
auf dem DFKI-Cluster (Partitionen und Hauptspeicher in
Tabelle~\ref{tab:app-einstellungen}); ihre Slurm-Abrechnung liegt als
\texttt{results/slurm\_accounting\_2026-08.csv} vor, die 27B-Läufe
umfassen die Jobs 3254977, 3254980 und 3254983 (wiederholt eingereiht)
sowie 3267440, 3267443 und 3267446.
Die drei Zero-Shot-Läufe und der Greedy-Kontrolllauf entstanden am
1.~August bzw. 29.~Juli 2026 (27B auf einer H100~NVL, 0.8B und 4B auf
einer A100); ihre Job-IDs sind in den Laufartefakten nicht
protokolliert.
Die CrossNER-Evaluationen mit 256 Tokens liefen im Oktober 2026.
Die Nachinferenz mit 1\,024 Tokens lief am 28.~Oktober 2026 (Jobs
3673515 bis 3673523, H100~HBM3, Commit \texttt{ff91342}), der
vollständige CrossNER-Neulauf mit 1\,024 Tokens vom 28.\ bis
29.~Oktober 2026 (Jobs 3673475 bis 3673477, Aggregation 3673478); ein
erster Versuch (Jobs 3672273 bis 3672276) brach wegen eines
abweichenden Codestands vor der Inferenz ab.

\paragraph{Explorative Vorläufe.}
Tabelle~\ref{tab:app-vorlaeufe} dokumentiert die explorativen Vorläufe
(Abschnitt~\ref{sec:datenaufteilung}), je ein Lauf mit Seed~42, die auf
dem Testsplit ausgewertet wurden und nicht in die Hauptauswertung
eingehen; ihre Abweichung vom jeweiligen Seed-42-Lauf der Hauptstudie
beträgt höchstens 0,12 Punkte.

\begin{table}[ht]
  \centering
  \footnotesize
  \caption[Explorative Vorläufe]{Explorative Vorläufe (Seed~42) und ihr
  Mikro-F$_1$ auf dem MultiNERD-Testsplit (256 Ausgabetokens) gegenüber
  dem Seed-42-Lauf der Hauptstudie.}
  \label{tab:app-vorlaeufe}
  \setlength{\tabcolsep}{4pt}
  \begin{tabular}{llllrrr}
    \toprule
    \textbf{Vorlauf} & \textbf{Commit} & \textbf{GPU} & \textbf{Epochen} & \textbf{Mikro-F$_1$} &
      \textbf{Hauptstudie} & $\Delta$ \\
    \midrule
    DeBERTa-v3-base & \texttt{cdf10db}$^{a}$ & A100 & 5 & 95,58 & 95,49 & $+$0,10 \\
    DeBERTa-v3-large & \texttt{06b97be} & A100 & 4 & 95,84 & 95,83 & $+$0,01 \\
    Qwen3.5-0.8B QLoRA & \texttt{4679bf1} & A100 & 3 & 93,32 & 93,24 & $+$0,07 \\
    Qwen3.5-4B QLoRA & \texttt{137bffd} & A100 & 3 & 94,05 & 93,93 & $+$0,12 \\
    Qwen3.5-27B QLoRA & \texttt{137bffd}$^{a}$ & H100 NVL & 2 & 94,54 & 94,53 & $+$0,00 \\
    \bottomrule
    \multicolumn{7}{@{}l}{\scriptsize $^{a}$ Arbeitskopie mit lokalen, nicht erhaltenen \"Anderungen. Epochen: gew\"ahlte} \\
    \multicolumn{7}{@{}l}{\scriptsize Epoche (QLoRA) bzw. trainierte Epochen (Encoder); Hauptstudie: Seed-42-Lauf; $\Delta$ =} \\
    \multicolumn{7}{@{}l}{\scriptsize Vorlauf minus Hauptstudie.} \\
\end{tabular}

\end{table}

\begin{table}[ht]
  \centering
  \footnotesize
  \caption[Vollständige Commits und Revisionen]{Vollständige Commits und Revisionen.
  Die Modell- und Datensatzrevisionen gelten für die Trainingsläufe der
  Hauptstudie, den CrossNER-Versuch und die Sensitivitätsanalysen; die
  Zero-Shot-Läufe luden Modelle und Datensatz ohne explizite
  Revisionsangabe.}
  \label{tab:app-revisionen}
  \begin{tabular}{ll}
    \toprule
    \textbf{Komponente} & \textbf{Commit bzw. Revision} \\
    \midrule
    Code: Trainingsläufe der Hauptstudie     & \texttt{179eed8c68b49f51fff5498e411db22c88ef3c9a} \\
    Code: Fortsetzung und Inferenz 27B-QLoRA & \texttt{c01bcbe0f1be4356a2cf24b1bb501e03b7ab63be} \\
    Code: Zero-Shot-Läufe                    & \texttt{137bffdd9b4351b35a31e4550bde0c15b3c53a19} \\
    Code: CrossNER-Evaluationen              & \texttt{b620041ce75878603b37763831204f26fac7ab91} \\
    Code: CrossNER mit 1\,024 Tokens (Vorstudie) & \texttt{ec492799551cbfc9c47550e25006f11386f99ceb} \\
    Code: Sensitivitätsanalysen mit 1\,024 Tokens & \texttt{ff91342f1f206572c5e930e0be7d1929362ccb9c} \\
    Datensatz \texttt{Babelscape/multinerd}  & \texttt{2814b78e7af4b5a1f1886fe7ad49632de4d9dd25} \\
    CrossNER-Daten (Commit in \texttt{zliucr/CrossNER}) & \texttt{2e7ba2a7798c961e3f29fbc51252c5a8d40224bf} \\
    \texttt{Qwen/Qwen3.5-0.8B}               & \texttt{2fc06364715b967f1860aea9cf38778875588b17} \\
    \texttt{Qwen/Qwen3.5-4B}                 & \texttt{851bf6e806efd8d0a36b00ddf55e13ccb7b8cd0a} \\
    \texttt{Qwen/Qwen3.5-27B}                & \texttt{fc05daec18b0a78c049392ed2e771dde82bdf654} \\
    \texttt{microsoft/deberta-v3-base}       & \texttt{8ccc9b6f36199bec6961081d44eb72fb3f7353f3} \\
    \texttt{microsoft/deberta-v3-large}      & \texttt{64a8c8eab3e352a784c658aef62be1662607476f} \\
    \bottomrule
  \end{tabular}
\end{table}

\paragraph{Nachrangige Einstellungen.}
Tabelle~\ref{tab:app-einstellungen} ergänzt die Hyperparameter aus
Kapitel~\ref{chap:methodik} um Einstellungen, die das Ergebnis nur
mittelbar betreffen oder die Ausführung auf dem Cluster beschreiben.

\begin{table}[ht]
  \centering
  \footnotesize
  \caption[Nachrangige Trainings-, Auswertungs- und Ausführungseinstellungen]{Nachrangige Trainings-, Auswertungs- und
  Ausführungseinstellungen der trainierten Konfigurationen.}
  \label{tab:app-einstellungen}
  \setlength{\tabcolsep}{4pt}
  \begin{tabular}{lccccc}
    \toprule
    & \multicolumn{2}{c}{\textbf{DeBERTa-v3}} &
      \multicolumn{3}{c}{\textbf{Qwen3.5 (QLoRA)}} \\
    \textbf{Einstellung} & \textbf{base} & \textbf{large} &
      \textbf{0.8B} & \textbf{4B} & \textbf{27B} \\
    \midrule
    Batchgröße Validierung & 32 & 16 & 8 & 4 & 1 \\
    Generative Validierung & -- & -- &
      \multicolumn{3}{c}{200 Sätze, Batchgröße 1, greedy} \\
    Checkpoint-Intervall & Epoche & Epoche &
      \multicolumn{3}{c}{500 Schritte} \\
    Max.\ gespeicherte Checkpoints & \multicolumn{5}{c}{2} \\
    Logging-Intervall (Schritte) & 50 & 50 & 25 & 25 & 25 \\
    AdamW $\beta_1$; $\beta_2$; $\epsilon$ &
      \multicolumn{5}{c}{0,9; 0,999; $10^{-8}$} \\
    Gradient Clipping (max.\ Norm) & \multicolumn{5}{c}{1,0} \\
    Gradient Checkpointing & nein & nein &
      \multicolumn{3}{c}{ja (nicht-reentrant)} \\
    Dropout & \multicolumn{2}{c}{0,1 (Encoder und Kopf)} &
      \multicolumn{3}{c}{0,05 (LoRA)} \\
    Optimierungsschritte gesamt & 41\,025 & 32\,820 & 24\,615 & 24\,615 &
      24\,615 \\
    SLURM-Partition & H100 & H100 & H100 & A100-80GB & H100-Trails \\
    Angeforderter Hauptspeicher & 32\,GB & 48\,GB & 64\,GB & 96\,GB &
      160\,GB \\
    \bottomrule
  \end{tabular}
\end{table}

\paragraph{Zero-Shot-Prompt.}
Listing~\ref{lst:prompt-glossar} zeigt die Systemnachricht des
Glossar-Profils; Nutzernachricht und Chat-Template entsprechen
Listing~\ref{lst:prompt-task}.

\begin{lstlisting}[basicstyle=\ttfamily\scriptsize, breaklines=true,
  breakatwhitespace=true, columns=fullflexible, keepspaces=true,
  frame=single, captionpos=b, float=htbp,
  postbreak=\mbox{\textcolor{gray}{$\hookrightarrow$}\space},
  caption={Systemnachricht des im Zero-Shot-Regime verwendeten
  Glossar-Profils.},
  label={lst:prompt-glossar}]
You are an expert Named Entity Recognition (NER) system. Find every named-entity mention in the numbered token list. Named entities are specific named people, organizations, places, animals, biological entities, celestial bodies, diseases, events, foods, instruments, media or creative works, mythological entities, plants, times or dates, and vehicles. Do not label ordinary common nouns unless they are part of a specific named entity.
Use these exact MultiNERD labels: PER = Person; ORG = Organization; LOC = Location; ANIM = Animal; BIO = Biological entity; CEL = Celestial body; DIS = Disease; EVE = Event; FOOD = Food; INST = Instrument; MEDIA = Media or creative work; MYTH = Mythological entity; PLANT = Plant; TIME = Time or date; VEHI = Vehicle.
Return ONLY a JSON array. Do not use Markdown code blocks and do not add explanations. Extract all occurrences, including multi-token entities.
Each entity object must have exactly these fields: "start_token" (integer, inclusive), "end_token" (integer, exclusive), "text" (the exact text covered by the token span), and "type" (one of: PER, ORG, LOC, ANIM, BIO, CEL, DIS, EVE, FOOD, INST, MEDIA, MYTH, PLANT, TIME, VEHI).
Token indices must refer exactly to the numbered token list shown by the user. Check the entire token list before answering. Return [] only when none of the listed named-entity categories occurs.
\end{lstlisting}

\clearpage
\section{Ergänzende Ergebnistabellen}
\label{app:ergebnistabellen}

% Tabellenkörper werden von thesis/scripts/aggregate_results.py erzeugt
% (thesis/generated/); die Beispiele in tab:app-beispiele stammen aus den
% Testvorhersagen der genannten Läufe (Listenindex in test_predictions.json).

Tabelle~\ref{tab:app-seeds} nennt die Testwerte der einzelnen Seeds, aus
denen die Mittelwerte in Tabelle~\ref{tab:haupt} gebildet wurden, sowie
den Validierungs-F$_1$ und die Epoche des jeweils gewählten Checkpoints;
Tabelle~\ref{tab:app-limit-seeds} ergänzt die Nachinferenz mit
1\,024 Tokens (Tabelle~\ref{tab:limit-multinerd}) um die Werte je Seed;
Tabelle~\ref{tab:app-typen-prf} ergänzt Tabelle~\ref{tab:typen} um
Precision und Recall je Entitätstyp; die
Tabellen~\ref{tab:app-crossner-seeds} und~\ref{tab:app-crossner-typen}
ergänzen Tabelle~\ref{tab:crossner-erg} um die Werte der einzelnen Seeds
und um die Grobtypen je Domäne; Tabelle~\ref{tab:app-beispiele}
zeigt Beispiele typischer Fehler je Fehlerklasse aus
Abschnitt~\ref{sec:erg-fehlertypen}.

\begin{table}[ht]
  \centering
  \footnotesize
  \caption[Mikro-F$_1$ je Seed und Validierungs-F$_1$ der Checkpoints]{Mikro-F$_1$ (strikt, in Punkten) der einzelnen Seeds auf dem
  MultiNERD-Testsplit (Hauptläufe mit 256 Ausgabetokens) sowie
  Validierungs-F$_1$ des gewählten Checkpoints mit Epoche in Klammern
  (Encoder: vollständiger Validierungssplit; QLoRA: generative
  Validierung auf 200 Sätzen).}
  \label{tab:app-seeds}
  \begin{tabular}{llccc}
    \toprule
    \textbf{Modell} & \textbf{Regime} & \textbf{Seed 42} & \textbf{Seed 123} & \textbf{Seed 456} \\
    \midrule
    DeBERTa-v3-base & Fine-Tuning & 95,49 & 95,55 & 95,47 \\
    \quad Validierung (Epoche) & & 92,65 (4) & 92,57 (5) & 92,66 (5) \\
    DeBERTa-v3-large & Fine-Tuning & 95,83 & 95,78 & 95,84 \\
    \quad Validierung (Epoche) & & 93,08 (3) & 92,89 (4) & 93,05 (4) \\
    Qwen3.5-0.8B & QLoRA & 93,24 & 93,14 & 93,04 \\
    \quad Validierung (Epoche) & & 92,25 (3) & 93,15 (3) & 93,64 (3) \\
    Qwen3.5-4B & QLoRA & 93,93 & 94,02 & 93,96 \\
    \quad Validierung (Epoche) & & 95,45 (2) & 94,18 (3) & 95,05 (3) \\
    Qwen3.5-27B & QLoRA & 94,53 & 94,39 & 94,52 \\
    \quad Validierung (Epoche) & & 96,13 (3) & 94,68 (3) & 95,58 (3) \\
    \bottomrule
\end{tabular}

\end{table}

\begin{table}[ht]
  \centering
  \footnotesize
  \caption[Nachinferenz mit 1\,024 Tokens je Seed]{Nachinferenz der
  abgeschnittenen Antworten mit 1\,024 Ausgabetokens je Seed: Mikro-F$_1$
  bei 256 und 1\,024 Tokens, Differenz, Zahl der neu erzeugten und der
  weiterhin abgeschnittenen Antworten sowie Zahl der Antworten, deren
  erste 256 Tokens die gespeicherte Antwort reproduzieren.}
  \label{tab:app-limit-seeds}
  \setlength{\tabcolsep}{4pt}
  \begin{tabular}{llrrrrrr}
    \toprule
    \textbf{Modell} & \textbf{Seed} & \textbf{F$_1$ (256)} & \textbf{F$_1$ (1\,024)} & $\Delta$\textbf{F$_1$} &
      \textbf{\shortstack[c]{neu\\erzeugt}} & \textbf{\shortstack[c]{weiterhin\\abgeschnitten}} &
      \textbf{\shortstack[c]{Pr\"afix\\gleich}} \\
    \midrule
    Qwen3.5-0.8B & 42 & 93,24 & 94,35 & $+$1,11 & 68 & 0 & 67/68 \\
     & 123 & 93,14 & 94,29 & $+$1,14 & 68 & 0 & 67/68 \\
     & 456 & 93,04 & 94,16 & $+$1,12 & 67 & 0 & 67/67 \\
    \midrule
    Qwen3.5-4B & 42 & 93,93 & 95,06 & $+$1,14 & 72 & 0 & 70/72 \\
     & 123 & 94,02 & 95,15 & $+$1,14 & 67 & 0 & 66/67 \\
     & 456 & 93,96 & 95,11 & $+$1,15 & 67 & 0 & 66/67 \\
    \midrule
    Qwen3.5-27B & 42 & 94,53 & 95,68 & $+$1,15 & 65 & 0 & 63/65 \\
     & 123 & 94,39 & 95,57 & $+$1,18 & 68 & 0 & 67/68 \\
     & 456 & 94,52 & 95,67 & $+$1,15 & 62 & 0 & 61/62 \\
    \bottomrule
\end{tabular}

\end{table}

\begin{table}[ht]
  \centering
  \scriptsize
  \caption[Precision, Recall und F$_1$ je Entitätstyp]{Precision, Recall und F$_1$ je Entitätstyp (strikt, in Punkten,
  256 Ausgabetokens) für DeBERTa-v3-large und Qwen3.5-27B (QLoRA), Mittel
  über drei Seeds, sowie für Qwen3.5-27B im Zero-Shot-Regime (Einzellauf).
  Support nach Tabelle~\ref{tab:datensplits}.}
  \label{tab:app-typen-prf}
  \setlength{\tabcolsep}{3pt}
  \begin{tabular}{lrrrrrrrrrr}
    \toprule
    & & \multicolumn{3}{c}{\textbf{DeBERTa-v3-large}} &
      \multicolumn{3}{c}{\textbf{Qwen3.5-27B (QLoRA)}} &
      \multicolumn{3}{c}{\textbf{Qwen3.5-27B (Zero-Shot)}} \\
    \cmidrule(lr){3-5}\cmidrule(lr){6-8}\cmidrule(lr){9-11}
    \textbf{Typ} & \textbf{Support} & P & R & F$_1$ & P & R & F$_1$ & P & R & F$_1$ \\
    \midrule
    PER   & 5\,265 & 99,7 & 99,7 & 99,7 & 99,7 & 97,6 & 98,7 & 90,5 & 85,7 & 88,0 \\
    ORG   & 3\,309 & 98,6 & 99,0 & 98,8 & 99,0 & 98,8 & 98,9 & 87,7 & 76,8 & 81,9 \\
    LOC   & 12\,024 & 99,7 & 99,6 & 99,7 & 99,6 & 96,9 & 98,2 & 92,0 & 82,1 & 86,8 \\
    CEL   & 41 & 78,4 & 82,1 & 80,2 & 88,5 & 81,3 & 84,7 & 62,5 & 48,8 & 54,8 \\
    DIS   & 759 & 80,1 & 83,8 & 81,9 & 78,9 & 81,2 & 80,0 & 57,4 & 49,4 & 53,1 \\
    EVE   & 352 & 96,1 & 98,3 & 97,2 & 97,0 & 96,3 & 96,6 & 37,5 & 68,5 & 48,4 \\
    MEDIA & 458 & 97,4 & 98,3 & 97,8 & 97,4 & 98,3 & 97,8 & 62,9 & 41,5 & 50,0 \\
    ANIM  & 1\,604 & 76,7 & 82,3 & 79,4 & 74,4 & 73,6 & 74,0 & 54,0 & 38,7 & 45,0 \\
    BIO   & 8 & 72,6 & 87,5 & 79,0 & 59,2 & 79,2 & 67,5 & 1,6 & 100,0 & 3,2 \\
    FOOD  & 566 & 72,5 & 73,4 & 72,9 & 70,9 & 69,3 & 70,1 & 27,6 & 29,9 & 28,7 \\
    PLANT & 894 & 68,6 & 77,1 & 72,6 & 66,9 & 73,6 & 70,1 & 44,4 & 28,6 & 34,8 \\
    INST  & 12 & 85,5 & 80,6 & 82,7 & 97,2 & 97,2 & 97,2 & 4,9 & 33,3 & 8,6 \\
    MYTH  & 32 & 78,6 & 86,5 & 82,2 & 84,9 & 92,7 & 88,6 & 51,3 & 62,5 & 56,3 \\
    TIME  & 289 & 83,2 & 86,6 & 84,9 & 88,5 & 86,4 & 87,4 & 1,0 & 18,7 & 1,9 \\
    VEHI  & 32 & 90,0 & 92,7 & 91,3 & 92,0 & 93,8 & 92,8 & 10,0 & 65,6 & 17,4 \\
    \bottomrule
\end{tabular}

\end{table}

\begin{table}[ht]
  \centering
  \footnotesize
  \caption[CrossNER-Transfertest: F$_1$ je Domäne und Seed]{CrossNER-Transfertest
  mit 256 Ausgabetokens: F$_1$ (strikt, nur PER, ORG und LOC, in Punkten)
  je Domäne und Seed für DeBERTa-v3-large und Qwen3.5-27B (QLoRA).
  Gepoolt: Mikrowert über alle 2\,506 Sätze; Domänenmittel: ungewichtetes
  Mittel der fünf Domänen.}
  \label{tab:app-crossner-seeds}
  \setlength{\tabcolsep}{4pt}
  \begin{tabular}{lrrrrrrr}
    \toprule
    & & \multicolumn{3}{c}{\textbf{DeBERTa-v3-large}} &
      \multicolumn{3}{c}{\textbf{Qwen3.5-27B (QLoRA)}} \\
    \cmidrule(lr){3-5}\cmidrule(lr){6-8}
    \textbf{Dom\"ane} & \textbf{S\"atze} & \textbf{42} & \textbf{123} & \textbf{456} & \textbf{42} & \textbf{123} & \textbf{456} \\
    \midrule
    K\"unstliche Intelligenz & 431 & 82,14 & 81,48 & 82,32 & 75,16 & 74,40 & 74,87 \\
    Literatur & 416 & 83,28 & 84,07 & 84,65 & 77,21 & 78,67 & 76,96 \\
    Musik & 465 & 85,72 & 84,02 & 85,42 & 62,81 & 63,55 & 62,56 \\
    Politik & 651 & 86,79 & 86,18 & 87,81 & 71,33 & 72,44 & 70,99 \\
    Naturwissenschaften & 543 & 81,89 & 82,59 & 81,59 & 75,93 & 76,02 & 76,48 \\
    \midrule
    Gepoolt & 2\,506 & 84,98 & 84,52 & 85,46 & 71,49 & 72,27 & 71,31 \\
    Dom\"anenmittel & -- & 83,97 & 83,67 & 84,36 & 72,49 & 73,01 & 72,37 \\
    \bottomrule
\end{tabular}

\end{table}

\begin{table}[ht]
  \centering
  \footnotesize
  \caption[CrossNER-Transfertest: F$_1$ je Grobtyp und Domäne]{CrossNER-Transfertest
  mit 256 Ausgabetokens: F$_1$ je Grobtyp und Domäne (strikt, in Punkten,
  Mittel über drei Seeds) mit Support je Typ.}
  \label{tab:app-crossner-typen}
  \setlength{\tabcolsep}{3pt}
  \begin{tabular}{lrrrrrrrrr}
    \toprule
    & \multicolumn{3}{c}{\textbf{Support}} & \multicolumn{3}{c}{\textbf{DeBERTa-v3-large}} &
      \multicolumn{3}{c}{\textbf{Qwen3.5-27B (QLoRA)}} \\
    \cmidrule(lr){2-4}\cmidrule(lr){5-7}\cmidrule(lr){8-10}
    \textbf{Dom\"ane} & PER & ORG & LOC & PER & ORG & LOC & PER & ORG & LOC \\
    \midrule
    K\"unstliche Intelligenz & 227 & 173 & 83 & 86,6 & 72,5 & 92,8 & 76,5 & 68,0 & 85,8 \\
    Literatur & 742 & 110 & 200 & 87,1 & 69,9 & 82,1 & 80,6 & 67,3 & 73,0 \\
    Musik & 595 & 670 & 503 & 86,8 & 82,2 & 86,9 & 62,9 & 62,4 & 63,8 \\
    Politik & 839 & 1\,466 & 1\,017 & 91,4 & 86,7 & 84,0 & 72,4 & 75,7 & 64,5 \\
    Naturwissenschaften & 701 & 403 & 197 & 89,7 & 71,1 & 82,5 & 77,9 & 75,4 & 72,2 \\
    \bottomrule
\end{tabular}

\end{table}

\begin{table}[ht]
  \centering
  \footnotesize
  \caption[Beispiele typischer Fehler je Fehlerklasse]{Beispiele typischer Fehler je Fehlerklasse
  (Abschnitt~\ref{sec:erg-fehlertypen}) sowie der Formtreue-Kategorie
  \enquote{Unbekannter Typ} (Abschnitt~\ref{sec:erg-formtreue}) aus den
  Testvorhersagen mit 256 Ausgabetokens. Illustrative, nicht zufällig
  gezogene Auswahl; Nr.\ ist der nullbasierte Listenindex des Satzes in
  \texttt{test\_predictions.json}; Sätze aus dem MultiNERD-Testsplit
  \cite{tedeschi2022multinerd}. Referenz und Vorhersage als Span mit
  Typ; ein Strich bedeutet, dass die Referenz bzw. die Vorhersage keinen
  Span an dieser Stelle enthält. Seed~42, sofern nicht anders angegeben.}
  \label{tab:app-beispiele}
  \setlength{\tabcolsep}{3pt}
  \begin{tabularx}{\linewidth}{@{}>{\raggedleft\arraybackslash}p{0.85cm}
      >{\raggedright\arraybackslash}X
      >{\raggedright\arraybackslash}p{1.95cm}
      >{\raggedright\arraybackslash}p{1.95cm}
      >{\raggedright\arraybackslash}p{2.15cm}
      >{\raggedright\arraybackslash}p{2.55cm}@{}}
    \toprule
    \textbf{Nr.} & \textbf{Satz (gekürzt)} & \textbf{Referenz} & \textbf{Vorhersage} &
    \textbf{Fehlerklasse} & \textbf{Konfiguration} \\
    \midrule
    327 & He is also in the band Filthy Friends, playing bass and accordion. &
      Filthy Friends (PER) & Filthy Friends (ORG) & Typfehler &
      DeBERTa-v3-large und Qwen3.5-27B (QLoRA) \\
    396 & The lake is home to perch, pike-perch, bream, roaches, whitefishes,
      smelt and other species of fish. &
      perch (ANIM) & perch (FOOD) & Typfehler & DeBERTa-v3-large \\
    59 & In 2000, the entire campus was designated as an arboretum, the
      University of Georgia Campus Arboretum. &
      University of Georgia Campus Arboretum (ORG) &
      University of Georgia Campus (ORG); Arboretum (LOC) & Grenzfehler
      (mit Zusatz) & DeBERTa-v3-large \\
    59 & dito &
      University of Georgia Campus Arboretum (ORG) &
      University of Georgia (ORG); Campus Arboretum (LOC) &
      Grenz- und Typfehler & DeBERTa-v3-large (Seed~123) \\
    383 & The wilderness provides habitat for wild burro, desert tortoise,
      spotted bat and desert bighorn sheep. &
      burro (ANIM) & wild burro (ANIM) & Grenzfehler & Qwen3.5-27B (QLoRA) \\
    35 & It takes around 6,000 full-grown coconuts to produce one tonne of
      copra. &
      copra (FOOD) & -- & Auslassung & DeBERTa-v3-large \\
    4 & Within city limits, the southwestern fence lizard and whiptails are
      common. &
      whiptails (ANIM) & -- & Auslassung & Qwen3.5-27B (QLoRA) \\
    5 & Commonly seen mammals include the coyote, rock squirrel, Gunnison's
      prairie dog, desert cottontail, and black-tailed jackrabbit. &
      -- & desert cottontail (ANIM) & Zusatz &
      DeBERTa-v3-large und Qwen3.5-27B (QLoRA) \\
    10805 & Propranolol is a beta blocker that is mainly used to treat
      hypertension, but also prevents trembling. &
      -- & Propranolol (MED) & Unbekannter Typ & Qwen3.5-27B (Zero-Shot) \\
    2405 & The municipal holiday is Easter Monday. &
      Easter Monday (TIME) & Easter Monday (TME) & Unbekannter Typ &
      Qwen3.5-27B (Zero-Shot) \\
    \bottomrule
  \end{tabularx}
\end{table}


\end{document}